\documentclass[11pt]{article}

\usepackage[margin=1in]{geometry}
\usepackage{amsmath,amssymb,mathtools}
\usepackage{braket}
\usepackage{hyperref}
\usepackage{listings}
\usepackage{xcolor}

\usepackage{amsthm}
\newtheorem{lemma}{Lemma}
\newtheorem{proposition}{Proposition}
\newtheorem{theorem}{Theorem}
\newtheorem{remark}{Remark}
\newtheorem{corollary}{Corollary}

\newcommand{\wt}{\mathrm{wt}}
\newcommand{\supp}{\mathrm{supp}}
\newcommand{\Bit}{\mathrm{Bit}}

\hypersetup{
  colorlinks=true,
  linkcolor=blue,
  urlcolor=blue,
  citecolor=blue
}

\lstset{
  basicstyle=\ttfamily\small,
  columns=fullflexible,
  breaklines=true,
  frame=single,
  keywordstyle=\color{blue},
  commentstyle=\color{gray},
  stringstyle=\color{teal},
  showstringspaces=false
}

\title{Notes on the Distance-3 Exact BD16 Examples}
\author{(working notes)}
\date{}

\begin{document}
\maketitle

\section{Scope and goal}

These notes summarize the \emph{subset-sum--linear-programming} (SS--LP / SSLP) pipeline for discovering \(((n,2,3))\) nonadditive quantum codes with prescribed \emph{transversal diagonal} gates, and then work through one \emph{exact} BD16 (transversal \(T\)) example in detail.

We follow the setup in:
\begin{itemize}
\item Chao Zhang, Zipeng Wu, Shilin Huang, Bei Zeng, \emph{Transversal Gates in Nonadditive Quantum Codes}, April 30, 2025.
\end{itemize}

Throughout, \(n=7\), \(K=2\), and we write Pauli operators as tensor products from \(\{I,X,Y,Z\}\).

\section{The SSLP pipeline for distance-3 codes (Lean-friendly)}

This section rewrites the SSLP logic in a way that is meant to be directly formalizable:
all objects are finite, all constraints are explicit equalities over finite sums, and we
make systematic use of the \emph{complementary} convention
\[
\ket{1_L}=X^{\otimes n}\ket{0_L}.
\]

\subsection{Data, notation, and the complementary convention}

\paragraph{Bitstrings.}
Fix an integer \(n\ge 1\). We write computational basis labels as bitstrings
\[
s=(s_1,\dots,s_n)\in\{0,1\}^n,
\qquad
\ket{s}:=\ket{s_1}\otimes\cdots\otimes\ket{s_n}.
\]
For bitstrings \(s,t\in\{0,1\}^n\),
\[
(s\oplus t)_j := s_j+t_j\pmod 2,
\qquad
z\cdot s := \sum_{j=1}^n z_js_j\pmod 2,
\qquad
(-1)^{z\cdot s}\in\{\pm 1\}.
\]
Let \(\mathbf{1}:=(1,\dots,1)\in\{0,1\}^n\) and define the \emph{bitwise complement}
\[
\overline{s}:=s\oplus\mathbf{1}.
\]
Let \(e_i\in\{0,1\}^n\) denote the standard basis bitstring with \((e_i)_i=1\) and all other entries \(0\).

\paragraph{Logical basis with complementary symmetry.}
We focus on \(((n,2,3))\) codes with a fixed orthonormal basis \(\{\ket{0_L},\ket{1_L}\}\)
and impose the \emph{complementary} condition
\[
\boxed{\;\ket{1_L}=C\ket{0_L},\quad C:=X^{\otimes n}.\;}
\]
If
\[
\ket{0_L}=\sum_{s\in\{0,1\}^n}c_s\ket{s},
\]
then
\[
\ket{1_L}=X^{\otimes n}\ket{0_L}=\sum_{s\in\{0,1\}^n}c_s\ket{\overline{s}}
=\sum_{t\in\{0,1\}^n}c_{\overline{t}}\ket{t}.
\]
So the amplitude functions satisfy
\[
\boxed{\;c^{(1)}_t = c^{(0)}_{\overline{t}}.\;}
\]
In particular, all unknowns can be taken to be the amplitudes \(c_s\) of \(\ket{0_L}\); \(\ket{1_L}\) is then determined.

\subsection{Stage 1: Subset-sum support (SS)}

Fix a modulus \(m\in\mathbb{N}\) and an \emph{angle vector}
\(\vec a=(a_1,\dots,a_n)\in\mathbb{Z}_m^n\).
Let \(\omega:=e^{2\pi i/m}\) and define the transversal diagonal
\[
D(\vec a):=\bigotimes_{j=1}^n\left(\ket{0}\!\bra{0}+\omega^{a_j}\ket{1}\!\bra{1}\right).
\]
On a computational basis state,
\[
D(\vec a)\ket{s}=\omega^{\langle \vec a,s\rangle}\ket{s},
\qquad
\langle \vec a,s\rangle := \sum_{j=1}^n a_js_j\ \ (\mathrm{mod}\ m).
\]

\paragraph{Residue classes.}
If \(D(\vec a)\) acts \emph{logically diagonally} with residues \(b_0,b_1\in\mathbb{Z}_m\),
i.e.
\[
D(\vec a)\ket{k_L}=\omega^{b_k}\ket{k_L},\qquad k\in\{0,1\},
\]
then every basis component of \(\ket{k_L}\) must lie in a single residue class:
\[
c^{(k)}_s\neq 0\implies \langle \vec a,s\rangle\equiv b_k\pmod m.
\]
Define the SS sets
\[
S_k(\vec a):=\left\{s\in\{0,1\}^n:\ \langle \vec a,s\rangle\equiv b_k\pmod m\right\}.
\]
The SS constraint is
\[
\mathrm{supp}(\ket{k_L})\subseteq S_k(\vec a),\qquad k\in\{0,1\}.
\]

\paragraph{Complementary compatibility.}
Under \(\ket{1_L}=X^{\otimes n}\ket{0_L}\), supports satisfy
\[
\mathrm{supp}(\ket{1_L})=\overline{\mathrm{supp}(\ket{0_L})}.
\]
Thus, if \(\mathrm{supp}(\ket{0_L})\subseteq S_0(\vec a)\), then necessarily
\[
\mathrm{supp}(\ket{1_L})\subseteq \overline{S_0(\vec a)}:=\{\overline{s}:s\in S_0(\vec a)\}.
\]
Moreover,
\[
\langle \vec a,\overline{s}\rangle
=\langle \vec a,s\oplus\mathbf{1}\rangle
\equiv \langle \vec a,s\rangle + \langle \vec a,\mathbf{1}\rangle\pmod m,
\quad\text{where}\quad
\langle \vec a,\mathbf{1}\rangle:=\sum_{j=1}^n a_j\ (\mathrm{mod}\ m).
\]
So the residues must satisfy
\[
\boxed{\;b_1\equiv b_0+\sum_{j=1}^n a_j\pmod m,\qquad S_1(\vec a)=\overline{S_0(\vec a)}.\;}
\]
In particular, under the complementary convention one may choose \(b_0\) and then \emph{define} \(b_1\) by the above congruence.

\subsection{Stage 2: Linear-program filter (LP) = \(Z\)-type KL constraints}

\paragraph{Magnitude--phase parametrization.}
For each \(s\in S_0(\vec a)\), write
\[
c_s=\sqrt{p_s}\,e^{i\phi_s},
\qquad p_s\ge 0,\quad \phi_s\in\mathbb{R},
\]
and set \(c_s:=0\) for \(s\notin S_0(\vec a)\).
Normalization is
\[
\boxed{\;\sum_{s\in S_0(\vec a)} p_s = 1.\;}
\]
Under the complementary convention,
\[
c^{(1)}_t=c_{\overline{t}},
\qquad
p^{(1)}_t=p_{\overline{t}},
\qquad
\phi^{(1)}_t=\phi_{\overline{t}}.
\]

\paragraph{\(Z\)-type matrix elements are linear in \(p\).}
For any \(z\in\{0,1\}^n\), define the diagonal Pauli
\[
Z^z := \bigotimes_{j=1}^n Z^{z_j}.
\]
It acts as \(Z^z\ket{s}=(-1)^{z\cdot s}\ket{s}\), hence
\[
\bra{0_L}Z^z\ket{0_L}=\sum_{s\in S_0(\vec a)} (-1)^{z\cdot s}\,p_s.
\]
These expressions are linear in the LP variables \(\{p_s\}\).

\paragraph{Distance-3 \(Z\)-type KL constraints (weight \(\le 2\)).}
For distance \(3\), it is enough to enforce KL for all Paulis of weight \(\le 2\).
In the \(Z\)-type sector this means \(z\) with \(\mathrm{wt}(z)\in\{1,2\}\), i.e.
\(z=e_i\) and \(z=e_i\oplus e_j\) for \(i<j\).

Under the complementary convention we can simplify diagonal equality:
\[
\bra{1_L}Z^z\ket{1_L}
=\bra{0_L}C Z^z C\ket{0_L}
=\bra{0_L}(-1)^{\mathrm{wt}(z)}Z^z\ket{0_L}
=(-1)^{\mathrm{wt}(z)}\bra{0_L}Z^z\ket{0_L}.
\]
Therefore the KL diagonal condition
\(\bra{0_L}Z^z\ket{0_L}=\bra{1_L}Z^z\ket{1_L}\) is equivalent to
\[
\boxed{\;\mathrm{wt}(z)\ \text{odd}\ \Longrightarrow\ \bra{0_L}Z^z\ket{0_L}=0.\;}
\]
In particular, for all \(i\) we must have
\[
\boxed{\;\bra{0_L}Z_i\ket{0_L}=0
\iff
\sum_{s\in S_0(\vec a)}(-1)^{s_i}p_s=0
\iff
\sum_{s:s_i=0}p_s=\sum_{s:s_i=1}p_s=\frac12.\;}
\]
For \(Z_iZ_j\) (even \(\mathrm{wt}(z)=2\)), the diagonal equality is automatic under complement; the value \(\bra{0_L}Z_iZ_j\ket{0_L}\) is not forced to be \(0\).

\paragraph{LP stage summary.}
The LP stage searches for a feasible probability vector \(p:S_0(\vec a)\to\mathbb{Q}\) (or \(\mathbb{R}\)) satisfying:
\[
p_s\ge 0,\quad \sum_{s\in S_0}p_s=1,\quad \sum_{s\in S_0}(-1)^{s_i}p_s=0\ \ (i=1,\dots,n),
\]
(and optionally any additional linear \(Z\)-type constraints one chooses to include).
The output is a candidate \(p\); the phases \(\phi_s\) are deferred to Full-KL.

\subsection{Stage 3: Full-KL as explicit phase equations (all Paulis of weight \(\le 2\))}

\subsubsection{Pauli strings as \((x,z)\) pairs}

Represent a Pauli string (with \emph{no global phase}) by two bitstrings
\[
x,z\in\{0,1\}^n,
\]
with the per-qubit identification
\[
(x_i,z_i)=(0,0)\leftrightarrow I,\quad
(1,0)\leftrightarrow X,\quad
(0,1)\leftrightarrow Z,\quad
(1,1)\leftrightarrow Y.
\]
Define
\[
N_Y(x,z):=\sum_{i=1}^n x_i z_i\in\mathbb{N},
\qquad
P(x,z):= i^{N_Y(x,z)}X^x Z^z,
\]
where \(X^x:=\bigotimes_i X^{x_i}\) and \(Z^z:=\bigotimes_i Z^{z_i}\).
Then for any basis state \(\ket{s}\),
\[
\boxed{\;P(x,z)\ket{s}=i^{N_Y(x,z)}(-1)^{z\cdot s}\ket{s\oplus x}.\;}
\]
The (Pauli) weight is
\[
\mathrm{wt}(P(x,z)):=\#\{i:(x_i,z_i)\neq(0,0)\}.
\]

\subsubsection{General full-KL matrix element as a finite sum}

Write logical states as
\[
\ket{k_L}=\sum_{s\in\{0,1\}^n} c^{(k)}_s\ket{s},
\qquad c^{(k)}_s=0\ \text{if }s\notin S_k(\vec a).
\]
Then for any Pauli \(P(x,z)\),
\[
\boxed{
\bra{j_L}P(x,z)\ket{k_L}
=i^{N_Y(x,z)}\sum_{s\in\{0,1\}^n}(-1)^{z\cdot s}\,
\overline{c^{(j)}_{s\oplus x}}\;c^{(k)}_{s}.
}
\]
In magnitude--phase form \(c^{(k)}_s=\sqrt{p^{(k)}_s}e^{i\phi^{(k)}_s}\) this becomes
\[
\boxed{
\bra{j_L}P(x,z)\ket{k_L}
=i^{N_Y(x,z)}\sum_{s\in\{0,1\}^n}(-1)^{z\cdot s}\,
\sqrt{p^{(j)}_{s\oplus x}}\sqrt{p^{(k)}_{s}}\;
e^{i(\phi^{(k)}_{s}-\phi^{(j)}_{s\oplus x})},
}
\]
with the convention \(p^{(k)}_t=0\) (hence the summand vanishes) whenever \(t\notin S_k(\vec a)\).

\subsubsection{Distance-3 KL constraints specialized to \(K=2\)}

Let \(\mathcal{P}_{\le 2}\) be the set of phase-free Pauli strings \(\{I,X,Y,Z\}^{\otimes n}\)
of weight \(\le 2\).
A code corrects all single-qubit errors (distance \(3\)) iff for every \(P\in\mathcal{P}_{\le 2}\),
\[
\boxed{
\bra{0_L}P\ket{1_L}=0,
\qquad
\bra{0_L}P\ket{0_L}=\bra{1_L}P\ket{1_L}.
}
\]

\paragraph{Complementary simplification of diagonal constraints.}
Using \(\ket{1_L}=C\ket{0_L}\) and \(C Z C = -Z\), \(C Y C = -Y\), \(C X C = X\) (per qubit), we have
\[
C\,P(x,z)\,C = (-1)^{\mathrm{wt}(z)}P(x,z),
\qquad \mathrm{wt}(z):=\sum_i z_i.
\]
Hence
\[
\bra{1_L}P(x,z)\ket{1_L}=\bra{0_L}C P(x,z) C\ket{0_L}
=(-1)^{\mathrm{wt}(z)}\bra{0_L}P(x,z)\ket{0_L}.
\]
Therefore the KL diagonal equality is equivalent to
\[
\boxed{\;\mathrm{wt}(z)\ \text{odd}\ \Longrightarrow\ \bra{0_L}P(x,z)\ket{0_L}=0,\;}
\]
while for \(\mathrm{wt}(z)\) even the diagonal equality holds automatically.

\paragraph{Complementary reduction of unknowns in off-diagonals.}
Since \(c^{(1)}_s=c^{(0)}_{\overline{s}}\), we can rewrite all off-diagonal constraints using only
the amplitudes of \(\ket{0_L}\):
\[
\boxed{
\bra{0_L}P(x,z)\ket{1_L}
=i^{N_Y(x,z)}\sum_{s\in\{0,1\}^n}(-1)^{z\cdot s}\,
\overline{c_{s\oplus x}}\;c_{\overline{s}},
}
\]
where now \(c_s:=c^{(0)}_s\) and \(c_s=0\) if \(s\notin S_0(\vec a)\).

\subsubsection{Explicit weight-\(\le 2\) Pauli constraints (ready-to-code forms)}

Below are the concrete finite-sum formulas for all Pauli types needed at distance \(3\).
In every case, the full-KL requirements are:
\[
\bra{0_L}(\cdot)\ket{1_L}=0,\quad
\text{and if }\mathrm{wt}(z)\text{ is odd then }\bra{0_L}(\cdot)\ket{0_L}=0,
\]
with \(\mathrm{wt}(z)\) computed from the \(Z\)-components (i.e. number of \(Z\) or \(Y\) factors).

Throughout, sums may be taken over all \(s\in\{0,1\}^n\) with the convention \(c_t=0\) for \(t\notin S_0(\vec a)\).

\paragraph{Weight 0.}
\[
\bra{0_L}I\ket{0_L}=\sum_s |c_s|^2=1,
\qquad
\bra{0_L}I\ket{1_L}=\sum_s \overline{c_s}\,c_{\overline{s}}=0.
\]

\paragraph{Weight 1: single-qubit Paulis (for each \(i\)).}
\[
\bra{0_L}X_i\ket{0_L}=\sum_s \overline{c_{s\oplus e_i}}\,c_s,
\qquad
\bra{0_L}X_i\ket{1_L}=\sum_s \overline{c_{s\oplus e_i}}\,c_{\overline{s}}.
\]
\[
\bra{0_L}Z_i\ket{0_L}=\sum_s (-1)^{s_i}|c_s|^2,
\qquad
\bra{0_L}Z_i\ket{1_L}=\sum_s (-1)^{s_i}\overline{c_s}\,c_{\overline{s}}.
\]
\[
\bra{0_L}Y_i\ket{0_L}= i\sum_s (-1)^{s_i}\overline{c_{s\oplus e_i}}\,c_s,
\qquad
\bra{0_L}Y_i\ket{1_L}= i\sum_s (-1)^{s_i}\overline{c_{s\oplus e_i}}\,c_{\overline{s}}.
\]
Diagonal simplification: since \(Z_i\) and \(Y_i\) have \(\mathrm{wt}(z)=1\), complement forces
\(\bra{0_L}Z_i\ket{0_L}=0\) and \(\bra{0_L}Y_i\ket{0_L}=0\).

\paragraph{Weight 2: two-qubit Paulis (for each \(i<j\)).}
\[
\bra{0_L}X_iX_j\ket{0_L}=\sum_s \overline{c_{s\oplus e_i\oplus e_j}}\,c_s,
\qquad
\bra{0_L}X_iX_j\ket{1_L}=\sum_s \overline{c_{s\oplus e_i\oplus e_j}}\,c_{\overline{s}}.
\]
\[
\bra{0_L}Z_iZ_j\ket{0_L}=\sum_s (-1)^{s_i+s_j}|c_s|^2,
\qquad
\bra{0_L}Z_iZ_j\ket{1_L}=\sum_s (-1)^{s_i+s_j}\overline{c_s}\,c_{\overline{s}}.
\]
\[
\bra{0_L}X_iZ_j\ket{0_L}=\sum_s (-1)^{s_j}\overline{c_{s\oplus e_i}}\,c_s,
\qquad
\bra{0_L}X_iZ_j\ket{1_L}=\sum_s (-1)^{s_j}\overline{c_{s\oplus e_i}}\,c_{\overline{s}}.
\]
\[
\bra{0_L}Z_iX_j\ket{0_L}=\sum_s (-1)^{s_i}\overline{c_{s\oplus e_j}}\,c_s,
\qquad
\bra{0_L}Z_iX_j\ket{1_L}=\sum_s (-1)^{s_i}\overline{c_{s\oplus e_j}}\,c_{\overline{s}}.
\]
\[
\bra{0_L}X_iY_j\ket{0_L}= i\sum_s (-1)^{s_j}\overline{c_{s\oplus e_i\oplus e_j}}\,c_s,
\qquad
\bra{0_L}X_iY_j\ket{1_L}= i\sum_s (-1)^{s_j}\overline{c_{s\oplus e_i\oplus e_j}}\,c_{\overline{s}}.
\]
\[
\bra{0_L}Y_iX_j\ket{0_L}= i\sum_s (-1)^{s_i}\overline{c_{s\oplus e_i\oplus e_j}}\,c_s,
\qquad
\bra{0_L}Y_iX_j\ket{1_L}= i\sum_s (-1)^{s_i}\overline{c_{s\oplus e_i\oplus e_j}}\,c_{\overline{s}}.
\]
\[
\bra{0_L}Y_iY_j\ket{0_L}= -\sum_s (-1)^{s_i+s_j}\overline{c_{s\oplus e_i\oplus e_j}}\,c_s,
\qquad
\bra{0_L}Y_iY_j\ket{1_L}= -\sum_s (-1)^{s_i+s_j}\overline{c_{s\oplus e_i\oplus e_j}}\,c_{\overline{s}}.
\]
\[
\bra{0_L}Y_iZ_j\ket{0_L}= i\sum_s (-1)^{s_i+s_j}\overline{c_{s\oplus e_i}}\,c_s,
\qquad
\bra{0_L}Y_iZ_j\ket{1_L}= i\sum_s (-1)^{s_i+s_j}\overline{c_{s\oplus e_i}}\,c_{\overline{s}}.
\]
\[
\bra{0_L}Z_iY_j\ket{0_L}= i\sum_s (-1)^{s_i+s_j}\overline{c_{s\oplus e_j}}\,c_s,
\qquad
\bra{0_L}Z_iY_j\ket{1_L}= i\sum_s (-1)^{s_i+s_j}\overline{c_{s\oplus e_j}}\,c_{\overline{s}}.
\]

Diagonal simplification at weight \(2\):
\begin{itemize}
\item Operators with exactly one \(Z\)-component (i.e. \(\mathrm{wt}(z)=1\)) must have
\(\bra{0_L}(\cdot)\ket{0_L}=0\).
These are \(X_iZ_j,\ Z_iX_j,\ X_iY_j,\ Y_iX_j\).
\item Operators with \(\mathrm{wt}(z)=0\) or \(2\) have no forced diagonal vanishing from complement.
These are \(X_iX_j\) (\(\mathrm{wt}(z)=0\)) and \(Z_iZ_j,\ Y_iY_j,\ Y_iZ_j,\ Z_iY_j\) (\(\mathrm{wt}(z)=2\)).
\end{itemize}

\paragraph{Full-KL stage summary (what must be solved).}
Given a feasible LP output \(p_s\) on \(S_0(\vec a)\), the Full-KL stage is the problem of finding phases
\(\phi_s\in\mathbb{R}\) (or equivalently unit-modulus \(\omega_s=e^{i\phi_s}\)) such that all the above
\emph{off-diagonal} sums vanish for every Pauli of weight \(\le 2\), and additionally all \emph{odd-\(\mathrm{wt}(z)\)}
\emph{diagonal} sums vanish (which include \(Z_i\), \(Y_i\), and the mixed \(XZ/ZX/XY/YX\) terms).
Because each constraint is a finite sum with explicit coefficients \(\pm 1,\pm i\), this is directly suitable
for exact verification (e.g. in \(\mathbb{Q}(\sqrt{\cdot},i)\)) once a candidate closed form is proposed.

\section{The SSLP pipeline for distance-3 codes (\(n=7\))}

Throughout this section we fix \(n=7\) and \(K=2\). Computational basis strings are bitstrings
\(s=(s_1,\dots,s_7)\in\{0,1\}^7\), and we write \(\ket{s}\) for the corresponding basis ket.
Bitwise XOR is denoted by \(\oplus\), and the all-ones string is \(1^7=(1,1,1,1,1,1,1)\).
We also use the \(\mathbb{Z}_2\) dot product
\[
u\cdot v \ :=\ \sum_{j=1}^7 u_j v_j \pmod 2,
\qquad
(-1)^{u\cdot v}\in\{\pm 1\}.
\]
For each position \(i\in\{1,\dots,7\}\), let \(e_i\in\{0,1\}^7\) be the unit vector with a single \(1\) at \(i\).

\subsection{Distance-3 Knill--Laflamme constraints in a finite Pauli basis}

Let \(\mathcal{P}_7:=\{I,X,Y,Z\}^{\otimes 7}\) denote the set of \emph{Hermitian} Pauli strings on 7 qubits
(no \(\pm 1,\pm i\) phases).
The \emph{Pauli weight} \(\mathrm{wt}(P)\) is the number of tensor factors of \(P\) that are not \(I\).

A \(((7,2,3))\) code corrects all single-qubit errors. Let
\[
\mathcal{E}_1 := \{I\}\ \cup\ \{X_i,Y_i,Z_i:\ i=1,\dots,7\}
\]
where, e.g., \(X_i\) means \(X\) on qubit \(i\) and \(I\) elsewhere.
The Knill--Laflamme (KL) conditions are:
\[
\bra{j_L}E^\dagger E'\ket{k_L} = \alpha_{E,E'}\,\delta_{jk}
\qquad
\forall\,E,E'\in\mathcal{E}_1,\quad \forall\, j,k\in\{0,1\},
\]
for some scalars \(\alpha_{E,E'}\in\mathbb{C}\) depending only on \((E,E')\).

Since \(E^\dagger E'\) has Pauli weight \(\le 2\), it is equivalent (and Lean-friendly) to verify KL
on the finite set
\[
\mathcal{P}_{\le 2}
\ :=\ \{P\in\mathcal{P}_7:\ \mathrm{wt}(P)\le 2\}.
\]
Concretely, \(\mathcal{P}_{\le 2}\) consists of:
\begin{itemize}
\item weight \(0\): \(I^{\otimes 7}\);
\item weight \(1\): \(X_i,Y_i,Z_i\) for each \(i\);
\item weight \(2\): for each pair \(i<j\), the \(9\) strings \(A_iB_j\) with \(A,B\in\{X,Y,Z\}\).
\end{itemize}

\paragraph{KL constraints in a 2D logical basis.}
Fix an orthonormal logical basis \(\{\ket{0_L},\ket{1_L}\}\).
For each \(P\in\mathcal{P}_{\le 2}\), define the \(2\times 2\) \emph{logical matrix}
\[
M(P)_{jk} := \bra{j_L}P\ket{k_L}.
\]
Then KL for distance \(3\) is exactly the statement:
\[
\boxed{
\forall\,P\in\mathcal{P}_{\le 2},\quad
M(P)=\lambda_P I_2\ \text{for some scalar}\ \lambda_P\in\mathbb{C}.
}
\]
Equivalently, for every \(P\in\mathcal{P}_{\le 2}\),
\[
\boxed{
\bra{0_L}P\ket{1_L}=0,
\qquad
\bra{0_L}P\ket{0_L}=\bra{1_L}P\ket{1_L}.
}
\]
(For \(P=I^{\otimes 7}\) these reduce to normalization and orthogonality:
\(\braket{0_L|0_L}=\braket{1_L|1_L}=1\) and \(\braket{0_L|1_L}=0\).)

\subsection{Stage 1: Subset-sum support (SS)}

Fix a modulus \(m\) and an \emph{angle vector} \(\vec a=(a_1,\dots,a_7)\in \mathbb{Z}_m^7\).
Define the residue map
\[
\mathrm{res}_{\vec a}(s)\ :=\ \sum_{j=1}^7 a_j s_j\ \in\ \mathbb{Z}_m.
\]
Given target residues \((b_0,b_1)\in\mathbb{Z}_m^2\), define the SS blocks
\[
S_k(\vec a)\ :=\ \{\,s\in\{0,1\}^7:\ \mathrm{res}_{\vec a}(s)\equiv b_k \ (\mathrm{mod}\ m)\,\},
\qquad k\in\{0,1\}.
\]
The SS constraint is the \emph{support restriction}
\[
\ket{k_L}=\sum_{s\in S_k(\vec a)} c^{(k)}_s \ket{s},
\qquad k\in\{0,1\}.
\]

\paragraph{Why SS is natural for transversal diagonal gates.}
For
\[
D(\vec a)=\bigotimes_{j=1}^7\left(\ket{0}\!\bra{0}+e^{2\pi i a_j/m}\ket{1}\!\bra{1}\right),
\]
we have \(D(\vec a)\ket{s}=e^{2\pi i\,\mathrm{res}_{\vec a}(s)/m}\ket{s}\).
Thus any eigenstate of \(D(\vec a)\) must have support on residue classes of \(\mathrm{res}_{\vec a}\),
which is exactly what the SS step encodes.

\subsection{Stage 2: Linear-program filter (LP) = Z-type KL constraints as linear equalities}

Define probabilities \(p^{(k)}_s:=|c^{(k)}_s|^2\ge 0\).
The Z-type Paulis are those with no bit flips, i.e. strings consisting only of \(I\) and \(Z\).
For distance \(3\), the Z-type Paulis of weight \(\le 2\) are:
\[
Z_i,\qquad Z_iZ_j\ (i<j),\qquad \text{and }I^{\otimes 7}.
\]

Because these are diagonal in the computational basis, their logical matrix elements are linear in \(p^{(k)}\).
Explicitly, for \(k\in\{0,1\}\),
\[
\bra{k_L}Z_i\ket{k_L}=\sum_{s\in S_k(\vec a)} (-1)^{s_i}\,p^{(k)}_s,
\]
\[
\bra{k_L}Z_iZ_j\ket{k_L}=\sum_{s\in S_k(\vec a)} (-1)^{s_i+s_j}\,p^{(k)}_s,
\]
and normalization is
\[
\sum_{s\in S_k(\vec a)} p^{(k)}_s = 1.
\]

\paragraph{LP feasibility constraints (Z-type KL).}
The Z-type KL requirements
\[
\bra{0_L}Z_i\ket{0_L}=\bra{1_L}Z_i\ket{1_L},\qquad
\bra{0_L}Z_iZ_j\ket{0_L}=\bra{1_L}Z_iZ_j\ket{1_L}
\]
become a system of \emph{linear equalities} in the variables \(\{p^{(0)}_s,p^{(1)}_s\}\),
together with nonnegativity and normalization.
This is the LP stage: it filters SS supports that cannot satisfy Z-type KL.

\paragraph{Remark (off-diagonal Z-type constraints).}
For a diagonal \(P\) (only \(Z\)'s), \(\bra{0_L}P\ket{1_L}\) contains only terms
\(\overline{c_s^{(0)}}c_s^{(1)}\), so if \(S_0(\vec a)\cap S_1(\vec a)=\varnothing\)
(which is typical when \(b_0\neq b_1\)), then \(\bra{0_L}P\ket{1_L}=0\) holds automatically
for all diagonal \(P\). This is one major reason the LP stage is often ``just'' about diagonal constraints.

\subsection{Stage 3: Full-KL = all Pauli constraints of weight \(\le 2\)}

\subsubsection{Magnitude--phase parametrization}

For each \(k\in\{0,1\}\) and each \(s\in S_k(\vec a)\), write
\[
c^{(k)}_s = \sqrt{p^{(k)}_s}\,e^{i\phi^{(k)}_s},
\qquad
p^{(k)}_s\ge 0,\ \ \phi^{(k)}_s\in\mathbb{R}.
\]
For \(s\notin S_k(\vec a)\), we adopt the convention \(p^{(k)}_s=0\) (so \(c^{(k)}_s=0\)).

The LP stage produces candidates \(p^{(k)}\) satisfying all Z-type KL constraints.
Full-KL then searches for phases \(\phi^{(k)}\) such that the remaining KL equalities hold.

\subsubsection{A uniform formula for \texorpdfstring{$\bra{j_L}P\ket{k_L}$}{<jL|P|kL>} (symplectic form)}

Every \(P\in\mathcal{P}_7\) can be encoded by binary vectors
\[
x=x(P)\in\{0,1\}^7,\qquad z=z(P)\in\{0,1\}^7,
\]
and an integer \(N_Y=N_Y(P)\), by setting for each qubit \(t\):
\[
\begin{array}{c|cccc}
P_t & I & X & Z & Y\\ \hline
x_t & 0 & 1 & 0 & 1\\
z_t & 0 & 0 & 1 & 1
\end{array}
\qquad\text{and}\qquad
N_Y(P)=\#\{t:\ P_t=Y\}.
\]
Then the action of \(P\) on basis kets is:
\[
P\ket{s} = i^{N_Y(P)}\,(-1)^{z(P)\cdot s}\,\ket{s\oplus x(P)}.
\]
Hence, for any \(j,k\in\{0,1\}\),
\[
\boxed{
\bra{j_L}P\ket{k_L}
=
i^{N_Y(P)}\sum_{s\in\{0,1\}^7}
(-1)^{z(P)\cdot s}\,
\overline{c^{(j)}_{s\oplus x(P)}}\,c^{(k)}_{s}.
}
\]
In magnitude--phase variables this becomes:
\[
\boxed{
\bra{j_L}P\ket{k_L}
=
i^{N_Y(P)}\sum_{s\in\{0,1\}^7}
(-1)^{z(P)\cdot s}\,
\sqrt{p^{(j)}_{s\oplus x(P)}}\,\sqrt{p^{(k)}_{s}}\,
e^{i(\phi^{(k)}_{s}-\phi^{(j)}_{s\oplus x(P)})},
}
\]
with the convention \(p^{(j)}_t=0\) if \(t\notin S_j(\vec a)\) (so terms vanish automatically).

\subsubsection{Explicit specializations for all weight \texorpdfstring{$\le 2$}{<=2} Pauli strings}

Although the symplectic formula above is the cleanest for verification, it is useful to spell out
the common cases that appear when \(\mathrm{wt}(P)\le 2\).

\paragraph{Weight 1 (single-qubit).}
For each \(i\),
\[
\bra{j_L}X_i\ket{k_L}=\sum_{s}\overline{c^{(j)}_{s\oplus e_i}}\,c^{(k)}_{s},
\]
\[
\bra{j_L}Y_i\ket{k_L}=i\sum_{s}(-1)^{s_i}\,\overline{c^{(j)}_{s\oplus e_i}}\,c^{(k)}_{s},
\]
\[
\bra{j_L}Z_i\ket{k_L}=\sum_{s}(-1)^{s_i}\,\overline{c^{(j)}_{s}}\,c^{(k)}_{s}.
\]

\paragraph{Weight 2 (two-qubit, \(i<j\)).}
For each pair \(i<j\),
\[
\bra{j_L}X_iX_j\ket{k_L}=\sum_{s}\overline{c^{(j)}_{s\oplus e_i\oplus e_j}}\,c^{(k)}_{s},
\]
\[
\bra{j_L}X_iY_j\ket{k_L}=i\sum_{s}(-1)^{s_j}\,\overline{c^{(j)}_{s\oplus e_i\oplus e_j}}\,c^{(k)}_{s},
\]
\[
\bra{j_L}Y_iY_j\ket{k_L}=-\sum_{s}(-1)^{s_i+s_j}\,\overline{c^{(j)}_{s\oplus e_i\oplus e_j}}\,c^{(k)}_{s},
\]
\[
\bra{j_L}X_iZ_j\ket{k_L}=\sum_{s}(-1)^{s_j}\,\overline{c^{(j)}_{s\oplus e_i}}\,c^{(k)}_{s},
\]
\[
\bra{j_L}Y_iZ_j\ket{k_L}=i\sum_{s}(-1)^{s_i+s_j}\,\overline{c^{(j)}_{s\oplus e_i}}\,c^{(k)}_{s},
\]
\[
\bra{j_L}Z_iZ_j\ket{k_L}=\sum_{s}(-1)^{s_i+s_j}\,\overline{c^{(j)}_{s}}\,c^{(k)}_{s}.
\]
(The remaining mixed-order strings, e.g. \(Z_iX_j\) or \(Y_iX_j\), are obtained by swapping indices;
since Paulis on distinct qubits commute, e.g. \(X_iZ_j=Z_jX_i\) when \(i\neq j\).)

\paragraph{Full KL constraints (distance 3) in explicit form.}
For every Pauli string \(P\in\mathcal{P}_{\le 2}\), impose:
\[
\boxed{
\bra{0_L}P\ket{1_L}=0,
\qquad
\bra{0_L}P\ket{0_L}=\bra{1_L}P\ket{1_L}.
}
\]
Using the explicit summation formulas above, these become finite systems of complex equalities
in the unknown probabilities \(p^{(k)}\) and phases \(\phi^{(k)}\).

\subsection{How Full-KL simplifies under the complementary ansatz \texorpdfstring{$\ket{1_L}=X^{\otimes 7}\ket{0_L}$}{|1L>=X^7|0L>}}

Many BD16 constructions (including the ones in these notes) impose the \emph{complementary} relation
\[
\boxed{\ket{1_L} = X^{\otimes 7}\ket{0_L}.}
\]
Write
\[
\ket{0_L}=\sum_{s\in\{0,1\}^7} c_s\ket{s}
\qquad\text{and hence}\qquad
\ket{1_L}=\sum_{s} c_s\ket{s\oplus 1^7}.
\]
Equivalently, at the coefficient level,
\[
c^{(1)}_t = c^{(0)}_{t\oplus 1^7},
\qquad
p^{(1)}_t = p^{(0)}_{t\oplus 1^7},
\qquad
\phi^{(1)}_t=\phi^{(0)}_{t\oplus 1^7}.
\]
Thus \emph{only} \(\ket{0_L}\) is independent data; \(\ket{1_L}\) is derived.

\subsubsection{Orthogonality is often automatic under SS}

If SS gives \(\mathrm{supp}(\ket{0_L})\subseteq S_0(\vec a)\) and \(\mathrm{supp}(\ket{1_L})\subseteq S_1(\vec a)\)
with \(S_0(\vec a)\cap S_1(\vec a)=\varnothing\), then
\[
\braket{0_L|1_L}=\sum_{s}\overline{c^{(0)}_s}c^{(1)}_s=0
\]
holds automatically (no shared basis strings).

In BD16 with \((b_0,b_1)=(0,7)\), the compatibility condition
\[
\sum_{j=1}^7 a_j \equiv b_0+b_1 \equiv 7 \pmod 8
\]
ensures that \(\mathrm{res}_{\vec a}(s)=0\Rightarrow \mathrm{res}_{\vec a}(s\oplus 1^7)=7\),
i.e. complement maps \(S_0(\vec a)\) into \(S_1(\vec a)\), so orthogonality becomes SS-level.

\subsubsection{Diagonal KL equalities collapse via conjugation by \texorpdfstring{$X^{\otimes 7}$}{X^7}}

For any Pauli string \(P\in\mathcal{P}_7\),
\[
\bra{1_L}P\ket{1_L}
=
\bra{0_L}X^{\otimes 7} P X^{\otimes 7}\ket{0_L}.
\]
On one qubit, \(XZX=-Z\) and \(XYX=-Y\), while \(XXX=X\) and \(XIX=I\). Therefore,
\[
X^{\otimes 7} P X^{\otimes 7} = (-1)^{\nu(P)} P,
\qquad
\nu(P):=\#\{t:\ P_t\in\{Z,Y\}\}=\|z(P)\|_1.
\]
Hence
\[
\bra{1_L}P\ket{1_L}=(-1)^{\nu(P)}\bra{0_L}P\ket{0_L}.
\]

The diagonal KL equality \(\bra{0_L}P\ket{0_L}=\bra{1_L}P\ket{1_L}\) therefore becomes:
\[
\boxed{
\text{if }\nu(P)\text{ is even, the diagonal equality is automatic;}
\qquad
\text{if }\nu(P)\text{ is odd, it forces }\bra{0_L}P\ket{0_L}=0.
}
\]

\paragraph{Concrete consequence for \(\mathrm{wt}(P)\le 2\).}
For weight \(\le 2\), the only diagonal constraints that \emph{do not} become tautologies are those with \(\nu(P)\) odd, i.e.
\[
\boxed{
\bra{0_L}Z_i\ket{0_L}=0,\quad
\bra{0_L}Y_i\ket{0_L}=0,\quad
\bra{0_L}X_iZ_j\ket{0_L}=0,\quad
\bra{0_L}X_iY_j\ket{0_L}=0
\quad (i\neq j),
}
\]
(and equivalently with indices swapped, e.g. \(Z_iX_j\)).
All other diagonal equalities for \(\mathrm{wt}(P)\le 2\) hold automatically under complementarity.

\subsubsection{LP/Z-type constraints simplify drastically under complementarity}

For Z-type Paulis, \(\nu(P)\) equals their weight (since they contain only \(Z\)'s).
Thus:
\begin{itemize}
\item For \(Z_i\): \(\nu(Z_i)=1\) is odd, so KL forces \(\bra{0_L}Z_i\ket{0_L}=0\), i.e.
\[
\boxed{
\sum_{s\in S_0(\vec a)} (-1)^{s_i}\,p^{(0)}_s = 0\quad (i=1,\dots,7).
}
\]
This is the \emph{perfect-balance} constraint: the marginal probability that qubit \(i\) is \(1\) is \(1/2\).
\item For \(Z_iZ_j\): \(\nu(Z_iZ_j)=2\) is even, so the diagonal KL equality is automatic; no additional linear constraints appear.
\end{itemize}
Moreover, if \(S_0(\vec a)\cap S_1(\vec a)=\varnothing\), then \(\bra{0_L}Z_i\ket{1_L}=0\)
and \(\bra{0_L}Z_iZ_j\ket{1_L}=0\) are automatic as well (pure support disjointness).

So, under \(\ket{1_L}=X^{\otimes 7}\ket{0_L}\) and disjoint SS blocks, the LP stage can be taken to mean:
\[
\boxed{
p^{(0)}_s\ge 0,\ \sum_{s\in S_0}p^{(0)}_s=1,\ \ \sum_{s\in S_0}(-1)^{s_i}p^{(0)}_s=0\ \ (i=1,\dots,7),
}
\]
and then define \(p^{(1)}_t:=p^{(0)}_{t\oplus 1^7}\).

\subsubsection{Off-diagonal KL constraints reduce to finite ``pairing sums'' on \texorpdfstring{$S_0$}{S0}}

With \(\ket{1_L}=X^{\otimes 7}\ket{0_L}\), the off-diagonal KL constraint becomes
\[
\bra{0_L}P\ket{1_L}=\bra{0_L}P X^{\otimes 7}\ket{0_L}=0.
\]
In coefficient form, using the symplectic formula and reindexing \(t=s\oplus 1^7\), one convenient expression is:
\[
\bra{0_L}P\ket{1_L}
=
i^{N_Y(P)}(-1)^{\nu(P)}
\sum_{t\in\{0,1\}^7}
(-1)^{z(P)\cdot t}\,
\overline{c_{t\oplus(1^7\oplus x(P))}}\,c_{t}.
\]
Equivalently, in magnitude--phase variables (writing \(c_t=\sqrt{p_t}e^{i\phi_t}\) for \(\ket{0_L}\)):
\[
\boxed{
\bra{0_L}P\ket{1_L}
=
i^{N_Y(P)}(-1)^{\nu(P)}
\sum_{t\in\{0,1\}^7}
(-1)^{z(P)\cdot t}\,
\sqrt{p_{t\oplus(1^7\oplus x(P))}}\,\sqrt{p_t}\,
e^{i(\phi_t-\phi_{t\oplus(1^7\oplus x(P))})}.
}
\]
Because SS enforces \(p_t=0\) outside \(S_0(\vec a)\), the sum effectively runs over those
\(t\in S_0(\vec a)\) for which also \(t\oplus(1^7\oplus x(P))\in S_0(\vec a)\).
Hence:

\paragraph{Automatic vanishing criterion.}
If
\[
S_0(\vec a)\ \cap\ \bigl(S_0(\vec a)\oplus(1^7\oplus x(P))\bigr)=\varnothing,
\]
then \(\bra{0_L}P\ket{1_L}=0\) holds automatically (the sum is empty / all terms are zero).
This is one precise way in which many Full-KL constraints are ``automatically satisfied'' once SS is fixed.

\paragraph{Summary of what remains nontrivial under complementarity.}
Assuming \(\ket{1_L}=X^{\otimes 7}\ket{0_L}\) and disjoint SS blocks, Full-KL for distance \(3\) reduces to:
\begin{itemize}
\item \textbf{LP (linear):} normalization of \(p\) on \(S_0\) and the 7 balance constraints \(\bra{0_L}Z_i\ket{0_L}=0\).
\item \textbf{Full-KL diagonal (nonlinear but fewer):} only the ``odd \(\nu(P)\)'' diagonals must be forced to zero:
\(\bra{0_L}Y_i\ket{0_L}=0\) and \(\bra{0_L}X_iZ_j\ket{0_L}=\bra{0_L}X_iY_j\ket{0_L}=0\) for \(i\neq j\).
\item \textbf{Full-KL off-diagonal (phase cancellations):} \(\bra{0_L}P\ket{1_L}=0\) for every \(P\in\mathcal{P}_{\le 2}\).
Many of these are automatic if the relevant SS-induced intersections are empty; the remaining ones are explicit finite sums
of the form shown above.
\end{itemize}

\medskip
\noindent
This structure is designed to be translated directly into a finite verification workflow:
define SS blocks \(S_0,S_1\), define amplitudes on \(S_0\), derive \(\ket{1_L}\) by complement,
and then check a finite list of explicit equalities indexed by \(\mathcal{P}_{\le 2}\).

\section{BD16 (transversal \(T\)) specialization}

For BD16 we take \(m=8\). A transversal implementation of the logical gate \(Z(-2\pi/8)\) corresponds to choosing residues \((b_0,b_1)=(0,7)\) (i.e., \(b_1=m-1\)).

In Sec.~6.3.1 of the cited paper (BD16: transversal \(T\)), the SS--LP stage returns the following candidate angle vectors \(\vec a\):
\[
\begin{aligned}
&(1, 2, 2, 2, 2, 3, 3),\ (1, 1, 2, 2, 3, 3, 3),\ (1, 1, 2, 2, 2, 3, 4),\ (1, 1, 2, 2, 2, 2, 5),\\
&(1, 1, 1, 2, 3, 3, 4),\ (1, 1, 1, 2, 2, 4, 4),\ (1, 1, 1, 2, 2, 3, 5),\ (1, 1, 1, 2, 2, 2, 6),\\
&(1, 1, 1, 1, 3, 4, 4),\ (1, 1, 1, 1, 3, 3, 5),\ (0, 1, 2, 2, 3, 3, 4),\ (0, 1, 1, 2, 3, 3, 5).
\end{aligned}
\]
Each such \(\vec a\) has passed the SS and LP filters (i.e., the \(Z\)-type KL constraints), and one must proceed to the Full-KL stage to determine which ones yield actual distance-3 codes.

%\section{Explicit example from Sec.~6.3.1: \texorpdfstring{$\vec a=(1,2,2,2,2,3,3)$}{a=(1,2,2,2,2,3,3)}}

\section{SS sets, LP solution space, and Full-KL verification for $\vec a=(1,2,2,2,2,3,3)$}

The cited paper provides a concrete \(((7,2,3))\) example corresponding to
\[
\vec a=(1,2,2,2,2,3,3),
\qquad
\lambda^\ast=\sqrt{\frac{21}{8}}.
\]
The logical states are given (up to two free phases \(\theta_1,\theta_2\in\mathbb{R}\)) by
\[
\begin{aligned}
4\ket{0_L}
=&\ e^{i\theta_1}\ket{0000000}
+\sqrt{3}\ket{0111100}
+\sqrt{3}\ket{1001110}
+\sqrt{2}\ket{1010101}
+\sqrt{2}\ket{1011001}
+e^{i\theta_2}\ket{1110010}
+2i\,\ket{0100011},\\
\ket{1_L}=&\ X^{\otimes 7}\ket{0_L}.
\end{aligned}
\]

Writing \(Z(\theta)=\mathrm{diag}(1,e^{i\theta})\), the transversal logical generators reported there are
\[
\overline{X}=X^{\otimes 7},
\]
and
\[
\overline{Z\!\left(-\frac{2\pi}{8}\right)}
=
Z\!\left(\frac{2\pi}{8}\right)\otimes
Z\!\left(\frac{4\pi}{8}\right)\otimes
Z\!\left(\frac{4\pi}{8}\right)\otimes
Z\!\left(\frac{4\pi}{8}\right)\otimes
Z\!\left(\frac{4\pi}{8}\right)\otimes
Z\!\left(\frac{6\pi}{8}\right)\otimes
Z\!\left(\frac{6\pi}{8}\right).
\]

The (Shor--Laflamme) weight enumerators for this code are
\[
A=\Big[
1,\ 0,\ \frac{21}{8},\ \frac{2}{8}-\frac{1}{8}a,\ 
\frac{215}{16}+\frac{18}{16}a+\frac{1}{16}b,\ 
\frac{87}{16}-\frac{42}{16}a-\frac{3}{16}b,\ 
\frac{635}{16}+\frac{38}{16}a+\frac{3}{16}b,\ 
\frac{25}{16}-\frac{12}{16}a-\frac{1}{16}b
\Big],
\]
\[
B=\Big[
1,\ 0,\ \frac{21}{8},\ 
\frac{441}{16}+\frac{10}{16}a+\frac{1}{16}b,\ 
\frac{352}{16}-3a-\frac{1}{4}b,\ 
\frac{1590}{16}+\frac{84}{16}a+\frac{6}{16}b,\ 
\frac{846}{16}-4a-\frac{1}{4}b,\ 
\frac{809}{16}+\frac{18}{16}a+\frac{1}{16}b
\Big],
\]
with parameters (to avoid confusion with the angle vector \(\vec a\), note these are scalar parameters)
\[
a=\cos(2\theta_1)+\cos(2\theta_2),
\qquad
b=\cos(2\theta_1-2\theta_2).
\]

In this section we work out the SS blocks $S_0,S_1$, the LP constraints and their general solution,
and then verify in detail that the two-parameter family
\[
\begin{aligned}
4\ket{0_L}
=&\ e^{i\theta_1}\ket{0000000}
+\sqrt{3}\ket{0111100}
+\sqrt{3}\ket{1001110}
+\sqrt{2}\ket{1010101}
+\sqrt{2}\ket{1011001}
+e^{i\theta_2}\ket{1110010}
+2i\,\ket{0100011},\\
\ket{1_L}=&\ X^{\otimes 7}\ket{0_L}
\end{aligned}
\]
satisfies the \emph{full} distance-$3$ Knill--Laflamme conditions.

Throughout, $m=8$ and we use residues $(b_0,b_1)=(0,7)$.

\subsection{Stage 1 (SS): explicit residue blocks $S_0(\vec a)$ and $S_1(\vec a)$}

Fix
\[
\vec a=(1,2,2,2,2,3,3)\in\mathbb{Z}_8^7,
\qquad
\mathrm{res}_{\vec a}(s):=\sum_{j=1}^7 a_j s_j\pmod 8.
\]
We take residues $(b_0,b_1)=(0,7)$, so
\[
S_k(\vec a)=\{\,s\in\{0,1\}^7:\ \mathrm{res}_{\vec a}(s)\equiv b_k\ (\mathrm{mod}\ 8)\,\}.
\]

\paragraph{Complementary compatibility.}
Since
\[
\sum_{j=1}^7 a_j = 1+2+2+2+2+3+3=15\equiv 7\pmod 8,
\]
the complementary convention $\ket{1_L}=X^{\otimes 7}\ket{0_L}$ forces
\[
b_1\equiv b_0+\sum_{j=1}^7 a_j\equiv 7\pmod 8,
\qquad
S_1(\vec a)=\overline{S_0(\vec a)}.
\]

\paragraph{The blocks in full.}
A direct enumeration gives
\[
\begin{aligned}
S_0(\vec a)=\{&
0000000,\ 0000111,\ 0001011,\ 0010011,\ 0100011,\ 0111100,\\
&1001101,\ 1001110,\ 1010101,\ 1010110,\ 1011001,\ 1011010,\\
&1100101,\ 1100110,\ 1101001,\ 1101010,\ 1110001,\ 1110010
\},
\end{aligned}
\]
and (by complement)
\[
\begin{aligned}
S_1(\vec a)=\{&
1111111,\ 1111000,\ 1110100,\ 1101100,\ 1011100,\ 1000011,\\
&0110010,\ 0110001,\ 0101010,\ 0101001,\ 0100110,\ 0100101,\\
&0011010,\ 0011001,\ 0010110,\ 0010101,\ 0001110,\ 0001101
\}.
\end{aligned}
\]

\paragraph{Transversal diagonal action.}
Let $\omega=e^{2\pi i/8}$ and
\[
D(\vec a)=\bigotimes_{j=1}^7\left(\ket{0}\!\bra{0}+\omega^{a_j}\ket{1}\!\bra{1}\right)
=
Z\!\left(\tfrac{2\pi}{8}\right)\otimes
Z\!\left(\tfrac{4\pi}{8}\right)^{\otimes 4}\otimes
Z\!\left(\tfrac{6\pi}{8}\right)^{\otimes 2}.
\]
If $\ket{0_L}$ is supported on $S_0(\vec a)$ then $D(\vec a)$ acts as $\omega^{0}=1$ on $\ket{0_L}$, while
$D(\vec a)$ acts as $\omega^{7}=e^{-2\pi i/8}$ on $\ket{1_L}$ (supported on $S_1(\vec a)$), hence the induced
logical gate is $\overline{Z(-2\pi/8)}=\mathrm{diag}(1,e^{-2\pi i/8})$ in the logical basis.

\subsection{Stage 2 (LP): explicit $Z$-type constraints and general solution}

Write
\[
\ket{0_L}=\sum_{s\in S_0(\vec a)} c_s \ket{s},
\qquad
c_s=\sqrt{p_s}\,e^{i\phi_s},
\qquad
p_s\ge 0,\ \ \sum_{s\in S_0}p_s=1.
\]

\paragraph{LP equalities (distance 3).}
Under the complementary convention $\ket{1_L}=X^{\otimes 7}\ket{0_L}$, the diagonal KL equalities
for all \emph{odd}-$Z$ Paulis reduce to vanishing:
\[
\bra{0_L}Z_i\ket{0_L}=0\qquad (i=1,\dots,7).
\]
Since $Z_i\ket{s}=(-1)^{s_i}\ket{s}$, these are linear equations in $\{p_s\}$:
\[
\boxed{\ \sum_{s\in S_0(\vec a)} (-1)^{s_i}\,p_s = 0\ \ (i=1,\dots,7),\qquad \sum_{s\in S_0(\vec a)}p_s=1.\ }
\]
Equivalently, for each $i$,
\[
\sum_{s\in S_0:\ s_i=1} p_s=\sum_{s\in S_0:\ s_i=0} p_s=\frac12.
\]

For clarity, the seven ``$\frac12$'' equations are (one per qubit, listing the $s\in S_0$ with $s_i=1$):

\begin{align*}
&\text{(qubit 1)} &&
p_{1001101}+p_{1001110}+p_{1010101}+p_{1010110}+p_{1011001}+p_{1011010}
+p_{1100101}+p_{1100110}+p_{1101001}+p_{1101010}+p_{1110001}+p_{1110010}
=\frac12,\\
&\text{(qubit 2)} &&
p_{0100011}+p_{0111100}+p_{1100101}+p_{1100110}+p_{1101001}+p_{1101010}+p_{1110001}+p_{1110010}
=\frac12,\\
&\text{(qubit 3)} &&
p_{0010011}+p_{0111100}+p_{1010101}+p_{1010110}+p_{1011001}+p_{1011010}+p_{1110001}+p_{1110010}
=\frac12,\\
&\text{(qubit 4)} &&
p_{0001011}+p_{0111100}+p_{1001101}+p_{1001110}+p_{1011001}+p_{1011010}+p_{1101001}+p_{1101010}
=\frac12,\\
&\text{(qubit 5)} &&
p_{0000111}+p_{0111100}+p_{1001101}+p_{1001110}+p_{1010101}+p_{1010110}+p_{1100101}+p_{1100110}
=\frac12,\\
&\text{(qubit 6)} &&
p_{0000111}+p_{0001011}+p_{0010011}+p_{0100011}+p_{1001110}+p_{1010110}+p_{1011010}+p_{1100110}+p_{1101010}+p_{1110010}
=\frac12,\\
&\text{(qubit 7)} &&
p_{0000111}+p_{0001011}+p_{0010011}+p_{0100011}+p_{1001101}+p_{1010101}+p_{1011001}+p_{1100101}+p_{1101001}+p_{1110001}
=\frac12.
\end{align*}

\paragraph{General affine solution of the LP equalities.}
The above eight linear equalities on the $18$ variables $\{p_s:s\in S_0\}$ have a $10$-dimensional
solution space (before imposing inequalities $p_s\ge 0$).
A convenient parametrization is obtained by choosing the following $10$ free variables:
\[
\begin{aligned}
&u_1=p_{1010101},\quad u_2=p_{1010110},\quad u_3=p_{1011001},\quad u_4=p_{1011010},\\
&u_5=p_{1100101},\quad u_6=p_{1100110},\quad u_7=p_{1101001},\quad u_8=p_{1101010},\\
&u_9=p_{1110001},\quad u_{10}=p_{1110010}.
\end{aligned}
\]
Then the remaining probabilities are forced to be:
\begin{align*}
p_{0000000}&=\frac{1}{16},\\
p_{0111100}&=\frac{3}{16},\\
p_{0000111}&=u_3+u_4+u_7+u_8+u_9+u_{10}-\frac{3}{16},\\
p_{0001011}&=u_1+u_2+u_5+u_6+u_9+u_{10}-\frac{3}{16},\\
p_{0010011}&=\frac{5}{16}-(u_1+u_2+u_3+u_4+u_9+u_{10}),\\
p_{0100011}&=\frac{5}{16}-(u_5+u_6+u_7+u_8+u_9+u_{10}),\\
p_{1001101}&=\frac{1}{4}-(u_1+u_3+u_5+u_7+u_9),\\
p_{1001110}&=\frac{1}{4}-(u_2+u_4+u_6+u_8+u_{10}).
\end{align*}
Imposing $p_s\ge 0$ cuts out a polytope in the $(u_1,\dots,u_{10})$ parameters.

\subsection{Stage 3 (Full-KL): complementary simplification and support-intersection method}

Let $\mathcal{P}_{\le 2}$ be the set of all Hermitian Pauli strings on $7$ qubits of weight $\le 2$.
Distance $3$ is equivalent to:
\[
\boxed{
\forall\,P\in\mathcal{P}_{\le 2},\quad
\bra{0_L}P\ket{1_L}=0,\qquad \bra{0_L}P\ket{0_L}=\bra{1_L}P\ket{1_L}.
}
\]

\paragraph{Complementary simplification of the diagonal constraints.}
Write a Pauli as $P(x,z)$ with $x,z\in\{0,1\}^7$ indicating which tensor factors contain $X/Y$
and which contain $Z/Y$, respectively (so $Z$-components are counted by $z$).
Then
\[
X^{\otimes 7} P(x,z) X^{\otimes 7} = (-1)^{\mathrm{wt}(z)} P(x,z),
\]
hence
\[
\bra{1_L}P(x,z)\ket{1_L}=(-1)^{\mathrm{wt}(z)}\bra{0_L}P(x,z)\ket{0_L}.
\]
Therefore the diagonal KL equality is automatic when $\mathrm{wt}(z)$ is even, and becomes a vanishing
condition when $\mathrm{wt}(z)$ is odd:
\[
\boxed{\ \mathrm{wt}(z)\ \text{odd}\ \Longrightarrow\ \bra{0_L}P(x,z)\ket{0_L}=0.\ }
\]

\paragraph{Matrix elements as sums over support intersections.}
Let $T:=\mathrm{supp}(\ket{0_L})\subseteq\{0,1\}^7$, so $\mathrm{supp}(\ket{1_L})=\overline{T}$.
For $P(x,z)$ one has the (finite) sum formulas
\[
\bra{0_L}P(x,z)\ket{0_L}
=i^{N_Y(x,z)}\sum_{s\in T\cap (T\oplus x)}(-1)^{z\cdot s}\,\overline{c_{s\oplus x}}\,c_s,
\]
\[
\bra{0_L}P(x,z)\ket{1_L}
=i^{N_Y(x,z)}\sum_{s\in \overline{T}\cap (T\oplus x)}(-1)^{z\cdot s}\,\overline{c_{s\oplus x}}\,c_{\overline{s}},
\]
where $N_Y(x,z)=\sum_i x_i z_i$ is the number of $Y$ factors.

Thus \emph{most} KL constraints are automatic whenever the relevant intersections are empty; the remaining
constraints reduce to short explicit sums over those intersections.

\subsection{Why the two-parameter family satisfies Full-KL (distance 3)}

For the reported solution, the nonzero amplitudes of $\ket{0_L}$ are:
\[
c_{0000000}=\frac{e^{i\theta_1}}{4},\quad
c_{0111100}=\frac{\sqrt3}{4},\quad
c_{1001110}=\frac{\sqrt3}{4},\quad
c_{1010101}=\frac{\sqrt2}{4},\quad
c_{1011001}=\frac{\sqrt2}{4},\quad
c_{1110010}=\frac{e^{i\theta_2}}{4},\quad
c_{0100011}=\frac{i}{2},
\]
and $c_s=0$ for all other $s\in S_0(\vec a)$.
In particular,
\[
T=\{0000000,\ 0111100,\ 1001110,\ 1010101,\ 1011001,\ 1110010,\ 0100011\}
\subseteq S_0(\vec a),
\]
so the SS condition is satisfied. Also $T\cap\overline{T}=\varnothing$, so $\braket{0_L|1_L}=0$.

\paragraph{(1) LP / $Z_i$ constraints.}
The squared magnitudes (LP variables) are
\[
p_{0000000}=\frac{1}{16},\ 
p_{0111100}=\frac{3}{16},\ 
p_{1001110}=\frac{3}{16},\ 
p_{1010101}=\frac{2}{16},\ 
p_{1011001}=\frac{2}{16},\ 
p_{1110010}=\frac{1}{16},\ 
p_{0100011}=\frac{4}{16},
\]
and all other $p_s$ on $S_0(\vec a)$ vanish. These sum to $1$.

Moreover, for each qubit $i$ one checks directly that the total probability on $s_i=1$ is $1/2$.
For example, on qubit~1,
\[
p_{1001110}+p_{1010101}+p_{1011001}+p_{1110010}
=\frac{3+2+2+1}{16}=\frac12,
\]
and similarly for $i=2,\dots,7$. Hence $\bra{0_L}Z_i\ket{0_L}=0$ for all $i$, so all odd-$Z$ \emph{diagonal}
constraints of the form $Z_i$ are satisfied.

\paragraph{(2) All diagonal constraints with $\mathrm{wt}(z)$ odd.}
By the complementary simplification, it suffices to show $\bra{0_L}P(x,z)\ket{0_L}=0$ when $\mathrm{wt}(z)$ is odd.

\begin{itemize}
\item If $x$ has Hamming weight $1$ (single bit flip) then $T\cap(T\oplus x)=\varnothing$, so
$\bra{0_L}P(x,z)\ket{0_L}=0$ automatically for all such Paulis.
\item The only nontrivial two-bit-flip overlap within $T$ is
\[
T\cap (T\oplus (e_4\oplus e_5))=\{1010101,\ 1011001\},
\qquad
(1010101)\oplus(e_4\oplus e_5)=1011001.
\]
For mixed Paulis on qubits $(4,5)$ with exactly one $Z$-component (e.g.\ $X_4Y_5$ or $Y_4X_5$),
$\mathrm{wt}(z)=1$ is odd and the diagonal sum has exactly two terms which cancel because
$c_{1010101}=c_{1011001}=\sqrt2/4$ is real and the $(-1)^{z\cdot s}$ sign flips between the two strings
(which differ on exactly one of the $z$-supported positions). Thus $\bra{0_L}P(x,z)\ket{0_L}=0$ for these cases.
\end{itemize}

This exhausts all odd-$\mathrm{wt}(z)$ diagonal constraints at weight $\le 2$.

\paragraph{(3) All off-diagonal constraints $\bra{0_L}P\ket{1_L}=0$.}
Using the intersection formula, it suffices to analyze $\overline{T}\cap (T\oplus x)$.

\begin{itemize}
\item If $x=0$ (purely diagonal Paulis $Z^z$), then $\bra{0_L}Z^z\ket{1_L}=0$ because $T\cap\overline{T}=\varnothing$.
\item If $x$ has Hamming weight $1$, then $\overline{T}\cap(T\oplus x)=\varnothing$, so all weight-$1$ off-diagonals vanish.
\item For two-bit flips, the only $x$ for which $\overline{T}\cap(T\oplus x)\neq\varnothing$ are:
\[
\begin{aligned}
&x=e_1\oplus e_2:\quad \overline{T}\cap (T\oplus x)=\{0111100,\ 0100011\},\\
&x=e_3\oplus e_6:\quad \overline{T}\cap (T\oplus x)=\{1001110,\ 0100011\},\\
&x=e_4\oplus e_7:\quad \overline{T}\cap (T\oplus x)=\{1010101,\ 0100011\},\\
&x=e_5\oplus e_7:\quad \overline{T}\cap (T\oplus x)=\{1011001,\ 0100011\}.
\end{aligned}
\]
In each case the off-diagonal sum has exactly two terms, and it vanishes for \emph{all} choices of $z$ supported
on the same flipped qubits (i.e.\ for $X\!X$, $X\!Y$, $Y\!X$, and $Y\!Y$ on that pair), because:
\begin{enumerate}
\item one coefficient is real (one of $\sqrt3/4$ or $\sqrt2/4$) while $c_{0100011}=i/2$ is purely imaginary, so
\[
\overline{c_{\mathrm{real}}}\,c_{0100011}+\overline{c_{0100011}}\,c_{\mathrm{real}}=0,
\]
\item and the $(-1)^{z\cdot s}$ signs match on the two contributing terms since the two intermediate strings
$s$ appearing in the sum agree on the $z$-supported positions (the positions where a $Y$ contributes a $Z$-component).
\end{enumerate}
Hence $\bra{0_L}P(x,z)\ket{1_L}=0$ for all weight-$2$ Paulis as well.
\end{itemize}

\paragraph{Conclusion.}
We have shown:
\begin{itemize}
\item normalization and orthogonality;
\item all odd-$\mathrm{wt}(z)$ diagonal constraints (including all $Z_i$);
\item all off-diagonal constraints for every $P\in\mathcal{P}_{\le 2}$.
\end{itemize}
Therefore the logical matrix $M(P)$ is proportional to the identity for every Pauli string $P$ of weight $\le 2$,
so the code has distance $3$.

\paragraph{Why $\theta_1,\theta_2$ are free.}
The only basis states carrying the phases $e^{i\theta_1}$ and $e^{i\theta_2}$ are $0000000$ and $1110010$.
In the above intersection analysis, neither of these strings participates in any nontrivial overlap contributing to
weight-$\le 2$ off-diagonal sums or odd-$\mathrm{wt}(z)$ diagonal sums. Hence all KL constraints are independent of
$\theta_1,\theta_2$, yielding a genuine two-parameter family.

\subsection{Intersection test for full KL: only $25$ Paulis to check in the $(7,2,3)$ example}
\label{subsec:intersection-test-25}

We use the standard Pauli parametrization (up to global phase)
\[
P(x,z)\;:=\; i^{x\cdot z} X^x Z^z,\qquad x,z\in\{0,1\}^7,
\]
so that on qubit $i$ we have $(x_i,z_i)=(0,0)\mapsto I$, $(1,0)\mapsto X$, $(0,1)\mapsto Z$, $(1,1)\mapsto Y$.
Let
\[
\ket{0_L}=\sum_{s\in T} c_s \ket{s},
\qquad
\ket{1_L}=X^{\otimes 7}\ket{0_L}=\sum_{s\in T} c_s \ket{\overline{s}},
\]
where $T:=\mathrm{supp}(\ket{0_L})\subseteq\{0,1\}^7$ and $\overline{s}:=s\oplus 1^7$ is the bitwise complement.
In the concrete code,
\[
T=\{0000000,\ 0111100,\ 1001110,\ 1010101,\ 1011001,\ 1110010,\ 0100011\}.
\]

\paragraph{Finite-sum formulas.}
For any $P(x,z)$, straightforward basis expansion gives
\begin{align}
\bra{0_L}P(x,z)\ket{0_L}
&=
i^{x\cdot z}\!\!\sum_{s\in T\cap(T\oplus x)}\!\!
(-1)^{z\cdot s}\ \overline{c_{s\oplus x}}\,c_s,
\label{eq:diag-intersection}\\
\bra{0_L}P(x,z)\ket{1_L}
&=
i^{x\cdot z}(-1)^{z\cdot x}\!\!\sum_{s\in T\cap(T\oplus \Delta(x))}\!\!
(-1)^{z\cdot s}\ \overline{c_{s\oplus \Delta(x)}}\,c_s,
\qquad \Delta(x):=1^7\oplus x .
\label{eq:offdiag-intersection}
\end{align}
Thus, if the relevant intersection is empty, the corresponding matrix element vanishes \emph{automatically}.

\paragraph{Complementary simplification of diagonal KL.}
Since $X^{\otimes 7} P(x,z) X^{\otimes 7}=(-1)^{\mathrm{wt}(z)}P(x,z)$, we have
\[
\bra{1_L}P(x,z)\ket{1_L}=(-1)^{\mathrm{wt}(z)}\bra{0_L}P(x,z)\ket{0_L}.
\]
Hence the diagonal KL equality $\bra{0_L}P\ket{0_L}=\bra{1_L}P\ket{1_L}$ is automatic for even $\mathrm{wt}(z)$,
and reduces to a \emph{vanishing} condition for odd $\mathrm{wt}(z)$:
\begin{equation}
\mathrm{wt}(z)\ \text{odd}\quad\Longrightarrow\quad \bra{0_L}P(x,z)\ket{0_L}=0.
\label{eq:oddz-vanish}
\end{equation}

\begin{lemma}[Diagonal intersection screen for $T$]
\label{lem:diag-screen}
For the support set $T$ above, among all $x$ with $\mathrm{wt}(x)\le 2$ we have
\[
T\cap(T\oplus x)\neq\emptyset
\quad\Longleftrightarrow\quad
x\in\{0,\ e_4\oplus e_5\}.
\]
Moreover,
\[
T\cap(T\oplus(e_4\oplus e_5))=\{1010101,\ 1011001\}.
\]
\end{lemma}

\begin{proof}
One checks directly that no element of $T$ has a Hamming-$1$ neighbor in $T$, so $T\cap(T\oplus e_i)=\emptyset$ for all $i$.
Among Hamming-$2$ flips, the only pair of strings in $T$ at distance $2$ is
$1010101\leftrightarrow 1011001$, which differ exactly on qubits $4$ and $5$.
All other pairs in $T$ have distance $\ge 3$, hence yield empty intersections for $\mathrm{wt}(x)\le 2$.
\end{proof}

\begin{lemma}[Off-diagonal intersection screen for $T$]
\label{lem:offdiag-screen}
For the same $T$, among all $x$ with $\mathrm{wt}(x)\le 2$ we have
\[
T\cap(T\oplus\Delta(x))\neq\emptyset
\quad\Longleftrightarrow\quad
x\in\{e_1\oplus e_2,\ e_3\oplus e_6,\ e_4\oplus e_7,\ e_5\oplus e_7\}.
\]
In these four cases, the intersection has size $2$, namely
\[
\begin{aligned}
T\cap\bigl(T\oplus\Delta(e_1\oplus e_2)\bigr)&=\{0111100,\ 0100011\},\\
T\cap\bigl(T\oplus\Delta(e_3\oplus e_6)\bigr)&=\{1001110,\ 0100011\},\\
T\cap\bigl(T\oplus\Delta(e_4\oplus e_7)\bigr)&=\{1010101,\ 0100011\},\\
T\cap\bigl(T\oplus\Delta(e_5\oplus e_7)\bigr)&=\{1011001,\ 0100011\}.
\end{aligned}
\]
\end{lemma}

\begin{proof}
This is a finite check: for each $x$ with $\mathrm{wt}(x)\le 2$, compute $\Delta(x)=1^7\oplus x$ and test whether
there exist $s,s'\in T$ with $s'=s\oplus\Delta(x)$.
Exactly the four displayed choices admit such pairs, and each yields precisely one undirected edge inside $T$,
hence an intersection of size $2$.
\end{proof}

\begin{proposition}[Only $25$ Pauli constraints can be nontrivial]
\label{prop:25-paulis}
Consider the full distance-$3$ KL conditions for all Pauli operators of weight $\le 2$.
For the above $T$ and the complementary convention $\ket{1_L}=X^{\otimes 7}\ket{0_L}$, all KL constraints are
\emph{automatic} except the following $25$ operators:
\begin{enumerate}
\item (Odd-$z$ diagonal constraints at $x=0$) the seven single-$Z$ operators
\[
Z_1,\ Z_2,\ Z_3,\ Z_4,\ Z_5,\ Z_6,\ Z_7,
\]
which enforce $\bra{0_L}Z_i\ket{0_L}=0$ for all $i$ (equivalently, each bit-position is balanced in $|c_s|^2$).
\item (Odd-$z$ diagonal constraints at $x=e_4\oplus e_5$) the two mixed $XY$ operators
\[
X_4Y_5,\qquad Y_4X_5,
\]
which enforce $\bra{0_L}X_4Y_5\ket{0_L}=0$ and $\bra{0_L}Y_4X_5\ket{0_L}=0$.
\item (Off-diagonal constraints) the sixteen two-qubit $\{X,Y\}^{\otimes 2}$ operators on the four pairs
\[
(1,2),\ (3,6),\ (4,7),\ (5,7),
\]
i.e.
\[
\{X_iX_j,\ X_iY_j,\ Y_iX_j,\ Y_iY_j\}\quad \text{for each of the above pairs }(i,j),
\]
which enforce $\bra{0_L}P\ket{1_L}=0$ for all such $P$.
\end{enumerate}
Every other Pauli of weight $\le 2$ yields KL equations that hold automatically.
\end{proposition}

\begin{proof}
We separate diagonal and off-diagonal KL conditions.

\smallskip
\noindent\emph{Diagonal.}
By the complementary relation, the diagonal equality is automatic when $\mathrm{wt}(z)$ is even, and reduces to the vanishing
condition \eqref{eq:oddz-vanish} when $\mathrm{wt}(z)$ is odd.
For $\mathrm{wt}(P)\le 2$, odd $\mathrm{wt}(z)$ occurs only for
\[
Z_i\ (x=0,z=e_i),\qquad\text{and}\qquad X_iY_j,\ Y_iX_j\ (x=e_i\oplus e_j,\ z\in\{e_i,e_j\}),
\]
with $i\neq j$.
Now apply Lemma~\ref{lem:diag-screen}:
\begin{itemize}
\item If $x=e_i$ or $x=e_i\oplus e_j$ with $(i,j)\neq(4,5)$, then $T\cap(T\oplus x)=\emptyset$ and
$\bra{0_L}P(x,z)\ket{0_L}=0$ holds automatically by \eqref{eq:diag-intersection}.
\item If $x=0$, the only odd-$z$ cases are exactly $Z_i$ ($i=1,\dots,7$), giving $7$ potentially nontrivial constraints.
\item If $x=e_4\oplus e_5$, the odd-$z$ cases are exactly $X_4Y_5$ and $Y_4X_5$, giving $2$ potentially nontrivial constraints.
\end{itemize}

\smallskip
\noindent\emph{Off-diagonal.}
For weight $\le 2$, $\bra{0_L}P(x,z)\ket{1_L}$ can be nonzero only if
$T\cap(T\oplus\Delta(x))\neq\emptyset$ by \eqref{eq:offdiag-intersection}.
By Lemma~\ref{lem:offdiag-screen} this happens only for
$x\in\{e_1\oplus e_2,\,e_3\oplus e_6,\,e_4\oplus e_7,\,e_5\oplus e_7\}$.
Since $\mathrm{wt}(x)=2$ in each case, any Pauli of weight $\le 2$ with this $x$ must have $z$ supported only on the same two qubits,
so it is one of $\{X_iX_j,X_iY_j,Y_iX_j,Y_iY_j\}$, yielding $4$ operators per pair and $16$ total.

\smallskip
Combining diagonal ($7+2$) and off-diagonal ($16$) cases gives $25$ total potentially nontrivial constraints.
All other constraints are automatic either by even-$\mathrm{wt}(z)$ (diagonal) or by empty intersections (diagonal/off-diagonal).
\end{proof}

\begin{remark}[Where $\theta_1,\theta_2$ drop out]
In the explicit code, the only phased basis states are $0000000$ and $1110010$.
Neither participates in the nontrivial intersections of Lemmas~\ref{lem:diag-screen}--\ref{lem:offdiag-screen},
so none of the $25$ constraints depends on $\theta_1$ or $\theta_2$.
\end{remark}

% (Context/citation for the SSLP/KL setup and complementary simplification.)
% See the SSLP framework and KL discussion in the cited manuscript.
% :contentReference[oaicite:0]{index=0}


\section{Exact BD16 examples for \(\vec a=(1,1,2,2,2,2,5)\)}

We now explain in detail \emph{exact} \(((7,2,3))\) codes with transversal group BD16 and angle vector
\[
\vec a=(1,1,2,2,2,2,5)\in\mathbb{Z}_8^7,\qquad (b_0,b_1)=(0,7).
\]

\subsection{An example}

\subsubsection{Logical states}

Define \(\ket{0_L}\) by the following \emph{exact} computational-basis expansion:
\[
\begin{aligned}
\ket{0_L}
=&\ \frac{1}{4}\ket{0000000}
-\frac{\sqrt{3}}{4}\ket{0011110}\\
&+\frac{i\sqrt{2}}{4}\Big(
\ket{0100101}-\ket{0110001}-\ket{1000011}+\ket{1001001}
\Big)\\
&+\frac{i}{4}\Big(
\ket{1101110}+\ket{1110110}+\ket{1111010}+\ket{1111100}
\Big).
\end{aligned}
\]
Let \(\overline{s}\) denote bitwise complement. Define
\[
\ket{1_L}:=X^{\otimes 7}\ket{0_L},
\]
so each basis string \(\ket{s}\) in \(\ket{0_L}\) is mapped to \(\ket{\overline{s}}\) with the same coefficient.

\subsubsection{Why the coefficients are ``square roots of rationals''}

All amplitudes have the form \(\frac{1}{4}\times \{1,\sqrt{2},\sqrt{3}\}\) times a root of unity \(\{1,i,-1,-i\}\).
Hence \(|c_s|^2\in \{1/16,2/16,3/16\}\subset \mathbb{Q}\), and all KL checks can be done exactly in \(\mathbb{Q}(\sqrt{2},\sqrt{3},i)\).

\subsubsection{Normalization and orthogonality}

Normalization is immediate from the squared magnitudes:
\[
\|\ket{0_L}\|^2
=
\frac{1}{16}
+\frac{3}{16}
+4\cdot \frac{2}{16}
+4\cdot \frac{1}{16}
=1.
\]
Orthogonality \(\braket{0_L|1_L}=0\) holds because the supports are disjoint: if \(\ket{s}\) appears in \(\ket{0_L}\), then \(\ket{\overline{s}}\) appears in \(\ket{1_L}\), and no basis string appears in both.

\subsubsection{Support as a residue class: the SS condition}

For \(\vec a=(1,1,2,2,2,2,5)\) and \(m=8\), define the residue map
\[
\mathrm{res}(s):=\sum_{j=1}^7 a_js_j \ (\mathrm{mod}\ 8).
\]
A direct check shows:
\[
\mathrm{res}(s)=0\ \text{for every basis string } s \text{ appearing in } \ket{0_L},
\]
and consequently
\[
\mathrm{res}(\overline{s})\equiv \left(\sum_{j=1}^7 a_j\right)-\mathrm{res}(s)\equiv 7\ (\mathrm{mod}\ 8)
\quad\text{for every basis string } \overline{s}\text{ in }\ket{1_L}.
\]
Thus \(\mathrm{supp}(\ket{0_L})\subseteq S_0(\vec a)\) and \(\mathrm{supp}(\ket{1_L})\subseteq S_1(\vec a)\) with \((b_0,b_1)=(0,7)\), exactly as required in the SS stage.

\subsubsection{Transversal logical gate (BD16)}

Define the transversal diagonal gate
\[
D(\vec a)=Z(\tfrac{\pi}{4})\otimes Z(\tfrac{\pi}{4})\otimes Z(\tfrac{\pi}{2})^{\otimes 4}\otimes Z(\tfrac{5\pi}{4}),
\qquad Z(\theta)=\mathrm{diag}(1,e^{i\theta}).
\]
On each computational basis state \(\ket{s}\),
\[
D(\vec a)\ket{s}=e^{2\pi i\,\mathrm{res}(s)/8}\ket{s}.
\]
Hence
\[
D(\vec a)\ket{0_L}=\ket{0_L},\qquad
D(\vec a)\ket{1_L}=e^{2\pi i\cdot 7/8}\ket{1_L}=e^{-i\pi/4}\ket{1_L},
\]
so the induced logical action is a \(Z\)-axis rotation by \(-\pi/4\), i.e., transversal \(T\)-type behavior (BD16).

\subsubsection{Full KL verification idea}

Because the code corrects all single-qubit errors, it suffices to verify the KL relations for all Pauli operators of weight \(\le 2\):
\[
\bra{0_L}P\ket{1_L}=0,\qquad
\bra{0_L}P\ket{0_L}=\bra{1_L}P\ket{1_L},
\qquad \forall\, P\in \{I,X,Y,Z\}^{\otimes 7},\ \mathrm{wt}(P)\le 2.
\]
The nontrivial part is the \(X/Y\) sector: \(X_i\) and \(Y_i\) create Hamming-1 transitions inside the residue class graph, and the phases in \(\ket{0_L}\) are tuned so that all such contributions cancel exactly.

\subsection{Exact verification in SymPy (including weight enumerators and \(\lambda^\ast\))}

The following script encodes \(\ket{0_L}\) in a dictionary (your exact coefficients), defines \(\ket{1_L}=X^{\otimes 7}\ket{0_L}\), and then checks KL for all Pauli strings of weight \(\le 2\). It also computes the quantum weight enumerators \(A,B\) and a natural \(\lambda^\ast\) (see below).

\begin{lstlisting}[language=Python]
import sympy as sp
from sympy import Rational, sqrt, I
import itertools

# Define amplitudes for |0_L> with 1/4 scaling
c0 = {
 "0000000": Rational(1,4),
 "0011110": -sqrt(3)/4,
 "0100101": I*sqrt(2)/4,
 "0110001": -I*sqrt(2)/4,
 "1000011": -I*sqrt(2)/4,
 "1001001": I*sqrt(2)/4,
 "1101110": I/4,
 "1110110": I/4,
 "1111010": I/4,
 "1111100": I/4,
}

# |1_L> = X^{\otimes 7}|0_L> = bitwise complement of each basis string, same coefficient
def complement(s):
    return "".join("1" if b=="0" else "0" for b in s)
c1 = { complement(s):coef for s,coef in c0.items() }

# Pauli action on a single bit
# Y|0> = i|1>, Y|1> = -i|0>
def apply_single_pauli(p, bit):
    if p=='I': return sp.Integer(1), bit
    if p=='X': return sp.Integer(1), ('1' if bit=='0' else '0')
    if p=='Z': return (sp.Integer(1) if bit=='0' else -sp.Integer(1)), bit
    if p=='Y':
        if bit=='0': return I, '1'
        else: return -I, '0'
    raise ValueError(p)

def apply_pauli_string(P, s):
    phase = sp.Integer(1)
    out = list(s)
    for i,p in enumerate(P):
        ph, nb = apply_single_pauli(p, out[i])
        phase *= ph
        out[i] = nb
    return phase, "".join(out)

def inner(cbra, P, cket):
    tot = sp.Integer(0)
    for s, amp in cket.items():
        ph, t = apply_pauli_string(P, s)
        if t in cbra:
            tot += sp.conjugate(cbra[t]) * (ph*amp)
    return sp.simplify(tot)

# KL check for all Pauli strings of weight <= 2
paulis = ['I','X','Y','Z']
def wt(P): return sum(1 for p in P if p!='I')

viol = []
for P in itertools.product(paulis, repeat=7):
    P = "".join(P)
    if wt(P) <= 2:
        off = sp.simplify(inner(c0,P,c1))           # <0|P|1>
        d0  = sp.simplify(inner(c0,P,c0))           # <0|P|0>
        d1  = sp.simplify(inner(c1,P,c1))           # <1|P|1>
        if off != 0 or sp.simplify(d0-d1) != 0:
            viol.append((P, off, sp.simplify(d0-d1)))

assert len(viol) == 0

# Weight enumerators A_j and B_j (Rains-type definitions)
# A_j = (1/K^2) sum_{wt(E)=j} |Tr(PiE)|^2
# B_j = (1/K)   sum_{wt(E)=j} Tr(Pi E Pi E^\dagger)
K = 2
A = [sp.Integer(0) for _ in range(8)]
B = [sp.Integer(0) for _ in range(8)]

def mat_elements(P):
    m00 = inner(c0,P,c0)
    m11 = inner(c1,P,c1)
    m01 = inner(c0,P,c1)
    m10 = inner(c1,P,c0)
    return m00,m01,m10,m11

for P in itertools.product(paulis, repeat=7):
    P = "".join(P)
    j = wt(P)
    m00,m01,m10,m11 = mat_elements(P)
    tr = sp.simplify(m00+m11)  # Tr(PiE) = <0|E|0> + <1|E|1>
    A[j] += sp.simplify(tr*sp.conjugate(tr))
    B[j] += sp.simplify(m00*sp.conjugate(m00) + m01*sp.conjugate(m01)
                        + m10*sp.conjugate(m10) + m11*sp.conjugate(m11))

A = [sp.simplify(a/(K**2)) for a in A]
B = [sp.simplify(b/K) for b in B]

# Signature norm for d=3 using E_D = {Paulis of weight 1 or 2}:
# lambda*^2 = A_1 + A_2
lambda_sq = sp.simplify(A[1] + A[2])
lambda_star = sp.simplify(sp.sqrt(lambda_sq))

print("A =", A)
print("B =", B)
print("lambda*^2 =", lambda_sq)
print("lambda* =", lambda_star)
\end{lstlisting}

\subsection{Weight enumerators and \(\lambda^\ast\) for this example}

Running the script yields:
\[
A=(A_0,\dots,A_7)=\left(1,\ 0,\ \frac{33}{16},\ 0,\ \frac{135}{8},\ 0,\ \frac{705}{16},\ 0\right),
\]
\[
B=(B_0,\dots,B_7)=\left(1,\ 0,\ \frac{33}{16},\ \frac{435}{16},\ \frac{135}{8},\ \frac{909}{8},\ \frac{705}{16},\ \frac{819}{16}\right).
\]
Equivalently, as polynomials:
\[
A(z)=1+\frac{33}{16}z^2+\frac{135}{8}z^4+\frac{705}{16}z^6,
\]
\[
B(z)=1+\frac{33}{16}z^2+\frac{435}{16}z^3+\frac{135}{8}z^4+\frac{909}{8}z^5+\frac{705}{16}z^6+\frac{819}{16}z^7.
\]
Following the BD16 examples (where \(\lambda^\ast\) is reported alongside the weight enumerator), we record
\[
\lambda^{\ast 2}=A_1+A_2=\frac{33}{16},\qquad \lambda^\ast=\frac{\sqrt{33}}{4}.
\]



\subsection{A discrete family of exact Full-KL solutions for the same \texorpdfstring{$\vec a$}{a}}

This section records an additional structural observation for the \emph{same} BD16 angle vector
\[
\vec a=(1,1,2,2,2,2,5)\in \mathbb{Z}_8^7,
\]
beyond the single worked example in the previous sections.

\subsubsection{Signature vector and signature norm \texorpdfstring{$\lambda^\ast$}{lambda*}}

Let \(\Pi\) be the projector onto a \(((n,K,d))\) code space. For a Pauli operator \(E\), the KL coefficient
\[
\lambda_E := \frac{1}{K}\mathrm{Tr}(\Pi E)
\]
is characterized by \(\Pi E \Pi = \lambda_E \Pi\) whenever \(E\) is detectable. The \emph{signature vector} on a chosen detectable set \(E_D\) is
\[
\tilde{\lambda} := (\lambda_E)_{E\in E_D},
\]
and the \emph{signature norm} is its Euclidean norm
\[
\lambda^\ast := \|\tilde{\lambda}\|_2 = \left(\sum_{E\in E_D}|\lambda_E|^2\right)^{1/2}.
\]
This quantity is invariant under local unitaries (LU): if two codes are LU-equivalent then they have the same \(\lambda^\ast\).

For distance \(d=3\), a natural choice is
\[
E_D=\{E\in\mathcal{P}_n:\ 1\le \mathrm{wt}(E)\le 2\},
\]
so that \((\lambda^\ast)^2=A_1+A_2\) in terms of the quantum weight enumerator coefficients \(A_j\).

\subsubsection{Restricted exact ansatz}

Motivated by the closed-form examples in the paper, we restrict attention to an exact ansatz with:
\begin{itemize}
\item amplitudes in \(\mathbb{Q}(\sqrt2,\sqrt3,i)\);
\item magnitudes in \(\{1/4,\sqrt2/4,\sqrt3/4\}\) and phases among \(\{\pm 1,\pm i\}\);
\item transversal \(X\) symmetry \(\ket{1_L}=X^{\otimes 7}\ket{0_L}\);
\item subset-sum support constraint \(\mathrm{supp}(\ket{0_L})\subseteq S_0(\vec a)\), hence \(\mathrm{supp}(\ket{1_L})\subseteq S_1(\vec a)\).
\end{itemize}

Within this restricted ansatz, one can enumerate exact solutions and then verify the full KL equalities (all Pauli errors of weight \(\le 2\)) symbolically.

\subsubsection{A 192-element family (explicit parametrization)}

Define the following residue-0 basis strings (all lie in \(S_0(\vec a)\)):

\paragraph{High-weight block.}
\[
H_3=1101110,\quad H_4=1110110,\quad H_5=1111010,\quad H_6=1111100.
\]

\paragraph{One-hot-with-\(s_7=1\) pairs.}
For each \(j\in\{3,4,5,6\}\), define a pair \((P^{(01)}_j,P^{(10)}_j)\) by
\[
\begin{aligned}
j=3:&\quad (P^{(01)}_3,P^{(10)}_3)=(0110001,\ 1010001),\\
j=4:&\quad (P^{(01)}_4,P^{(10)}_4)=(0101001,\ 1001001),\\
j=5:&\quad (P^{(01)}_5,P^{(10)}_5)=(0100101,\ 1000101),\\
j=6:&\quad (P^{(01)}_6,P^{(10)}_6)=(0100011,\ 1000011).
\end{aligned}
\]

We consider logical states of the form
\[
\boxed{
\ket{0_L}=
\frac{1}{4}\ket{0000000}
+\sigma\,\frac{\sqrt3}{4}\ket{0011110}
+\frac{i}{4}\sum_{j\in\{3,4,5,6\}}\eta_j\,\ket{H_j}
+\frac{i\sqrt2}{4}\sum_{j\in\{3,4,5,6\}}\alpha_j\,\ket{P_j},
\qquad
\ket{1_L}=X^{\otimes 7}\ket{0_L},
}
\]
where the discrete parameters are chosen as follows.

\begin{itemize}
\item \(\sigma\in\{\pm 1\}\) is the sign on the \(\sqrt3\) term.
\item \(\eta_j\in\{\pm 1\}\) satisfy the single constraint
\[
\eta_3\eta_4\eta_5\eta_6 = +1.
\]
\item For each \(j\in\{3,4,5,6\}\), choose \(\ket{P_j}\) to be either \(\ket{P^{(01)}_j}\) or \(\ket{P^{(10)}_j}\), with the \emph{balancing rule} that among the four chosen \(\{\ket{P_j}\}\), exactly two have prefix \(01\) and exactly two have prefix \(10\).
\item The signs \(\alpha_j\in\{\pm 1\}\) are constrained by the prefix grouping: within the two chosen prefix-\(01\) positions, the \(\alpha\)'s are opposite; and within the two chosen prefix-\(10\) positions, the \(\alpha\)'s are opposite (equivalently: pick which one in each prefix group gets \(+\), the other gets \(-\)).
\end{itemize}

\paragraph{Counting (why 192).}
There are \(2\) choices for \(\sigma\). There are \(8\) choices of \((\eta_3,\eta_4,\eta_5,\eta_6)\) with product \(+1\).
For a fixed \(\eta\), the admissible choices of which two indices use prefix \(01\) (and which two use prefix \(10\)) depend on the pattern of \(\eta\):
\begin{itemize}
\item If \(\eta\) has two \(+1\) and two \(-1\), then the prefix-\(01\) index set must equal either the \(+1\) set or the \(-1\) set (2 choices); otherwise any 2-subset works (6 choices).
\item For each admissible prefix split, the \(\alpha\)-constraints contribute \(2\times 2=4\) sign choices.
\end{itemize}
This yields \(192\) distinct solutions within the ansatz.

\subsubsection{Observed invariants: all solutions share the same \texorpdfstring{$\lambda^\ast$}{lambda*}}

For every solution in the 192-element family (within this ansatz), symbolic verification confirms the full KL constraints for all Pauli errors of weight \(\le 2\), i.e., each is a valid \(((7,2,3))\) code.

Moreover, the computed signature norm is constant across the family:
\[
\boxed{\lambda^{\ast 2}=\frac{33}{16},\qquad \lambda^\ast=\frac{\sqrt{33}}{4}.}
\]
(Equivalently, \(A_1=0\) and \(A_2=33/16\), so \(\lambda^{\ast 2}=A_1+A_2=A_2\).)

\subsubsection{A concrete additional example}

An explicit member of the family (different \(\sqrt2\)-support and sign choices compared to the earlier worked example) is:
\[
\begin{aligned}
\ket{0_L}
=&\ \frac{1}{4}\ket{0000000}
-\frac{\sqrt3}{4}\ket{0011110}\\
&+\frac{i\sqrt2}{4}\Big(
\ket{0110001}-\ket{0101001}+\ket{1000101}-\ket{1000011}
\Big)\\
&+\frac{i}{4}\Big(
\ket{1101110}+\ket{1110110}+\ket{1111010}+\ket{1111100}
\Big),
\qquad \ket{1_L}=X^{\otimes 7}\ket{0_L}.
\end{aligned}
\]

\subsubsection{Reproducible SymPy script to generate all 192 and compute \texorpdfstring{$\lambda^\ast$}{lambda*}}

The following SymPy code enumerates the 192 solutions within the above ansatz, verifies the full KL conditions for all Paulis of weight \(\le 2\), and computes \((\lambda^\ast)^2\) on the detectable set \(E_D=\{1\le \mathrm{wt}(E)\le 2\}\).
It prints the number of solutions found and the set of distinct \((\lambda^\ast)^2\) values.

\begin{lstlisting}[language=Python]
import sympy as sp
from sympy import Rational, sqrt, I
import itertools
from itertools import combinations

paulis = ['I','X','Y','Z']
K = 2

def wt(P):
    return sum(1 for p in P if p != 'I')

def complement(s):
    return "".join('1' if b=='0' else '0' for b in s)

def apply_single_pauli(p, bit):
    if p=='I': return sp.Integer(1), bit
    if p=='X': return sp.Integer(1), ('1' if bit=='0' else '0')
    if p=='Z': return (sp.Integer(1) if bit=='0' else -sp.Integer(1)), bit
    # Y|0> = i|1>, Y|1> = -i|0>
    return (I if bit=='0' else -I), ('1' if bit=='0' else '0')

def apply_pauli(P, s):
    phase = sp.Integer(1)
    out = list(s)
    for i,p in enumerate(P):
        ph, nb = apply_single_pauli(p, out[i])
        phase *= ph
        out[i] = nb
    return phase, "".join(out)

all_strings = ["".join(bits) for bits in itertools.product('01', repeat=7)]
P_list_w2  = [P for P in itertools.product(paulis, repeat=7) if wt(P) <= 2]
P_list_w12 = [P for P in itertools.product(paulis, repeat=7) if 1 <= wt(P) <= 2]

action = {}
for P in P_list_w2:
    d={}
    for s in all_strings:
        d[s]=apply_pauli(P,s)
    action[P]=d

def inner(cbra, P, cket):
    tot = sp.Integer(0)
    dP = action[P]
    for s, amp in cket.items():
        ph, t = dP[s]
        amp2 = cbra.get(t)
        if amp2 is not None:
            tot += sp.conjugate(cbra[t]) * (ph*amp)
    return sp.simplify(tot)

def kl_ok(c0):
    c1 = {complement(s): coef for s,coef in c0.items()}
    for P in P_list_w2:
        P = "".join(P)
        off = sp.simplify(inner(c0,P,c1))
        d0  = sp.simplify(inner(c0,P,c0))
        d1  = sp.simplify(inner(c1,P,c1))
        if off != 0 or sp.simplify(d0-d1) != 0:
            return False
    return True

def lambda_star_sq(c0):
    c1 = {complement(s): coef for s,coef in c0.items()}
    tot = sp.Integer(0)
    for P in P_list_w12:
        P = "".join(P)
        tr = sp.simplify(inner(c0,P,c0) + inner(c1,P,c1))
        tot += sp.simplify(tr*sp.conjugate(tr))
    return sp.simplify(tot/(K**2))

pair_by_hot = {
    3: ("0110001","1010001"),
    4: ("0101001","1001001"),
    5: ("0100101","1000101"),
    6: ("0100011","1000011"),
}

H_terms = {
    3: "1101110",
    4: "1110110",
    5: "1111010",
    6: "1111100",
}

def make_code(S01, plus01_hot, plus10_hot, eta, sigma):
    c0 = {
        "0000000": Rational(1,4),
        "0011110": sigma*sqrt(3)/4,
    }
    for j in [3,4,5,6]:
        c0[H_terms[j]] = eta[j]*I/4
    for j in [3,4,5,6]:
        s01, s10 = pair_by_hot[j]
        if j in S01:
            s = s01
            sign = +1 if j == plus01_hot else -1
        else:
            s = s10
            sign = +1 if j == plus10_hot else -1
        c0[s] = sign*I*sqrt(2)/4
    return c0

def eta_patterns():
    pats=[]
    for bits in itertools.product([+1,-1], repeat=4):
        eta = {j:bits[k] for k,j in enumerate([3,4,5,6])}
        if eta[3]*eta[4]*eta[5]*eta[6] == 1:
            pats.append(eta)
    return pats

solutions = []
etas = eta_patterns()

for sigma in [+1,-1]:
    for eta in etas:
        plus_set = {j for j in [3,4,5,6] if eta[j]==+1}
        minus_set = set([3,4,5,6]) - plus_set

        if len(plus_set)==2:
            S01_candidates = [plus_set, minus_set]
        else:
            S01_candidates = [set(S01) for S01 in combinations([3,4,5,6], 2)]

        for S01 in S01_candidates:
            comp = set([3,4,5,6]) - set(S01)
            for plus01_hot in S01:
                for plus10_hot in comp:
                    c0 = make_code(set(S01), plus01_hot, plus10_hot, eta, sigma)
                    if kl_ok(c0):
                        lam2 = lambda_star_sq(c0)
                        solutions.append(lam2)

print("number of solutions found:", len(solutions))
print("distinct lambda*^2 values:", sorted(set(solutions)))
if solutions:
    print("lambda* =", sp.sqrt(solutions[0]))
\end{lstlisting}


\section{Additional exact BD16 examples for \texorpdfstring{$\vec a=(0,1,2,2,3,3,4)$}{a=(0,1,2,2,3,3,4)}}

We now record three further \emph{closed-form} \(((7,2,3))\) BD16 examples associated with the SS--LP-passing angle vector
\[
\vec a=(0,1,2,2,3,3,4)\in\mathbb{Z}_8^7.
\]
This vector appears in the candidate list in Sec.~3, hence it passes the SS and LP filters (i.e., the \(Z\)-type KL constraints). The three codes below satisfy the \emph{full} KL conditions for all Paulis of weight \(\le 2\) and therefore are valid distance-3 codes.

\subsection{Conventions}

We use the same conventions as above:
\begin{itemize}
\item Computational basis strings are length-7 bitstrings.
\item Logical basis is fixed by transversal \(X\):
\[
\ket{1_L}=X^{\otimes 7}\ket{0_L}.
\]
\item Distance \(3\) is verified by checking KL for all Pauli strings \(P\in\{I,X,Y,Z\}^{\otimes 7}\) with \(\mathrm{wt}(P)\le 2\):
\[
\langle j_L|P|k_L\rangle=\lambda_P\,\delta_{jk}.
\]
\end{itemize}

Each of the following solutions has coefficients in \(\mathbb{Q}(\sqrt2,\sqrt3,i)\) (``square roots of rationals times 4th roots of unity'') and comes with a different signature norm \(\lambda^\ast\).

\subsection{Solution A: \texorpdfstring{$\lambda^\ast=\sqrt{51}/4$}{lambda*=sqrt(51)/4}}

\[
\begin{aligned}
\ket{0_L}^{(A)} \;=\;
&-\frac{i}{4}\ket{0100011}
-\frac{i}{2}\ket{0100101}
+\frac{\sqrt3}{4}\ket{0111010}\\
&-\frac{i}{4}\ket{1000000}
+\frac{\sqrt2}{4}\ket{1001110}
+\frac{\sqrt2}{4}\ket{1010110}
+\frac{\sqrt3}{4}\ket{1011001},
\qquad
\ket{1_L}^{(A)} = X^{\otimes 7}\ket{0_L}^{(A)}.
\end{aligned}
\]

The signature norm and Shor--Laflamme weight enumerators are:
\[
\boxed{\lambda^{\ast 2}=\frac{51}{16},\qquad \lambda^\ast=\frac{\sqrt{51}}{4}.}
\]
\[
A(z)=1+\frac{51}{16}z^2+\frac{117}{8}z^4+\frac{723}{16}z^6,
\]
\[
B(z)=1+\frac{51}{16}z^2+\frac{489}{16}z^3+\frac{117}{8}z^4+\frac{855}{8}z^5+\frac{723}{16}z^6+\frac{873}{16}z^7.
\]

\subsection{Solution B: \texorpdfstring{$\lambda^\ast=\sqrt{23}/4$}{lambda*=sqrt(23)/4}}

\[
\begin{aligned}
\ket{0_L}^{(B)} \;=\;
&\frac{1}{4}\ket{0000000}
+\frac{i\sqrt2}{4}\ket{0010110}
-\frac{i\sqrt3}{4}\ket{0100011}
-\frac{\sqrt2}{4}\ket{0111100}\\
&+\frac{\sqrt2}{4}\ket{1001110}
-\frac{\sqrt3}{4}\ket{1011001}
+\frac{\sqrt2}{4}\ket{1100101}
+\frac{i}{4}\ket{1111010},
\qquad
\ket{1_L}^{(B)} = X^{\otimes 7}\ket{0_L}^{(B)}.
\end{aligned}
\]

\[
\boxed{\lambda^{\ast 2}=\frac{23}{16},\qquad \lambda^\ast=\frac{\sqrt{23}}{4}.}
\]
\[
A(z)=1+\frac{23}{16}z^2+\frac{145}{8}z^4+\frac{695}{16}z^6,
\]
\[
B(z)=1+\frac{23}{16}z^2+\frac{405}{16}z^3+\frac{145}{8}z^4+\frac{939}{8}z^5+\frac{695}{16}z^6+\frac{789}{16}z^7.
\]

\subsection{Solution C: \texorpdfstring{$\lambda^\ast=\sqrt{30}/4$}{lambda*=sqrt(30)/4}}

\[
\begin{aligned}
\ket{0_L}^{(C)} \;=\;
&-\frac{\sqrt3}{4}\ket{0011001}
-\frac{\sqrt3}{4}\ket{0100101}
-\frac{i\sqrt2}{4}\ket{0111010}
+\frac{1}{4}\ket{1000000}\\
&+\frac{i\sqrt2}{4}\ket{1001110}
-\frac{i\sqrt2}{4}\ket{1010110}
-\frac{i\sqrt2}{4}\ket{1100011}
-\frac{1}{4}\ket{1111100},
\qquad
\ket{1_L}^{(C)} = X^{\otimes 7}\ket{0_L}^{(C)}.
\end{aligned}
\]

\[
\boxed{\lambda^{\ast 2}=\frac{15}{8}=\frac{30}{16},\qquad \lambda^\ast=\frac{\sqrt{30}}{4}.}
\]
\[
A(z)=1+\frac{15}{8}z^2+\frac{69}{4}z^4+\frac{351}{8}z^6,
\]
\[
B(z)=1+\frac{15}{8}z^2+\frac{213}{8}z^3+\frac{69}{4}z^4+\frac{459}{4}z^5+\frac{351}{8}z^6+\frac{405}{8}z^7.
\]

\subsection{Minimal SymPy verification snippet (exact)}

To verify any of the three solutions in the same style as Sec.~5--6, define a dictionary \texttt{c0} with the amplitudes of \(\ket{0_L}\) and then set \(\ket{1_L}\) by bitwise complement (transversal \(X^{\otimes 7}\)). For example, Solution~A can be encoded as:

\begin{lstlisting}[language=Python]
import sympy as sp
from sympy import Rational, sqrt, I

c0 = {
 "0100011": -I*Rational(1,4),
 "0100101": -I*Rational(1,2),
 "0111010":  sqrt(3)/4,
 "1000000": -I*Rational(1,4),
 "1001110":  sqrt(2)/4,
 "1010110":  sqrt(2)/4,
 "1011001":  sqrt(3)/4,
}
# then set c1 = {complement(s): coef for s,coef in c0.items()}
# and run the same weight<=2 Pauli KL check and weight-enumerator/lambda* code as above.
\end{lstlisting}

(Analogous dictionaries for Solutions~B and~C follow immediately from their displayed expansions.)

\section{Exact BD16 examples for \texorpdfstring{$\vec a=(1,1,2,2,3,3,3)$}{a=(1,1,2,2,3,3,3)}}
\label{sec:bd16_1122333}

We now treat the SS--LP-passing BD16 candidate
\[
\vec a=(1,1,2,2,3,3,3)\in\mathbb{Z}_8^7,\qquad (b_0,b_1)=(0,7).
\]
Here \(\sum_{j=1}^7 a_j=15\equiv 7\pmod 8\), so the usual convention \(\ket{1_L}=X^{\otimes 7}\ket{0_L}\) maps residue \(0\) support to residue \(7\) support.

\subsection{An exact one-parameter family}

Let \(\theta\in\mathbb{R}\). Define
\[
\boxed{
\begin{aligned}
\ket{0_L(\theta)} :=
&\ \frac{1}{4}\ket{0000000}
+\frac{\sqrt2}{4}e^{i\theta}\ket{0001011}
-\frac{\sqrt2}{4}e^{i\theta}\ket{0010011}
+i\frac{\sqrt3}{4}e^{i\theta}\ket{0111100}\\[2pt]
&\ +i\frac{\sqrt6}{8}e^{i\theta}\ket{1011001}
-i\frac{\sqrt6}{8}e^{i\theta}\ket{1011010}
-i\frac{\sqrt{10}}{8}e^{i\theta}\ket{1100101}
-i\frac{\sqrt{10}}{8}e^{i\theta}\ket{1100110},
\\[2pt]
\ket{1_L(\theta)} :=&\ X^{\otimes 7}\ket{0_L(\theta)}.
\end{aligned}
}
\]

\subsection{Basic checks}

All checks are exact and hold for every \(\theta\in\mathbb{R}\).

\paragraph{SS residue condition.}
Each basis string \(s\) appearing in \(\ket{0_L(\theta)}\) satisfies
\[
s_1+s_2+2(s_3+s_4)+3(s_5+s_6+s_7)\equiv 0\pmod 8,
\]
hence \(\mathrm{supp}(\ket{0_L(\theta)})\subseteq S_0(\vec a)\).
Since \(\sum_j a_j\equiv 7\pmod 8\), complements lie in residue \(7\), so \(\mathrm{supp}(\ket{1_L(\theta)})\subseteq S_1(\vec a)\).

\paragraph{Normalization.}
Independently of \(\theta\),
\[
\|\ket{0_L(\theta)}\|^2
=\frac{1}{16}+\frac{2}{16}+\frac{2}{16}+\frac{3}{16}+\frac{3}{32}+\frac{3}{32}+\frac{5}{32}+\frac{5}{32}=1.
\]

\paragraph{Full KL (distance 3).}
Symbolic verification (with \(\theta\) declared real in SymPy) confirms that for all Pauli strings
\(P\in\{I,X,Y,Z\}^{\otimes 7}\) of weight \(\mathrm{wt}(P)\le 2\),
\[
\langle j_L(\theta)|P|k_L(\theta)\rangle=\lambda_P\,\delta_{jk},
\]
so this is a valid \(((7,2,3))\) code for every \(\theta\).

\subsection{A \texorpdfstring{$\theta=0$}{theta=0} representative}

Setting \(\theta=0\) yields a particularly simple closed form:
\[
\begin{aligned}
\ket{0_L}
=&\ \frac{1}{4}\ket{0000000}
+\frac{\sqrt2}{4}\ket{0001011}
-\frac{\sqrt2}{4}\ket{0010011}
+i\frac{\sqrt3}{4}\ket{0111100}\\
&\ +i\frac{\sqrt6}{8}\ket{1011001}
-i\frac{\sqrt6}{8}\ket{1011010}
-i\frac{\sqrt{10}}{8}\ket{1100101}
-i\frac{\sqrt{10}}{8}\ket{1100110},
\qquad
\ket{1_L}=X^{\otimes 7}\ket{0_L}.
\end{aligned}
\]

\subsection{Signature norm and weight enumerators}

Using the convention \((\lambda^\ast)^2=A_1+A_2\) with \(E_D=\{1\le \mathrm{wt}(E)\le 2\}\), one finds (independent of \(\theta\)):
\[
\boxed{\lambda^{\ast 2}=\frac{81}{32},\qquad \lambda^\ast=\sqrt{\frac{81}{32}}=\frac{9\sqrt2}{8}.}
\]

The Shor--Laflamme weight enumerators depend only on \(c:=\cos(2\theta)\):
\[
\begin{aligned}
A=&\Big[
1,\ 0,\ \frac{81}{32},\ \frac{1+c}{8},\ \frac{59}{4}-\frac{19}{16}c,\ \frac{45}{16}(1+c),\ \frac{1343}{32}-\frac{41}{16}c,\ \frac{13}{16}(1+c)
\Big],\\[2pt]
B=&\Big[
1,\ 0,\ \frac{81}{32},\ \frac{893}{32}-\frac{11}{16}c,\ \frac{307}{16}+\frac{13}{4}c,\ \frac{1683}{16}-\frac{45}{8}c,\ \frac{1561}{32}+\frac{17}{4}c,\ \frac{1645}{32}-\frac{19}{16}c
\Big].
\end{aligned}
\]

\subsection{Minimal SymPy encoding (symbolic \texorpdfstring{$\theta$}{theta})}

\begin{lstlisting}[language=Python]
import sympy as sp
from sympy import Rational, sqrt, I

theta = sp.Symbol('theta', real=True)

c0 = {
 "0000000": Rational(1,4),
 "0001011": sqrt(2)/4 * sp.exp(I*theta),
 "0010011": -sqrt(2)/4 * sp.exp(I*theta),
 "0111100": I*sqrt(3)/4 * sp.exp(I*theta),
 "1011001": I*sqrt(6)/8 * sp.exp(I*theta),
 "1011010": -I*sqrt(6)/8 * sp.exp(I*theta),
 "1100101": -I*sqrt(10)/8 * sp.exp(I*theta),
 "1100110": -I*sqrt(10)/8 * sp.exp(I*theta),
}
# set c1 by bitwise complement and run the same wt<=2 KL check as in the previous sections.
\end{lstlisting}

\section{Exact BD16 examples and non-uniform families for \texorpdfstring{$\vec a=(1,1,2,2,2,3,4)$}{a=(1,1,2,2,2,3,4)}}
\label{sec:bd16_1122234}

We now treat the SS--LP-passing BD16 candidate
\[
\vec a=(1,1,2,2,2,3,4)\in\mathbb{Z}_8^7,\qquad (b_0,b_1)=(0,7).
\]
Again \(\sum_j a_j=15\equiv 7\pmod 8\), so \(\ket{1_L}=X^{\otimes 7}\ket{0_L}\) is compatible with the SS residue split.

\subsection{SS and LP preliminaries}

For this \(\vec a\), the residue-\(0\) class has size \(16\):
\[
\begin{aligned}
S_0(\vec a)=\{&
0000000,\ 0001101,\ 0010101,\ 0011001,\ 0100011,\ 0101110,\ 0110110,\ 0111010,\\
&1000011,\ 1001110,\ 1010110,\ 1011010,\ 1100101,\ 1101001,\ 1110001,\ 1111100
\}.
\end{aligned}
\]
A particularly simple LP-feasible probability vector is the uniform choice
\[
p_s=\frac{1}{16}\qquad (s\in S_0(\vec a)),
\]
since each qubit is balanced across \(S_0(\vec a)\) (eight \(0\)s and eight \(1\)s), implying \(\langle Z_j\rangle=0\) for all \(j\).

\subsection{Uniform-magnitude exact solutions (phase-only)}

Fix \(\ket{1_L}=X^{\otimes 7}\ket{0_L}\). Within the uniform-magnitude ansatz
\[
c_s=\frac{1}{4}\omega_s,\qquad \omega_s\in\{\pm 1,\pm i\},\qquad s\in S_0(\vec a),
\]
there are multiple exact full-KL solutions; in particular we record three convenient closed forms with distinct \(\lambda^\ast\).

\subsubsection{Solution 1: \texorpdfstring{$\lambda^\ast=\sqrt7/4$}{lambda*=sqrt7/4}}

\[
\begin{aligned}
\ket{0_L}^{(7)}=
&\ \frac{1}{4}\ket{0000000}
-\frac{i}{4}\ket{0001101}
+\frac{i}{4}\ket{0010101}
-\frac{i}{4}\ket{0011001}
-\frac{i}{4}\ket{0100011}
-\frac{1}{4}\ket{0101110}\\
&\ +\frac{1}{4}\ket{0110110}
+\frac{1}{4}\ket{0111010}
-\frac{i}{4}\ket{1000011}
+\frac{1}{4}\ket{1001110}
+\frac{1}{4}\ket{1010110}
-\frac{1}{4}\ket{1011010}\\
&\ -\frac{1}{4}\ket{1100101}
-\frac{1}{4}\ket{1101001}
-\frac{1}{4}\ket{1110001}
-\frac{i}{4}\ket{1111100},\qquad
\ket{1_L}^{(7)}=X^{\otimes 7}\ket{0_L}^{(7)}.
\end{aligned}
\]

\[
\boxed{(\lambda^\ast)^2=\frac{7}{16},\qquad \lambda^\ast=\frac{\sqrt7}{4}.}
\]
\[
A(z)=1+\frac{7}{16}z^2+\frac{161}{8}z^4+\frac{679}{16}z^6,\qquad
B(z)=1+\frac{7}{16}z^2+\frac{357}{16}z^3+\frac{161}{8}z^4+\frac{987}{8}z^5+\frac{679}{16}z^6+\frac{741}{16}z^7.
\]

\subsubsection{Solution 2: \texorpdfstring{$\lambda^\ast=\sqrt{11}/4$}{lambda*=sqrt11/4}}

\[
\begin{aligned}
\ket{0_L}^{(11)}=
&\ \frac{1}{4}\ket{0000000}
-\frac{i}{4}\ket{0001101}
+\frac{i}{4}\ket{0010101}
+\frac{i}{4}\ket{0011001}
-\frac{i}{4}\ket{0100011}
-\frac{1}{4}\ket{0101110}\\
&\ -\frac{1}{4}\ket{0110110}
-\frac{1}{4}\ket{0111010}
+\frac{i}{4}\ket{1000011}
+\frac{1}{4}\ket{1001110}
-\frac{1}{4}\ket{1010110}
-\frac{1}{4}\ket{1011010}\\
&\ +\frac{1}{4}\ket{1100101}
-\frac{1}{4}\ket{1101001}
-\frac{1}{4}\ket{1110001}
+\frac{i}{4}\ket{1111100},\qquad
\ket{1_L}^{(11)}=X^{\otimes 7}\ket{0_L}^{(11)}.
\end{aligned}
\]

\[
\boxed{(\lambda^\ast)^2=\frac{11}{16},\qquad \lambda^\ast=\frac{\sqrt{11}}{4}.}
\]
\[
A(z)=1+\frac{11}{16}z^2+\frac{157}{8}z^4+\frac{683}{16}z^6,\qquad
B(z)=1+\frac{11}{16}z^2+\frac{369}{16}z^3+\frac{157}{8}z^4+\frac{975}{8}z^5+\frac{683}{16}z^6+\frac{753}{16}z^7.
\]

\subsubsection{Solution 3: \texorpdfstring{$\lambda^\ast=\sqrt{19}/4$}{lambda*=sqrt19/4}}

\[
\begin{aligned}
\ket{0_L}^{(19)}=
&\ \frac{1}{4}\ket{0000000}
-\frac{i}{4}\ket{0001101}
+\frac{i}{4}\ket{0010101}
-\frac{i}{4}\ket{0011001}
+\frac{i}{4}\ket{0100011}
+\frac{1}{4}\ket{0101110}\\
&\ +\frac{1}{4}\ket{0110110}
-\frac{1}{4}\ket{0111010}
-\frac{i}{4}\ket{1000011}
-\frac{1}{4}\ket{1001110}
-\frac{1}{4}\ket{1010110}
+\frac{1}{4}\ket{1011010}\\
&\ -\frac{1}{4}\ket{1100101}
-\frac{1}{4}\ket{1101001}
+\frac{1}{4}\ket{1110001}
-\frac{i}{4}\ket{1111100},\qquad
\ket{1_L}^{(19)}=X^{\otimes 7}\ket{0_L}^{(19)}.
\end{aligned}
\]

\[
\boxed{(\lambda^\ast)^2=\frac{19}{16},\qquad \lambda^\ast=\frac{\sqrt{19}}{4}.}
\]
\[
A(z)=1+\frac{19}{16}z^2+\frac{149}{8}z^4+\frac{691}{16}z^6,\qquad
B(z)=1+\frac{19}{16}z^2+\frac{393}{16}z^3+\frac{149}{8}z^4+\frac{951}{8}z^5+\frac{691}{16}z^6+\frac{777}{16}z^7.
\]

\subsection{Non-uniform magnitudes: \texorpdfstring{$\sqrt2$}{sqrt2}-reweighting families}

Going beyond the uniform-magnitude ansatz on all \(16\) strings, we allow non-uniform magnitudes of ``paper-style'' closed form
\[
\ket{0_L}=\frac{1}{4}\sum_{s\in S_0(\vec a)} \omega_s\,\sqrt{w_s}\,\ket{s},
\qquad
\omega_s\in\{\pm 1,\pm i\},
\qquad
w_s\in\{0,1,2\},
\qquad
\ket{1_L}=X^{\otimes 7}\ket{0_L}.
\]
Thus each basis string is either dropped (\(w_s=0\)), kept at magnitude \(1/4\) (\(w_s=1\)), or boosted to magnitude \(\sqrt2/4\) (\(w_s=2\)).

\subsubsection{Linear (LP) constraints on the weights}

The \(Z\)-type KL constraints reduce to the linear balance conditions
\[
\boxed{
\sum_{s\in S_0(\vec a)} w_s = 16,
\qquad
\sum_{s\in S_0(\vec a)} w_s\, s_j = 8 \quad (j=1,\dots,7).
}
\]
There are exactly \(967\) solutions \((w_s)_{s\in S_0}\in\{0,1,2\}^{16}\) to these balance equations.
(Allowing \(w_s\in\{0,1,2,3,4\}\) instead yields \(3597\) LP-feasible patterns.)

A useful corollary for \(w_s\in\{0,1,2\}\) is that the number of dropped strings equals the number of boosted strings:
if \(n_0:=|\{s:w_s=0\}|\) and \(n_2:=|\{s:w_s=2\}|\), then \(n_0=n_2\).

\subsubsection{Compact generator constraints for full KL (fixed phase patterns)}

Let \(\omega^{(7)},\omega^{(11)},\omega^{(19)}\) denote the phase patterns (on \(S_0(\vec a)\)) coming from Solutions 1--3 above, i.e. \(\omega_s^{(\cdot)}=4c_s\in\{\pm1,\pm i\}\) when \(c_s\) is the corresponding uniform coefficient.

For any fixed such \(\omega\), define the two four-pair matchings in \(S_0(\vec a)\):
\[
\begin{aligned}
\text{(\(X_1\)-pairs)}\quad
&(1000011 \leftrightarrow 1111100),\ (1001110 \leftrightarrow 1110001),\ (1010110 \leftrightarrow 1101001),\ (1011010 \leftrightarrow 1100101),\\
\text{(\(X_2\)-pairs)}\quad
&(0100011 \leftrightarrow 1111100),\ (0101110 \leftrightarrow 1110001),\ (0110110 \leftrightarrow 1101001),\ (0111010 \leftrightarrow 1100101).
\end{aligned}
\]

Within the \(\{0,1,2\}\) reweighting ansatz, an efficient way to generate full-KL solutions is:
\begin{quote}
Pick \(\omega\in\{\omega^{(7)},\omega^{(11)},\omega^{(19)}\}\). Solve the linear balance constraints above, together with the two scalar cancellation equations coming from \(\langle 0_L|X_1|1_L\rangle=0\) and \(\langle 0_L|X_2|1_L\rangle=0\), which become signed sums of \(\sqrt{w_s w_t}\) over the \(X_1\)- and \(X_2\)-pairs (signs determined by \(\omega\)).
\end{quote}

Concretely, for \(\omega^{(7)}\) the two equations are
\[
\begin{aligned}
&\sqrt{w_{1000011}w_{1111100}}
-\sqrt{w_{1001110}w_{1110001}}
-\sqrt{w_{1010110}w_{1101001}}
+\sqrt{w_{1011010}w_{1100101}}=0,\\
&\sqrt{w_{0100011}w_{1111100}}
+\sqrt{w_{0101110}w_{1110001}}
-\sqrt{w_{0110110}w_{1101001}}
-\sqrt{w_{0111010}w_{1100101}}=0.
\end{aligned}
\]
For \(\omega^{(11)}\),
\[
\begin{aligned}
&\sqrt{w_{1000011}w_{1111100}}
-\sqrt{w_{1001110}w_{1110001}}
+\sqrt{w_{1010110}w_{1101001}}
-\sqrt{w_{1011010}w_{1100101}}=0,\\
&-\sqrt{w_{0100011}w_{1111100}}
+\sqrt{w_{0101110}w_{1110001}}
+\sqrt{w_{0110110}w_{1101001}}
-\sqrt{w_{0111010}w_{1100101}}=0.
\end{aligned}
\]
For \(\omega^{(19)}\),
\[
\begin{aligned}
&\sqrt{w_{1000011}w_{1111100}}
-\sqrt{w_{1001110}w_{1110001}}
+\sqrt{w_{1010110}w_{1101001}}
-\sqrt{w_{1011010}w_{1100101}}=0,\\
&-\sqrt{w_{0100011}w_{1111100}}
+\sqrt{w_{0101110}w_{1110001}}
-\sqrt{w_{0110110}w_{1101001}}
+\sqrt{w_{0111010}w_{1100101}}=0.
\end{aligned}
\]

\paragraph{Enumerated solutions (within the \(\{0,1,2\}\) ansatz).}
Solving these equations over \(w_s\in\{0,1,2\}\) yields \(25\) solutions for \(\omega^{(7)}\), \(25\) solutions for \(\omega^{(11)}\), and \(35\) solutions for \(\omega^{(19)}\) (each count includes the uniform pattern \(w_s\equiv 1\)).
All non-uniform solutions occur only when
\[
|\{s:w_s=0\}|=|\{s:w_s=2\}|\in\{4,6\}.
\]
Across all solutions found in this scan, \((\lambda^\ast)^2\) takes values in
\[
\left\{
\frac{7}{16},\ \frac{11}{16},\ \frac{17}{16},\ \frac{19}{16},\ \frac{21}{16},\ \frac{25}{16},\ \frac{17}{8},\ \frac{7}{4},\ 2,\ \frac{35}{16},\ \frac{9}{4}
\right\},
\]
so \(\lambda^\ast\) varies widely even at fixed \(\vec a=(1,1,2,2,2,3,4)\).

\subsubsection{Representative non-uniform exact codes}

All examples below satisfy the full KL constraints for all Pauli errors of weight \(\le 2\) exactly.

\paragraph{Example D: \texorpdfstring{$\lambda^\ast=\sqrt{17}/4$}{lambda*=sqrt17/4}.}
\[
\begin{aligned}
\ket{0_L}^{(17)}=
&\ \frac{1}{4}\ket{0000000}
-\frac{i}{4}\ket{0001101}
+\frac{i}{4}\ket{0010101}
-\frac{i}{4}\ket{0011001}
-\frac{\sqrt2}{4}\ket{0101110}
+\frac{\sqrt2}{4}\ket{0111010}\\
&\ -\frac{i\sqrt2}{4}\ket{1000011}
+\frac{\sqrt2}{4}\ket{1010110}
-\frac{1}{4}\ket{1100101}
-\frac{1}{4}\ket{1101001}
-\frac{1}{4}\ket{1110001}
-\frac{i}{4}\ket{1111100},\\
&\ket{1_L}^{(17)}=X^{\otimes 7}\ket{0_L}^{(17)}.
\end{aligned}
\]
\[
\boxed{(\lambda^\ast)^2=\frac{17}{16},\qquad \lambda^\ast=\frac{\sqrt{17}}{4}.}
\]
\[
A(z)=1+\frac{17}{16}z^2+\frac{151}{8}z^4+\frac{689}{16}z^6,\qquad
B(z)=1+\frac{17}{16}z^2+\frac{387}{16}z^3+\frac{151}{8}z^4+\frac{957}{8}z^5+\frac{689}{16}z^6+\frac{771}{16}z^7.
\]

\paragraph{Example E: \texorpdfstring{$\lambda^\ast=\sqrt2$}{lambda*=sqrt2}.}
\[
\begin{aligned}
\ket{0_L}^{(2)}=
&\ \frac{1}{4}\ket{0000000}
-\frac{i}{4}\ket{0001101}
+\frac{i}{4}\ket{0010101}
-\frac{i}{4}\ket{0011001}
-\frac{\sqrt2}{4}\ket{0101110}
+\frac{\sqrt2}{4}\ket{0111010}\\
&\ -\frac{i}{4}\ket{1000011}
+\frac{1}{4}\ket{1001110}
+\frac{1}{4}\ket{1010110}
-\frac{1}{4}\ket{1011010}
-\frac{\sqrt2}{4}\ket{1100101}
-\frac{\sqrt2}{4}\ket{1110001},\qquad\\
&\ket{1_L}^{(2)}=X^{\otimes 7}\ket{0_L}^{(2)}.
\end{aligned}
\]
\[
\boxed{(\lambda^\ast)^2=2,\qquad \lambda^\ast=\sqrt2.}
\]
\[
A(z)=1+2z^2+17z^4+44z^6,\qquad
B(z)=1+2z^2+27z^3+17z^4+114z^5+44z^6+51z^7.
\]

\paragraph{Example F: \texorpdfstring{$\lambda^\ast=\frac{3}{2}$}{lambda*=3/2}.}
\[
\begin{aligned}
\ket{0_L}^{(9/4)}=
&\ \frac{1}{4}\ket{0000000}
-\frac{i}{4}\ket{0001101}
+\frac{i}{4}\ket{0010101}
+\frac{i}{4}\ket{0011001}
-\frac{\sqrt2}{4}\ket{0110110}
-\frac{\sqrt2}{4}\ket{0111010}\\
&\ +\frac{i}{4}\ket{1000011}
+\frac{1}{4}\ket{1001110}
-\frac{1}{4}\ket{1010110}
-\frac{1}{4}\ket{1011010}
+\frac{\sqrt2}{4}\ket{1100101}
-\frac{\sqrt2}{4}\ket{1101001},\qquad\\
&\ket{1_L}^{(9/4)}=X^{\otimes 7}\ket{0_L}^{(9/4)}.
\end{aligned}
\]
\[
\boxed{(\lambda^\ast)^2=\frac{9}{4},\qquad \lambda^\ast=\frac{3}{2}.}
\]
\[
A(z)=1+\frac{9}{4}z^2+\frac{33}{2}z^4+\frac{177}{4}z^6,\qquad
B(z)=1+\frac{9}{4}z^2+\frac{111}{4}z^3+\frac{33}{2}z^4+\frac{225}{2}z^5+\frac{177}{4}z^6+\frac{207}{4}z^7.
\]

\paragraph{Example G: \texorpdfstring{$\lambda^\ast=\sqrt{35}/4$}{lambda*=sqrt35/4}.}
\[
\begin{aligned}
\ket{0_L}^{(35)}=
&\ \frac{1}{4}\ket{0000000}
+\frac{i}{4}\ket{0010101}
-\frac{i\sqrt2}{4}\ket{0011001}
+\frac{i}{4}\ket{0100011}
+\frac{\sqrt2}{4}\ket{0101110}
+\frac{1}{4}\ket{0110110}\\
&\ -\frac{i}{4}\ket{1000011}
-\frac{\sqrt2}{4}\ket{1001110}
-\frac{1}{4}\ket{1010110}
-\frac{1}{4}\ket{1101001}
+\frac{\sqrt2}{4}\ket{1110001}
-\frac{i}{4}\ket{1111100},\qquad\\
&\ket{1_L}^{(35)}=X^{\otimes 7}\ket{0_L}^{(35)}.
\end{aligned}
\]
\[
\boxed{(\lambda^\ast)^2=\frac{35}{16},\qquad \lambda^\ast=\frac{\sqrt{35}}{4}.}
\]
\[
A(z)=1+\frac{35}{16}z^2+\frac{133}{8}z^4+\frac{707}{16}z^6,\qquad
B(z)=1+\frac{35}{16}z^2+\frac{411}{16}z^3+\frac{133}{8}z^4+\frac{927}{8}z^5+\frac{707}{16}z^6+\frac{795}{16}z^7.
\]

\paragraph{Example H: \texorpdfstring{$\lambda^\ast=\sqrt{34}/4$}{lambda*=sqrt34/4}.}
\[
\begin{aligned}
\ket{0_L}^{(17/8)}=
&\ \frac{1}{4}\ket{0000000}
+\frac{i}{4}\ket{0010101}
-\frac{i\sqrt2}{4}\ket{0011001}
-\frac{i}{4}\ket{0100011}
-\frac{\sqrt2}{4}\ket{0101110}
+\frac{1}{4}\ket{0110110}\\
&\ -\frac{i\sqrt2}{4}\ket{1000011}
+\frac{\sqrt2}{4}\ket{1010110}
-\frac{\sqrt2}{4}\ket{1101001}
-\frac{i\sqrt2}{4}\ket{1111100},\qquad\\
&\ket{1_L}^{(17/8)}=X^{\otimes 7}\ket{0_L}^{(17/8)}.
\end{aligned}
\]
\[
\boxed{(\lambda^\ast)^2=\frac{17}{8},\qquad \lambda^\ast=\frac{\sqrt{34}}{4}.}
\]
\[
A(z)=1+\frac{17}{8}z^2+\frac{67}{4}z^4+\frac{353}{8}z^6,\qquad
B(z)=1+\frac{17}{8}z^2+\frac{217}{8}z^3+\frac{67}{4}z^4+\frac{459}{4}z^5+\frac{353}{8}z^6+\frac{405}{8}z^7.
\]

\section{Exact BD16 examples for \texorpdfstring{$\vec a=(1,1,1,2,3,3,4)$}{a=(1,1,1,2,3,3,4)}}

We give an explicit self-complementary BD16 code for
\(\vec a=(1,1,1,2,3,3,4)\) at level \(m=8\), with residues \((b_0,b_1)=(0,7)\).
Equivalently, the logical basis is related by a global bit-flip:
\[
\ket{1_L} = X^{\otimes 7}\ket{0_L}.
\]

\subsection{Residue constraint}
For a bitstring \(\mathbf{s}=(s_1,\dots,s_7)\in\{0,1\}^7\), define the residue function
\[
\mathrm{Res}(\mathbf{s}) \equiv s_1+s_2+s_3+2s_4+3s_5+3s_6+4s_7 \pmod 8.
\]
All computational-basis kets appearing below satisfy \(\mathrm{Res}(\mathbf{s})=0\).

\subsection{One clean analytical solution}
A normalized choice of logical \(\ket{0_L}\) is
\begin{align*}
\ket{0_L} \;=\;&
\frac{1}{4}\ket{0000000}
+\frac{i\sqrt2}{4}\ket{0001110}
+\frac{\sqrt2}{4}\ket{0110110}
+\frac{\sqrt3}{4}\ket{0111001}\\
&+\frac{i\sqrt{10}}{8}\ket{1000011}
+\frac{i\sqrt{10}}{8}\ket{1000101}
+\frac{i\sqrt6}{8}\ket{1111010}
-\frac{i\sqrt6}{8}\ket{1111100}.
\end{align*}
By self-complementarity,
\[
\ket{1_L} = X^{\otimes 7}\ket{0_L},
\]
so explicitly (bitwise complement of every basis string),
\begin{align*}
\ket{1_L} \;=\;&
\frac{1}{4}\ket{1111111}
+\frac{i\sqrt2}{4}\ket{1110001}
+\frac{\sqrt2}{4}\ket{1001001}
+\frac{\sqrt3}{4}\ket{1000110}\\
&+\frac{i\sqrt{10}}{8}\ket{0111100}
+\frac{i\sqrt{10}}{8}\ket{0111010}
+\frac{i\sqrt6}{8}\ket{0000101}
-\frac{i\sqrt6}{8}\ket{0000011}.
\end{align*}
This pair satisfies the full Knill--Laflamme conditions for all Pauli errors of weight \(\le 2\),
i.e.\ it is a distance-\(3\) \(((7,2,3))\) code.

\subsection{Invariants for this solution}
For this solution,
\[
\lambda_*^2 = A_1 + A_2 = \frac{75}{32},
\qquad\text{hence}\qquad
\lambda_* = \frac{5\sqrt6}{8}.
\]
The weight enumerators are
\begin{align*}
A(z) &= 1+\frac{75}{32}z^2+\frac{261}{16}z^4+\frac{1419}{32}z^6,\\
B(z) &= 1+\frac{75}{32}z^2+\frac{897}{32}z^3+\frac{261}{16}z^4+\frac{1791}{16}z^5+\frac{1419}{32}z^6+\frac{1665}{32}z^7.
\end{align*}

\subsection{Phase variants and symmetries}
Keeping the same support and magnitudes
\(\{1/4,\ \sqrt2/4,\ \sqrt3/4,\ \sqrt6/8,\ \sqrt{10}/8\}\),
there exist multiple other exact solutions obtained by changing certain \(\pm\) and \(\pm i\) phases.
Empirically, there are \(64\) distinct phase variants with these fixed magnitudes on the same support,
after fixing the \(\ket{0000000}\) coefficient to \(+1/4\).
Moreover, since \(\vec a\) is symmetric under permutations of qubits \((1,2,3)\) and independently \((5,6)\),
permuting those qubits produces further solutions.

\section{Three uniform-magnitude BD16 solutions for \texorpdfstring{$\vec a=(1,1,1,2,2,4,4)$}{a=(1,1,1,2,2,4,4)} at \texorpdfstring{$m=8$}{m=8}}

For BD16 we fix \(m=8\) and \((b_0,b_1)=(0,7)\).
Take
\[
\vec a=(1,1,1,2,2,4,4)\in \mathbb{Z}_8^{\,7},
\qquad
\sum_{j=1}^7 a_j = 15 \equiv 7 \pmod 8,
\]
so the usual self-complementary choice \(\ket{1_L}=X^{\otimes 7}\ket{0_L}\) is compatible with the residue split
\(S_0(\vec a)\leftrightarrow S_1(\vec a)\).
Below we record three closed-form exact \(((7,2,3))\) solutions for this \(\vec a\), each with coefficients in \(\mathbb{Q}(i)\),
and each verified \emph{exactly} to satisfy the full Knill--Laflamme conditions for every Pauli error of weight \(\le 2\).

\subsection{Subset-sum and LP facts}
Define the subset-sum residue map
\[
\mathrm{res}(s)=s_1+s_2+s_3+2(s_4+s_5)+4(s_6+s_7)\pmod 8,
\qquad s\in\{0,1\}^7.
\]
The residue-\(0\) class has size \(16\) and is
\[
\begin{aligned}
S_0(\vec a)=\{&
0000000,\ 0000011,\ 0001101,\ 0001110,\ 
0110101,\ 0110110,\ 0111001,\ 0111010,\\
&1010101,\ 1010110,\ 1011001,\ 1011010,\ 
1100101,\ 1100110,\ 1101001,\ 1101010
\}.
\end{aligned}
\]
Moreover, each qubit position is perfectly balanced over \(S_0(\vec a)\) (eight zeros and eight ones in every coordinate).
Consequently the uniform LP choice
\[
p_s=\frac{1}{16}\qquad (s\in S_0)
\]
automatically satisfies the \(Z\)-type Knill--Laflamme constraints (in particular \(\langle Z_j\rangle_{0_L}=0\) for all \(j\)).

Motivated by this, we seek uniform-magnitude solutions of the form
\[
\ket{0_L}=\frac{1}{4}\sum_{s\in S_0(\vec a)}\omega_s\ket{s},
\qquad \omega_s\in\{\pm 1,\pm i\},
\qquad \ket{1_L}=X^{\otimes 7}\ket{0_L}.
\]
All three solutions below are of this type.

\subsection{Solution A (\texorpdfstring{$\lambda^\ast=1$}{lambda*=1})}
\[
\begin{aligned}
\ket{0_L}^{(A)}=\frac{1}{4}\Big(
&\ket{0000000}-\ket{0000011}+\ket{0001101}+\ket{0001110}\\
&+i\ket{0110101}-i\ket{0110110}-i\ket{0111001}+i\ket{0111010}\\
&+i\ket{1010101}+i\ket{1010110}-i\ket{1011001}-i\ket{1011010}\\
&+i\ket{1100101}-i\ket{1100110}+i\ket{1101001}-i\ket{1101010}
\Big),
\qquad
\ket{1_L}^{(A)}=X^{\otimes 7}\ket{0_L}^{(A)}.
\end{aligned}
\]
Here \(A_1=0\), so \((\lambda^\ast)^2=A_2\), and
\[
(\lambda^\ast)^2=1,\qquad \lambda^\ast=1.
\]
The Shor--Laflamme enumerator coefficient lists are
\[
A^{(A)}=\Big[1,\ 0,\ 1,\ \tfrac14,\ 15,\ \tfrac{21}{2},\ 33,\ \tfrac{13}{4}\Big],
\qquad
B^{(A)}=\Big[1,\ 0,\ 1,\ 21,\ 32,\ 99,\ 58,\ 44\Big].
\]

\subsection{Solution B (\texorpdfstring{$\lambda^\ast=\sqrt{5}/2$}{lambda*=sqrt(5)/2})}
\[
\begin{aligned}
\ket{0_L}^{(B)}=\frac{1}{4}\Big(
&\ket{0000000}+\ket{0000011}+\ket{0001101}+\ket{0001110}\\
&+i\ket{0110101}+i\ket{0110110}-i\ket{0111001}-i\ket{0111010}\\
&-i\ket{1010101}-i\ket{1010110}-i\ket{1011001}-i\ket{1011010}\\
&-i\ket{1100101}+i\ket{1100110}-i\ket{1101001}+i\ket{1101010}
\Big),
\qquad
\ket{1_L}^{(B)}=X^{\otimes 7}\ket{0_L}^{(B)}.
\end{aligned}
\]
The invariant is
\[
(\lambda^\ast)^2=\frac{5}{4},\qquad \lambda^\ast=\frac{\sqrt5}{2}.
\]
Enumerator coefficient lists:
\[
A^{(B)}=\Big[1,\ 0,\ \tfrac54,\ \tfrac14,\ \tfrac{29}{2},\ \tfrac{21}{2},\ \tfrac{133}{4},\ \tfrac{13}{4}\Big],
\qquad
B^{(B)}=\Big[1,\ 0,\ \tfrac54,\ \tfrac{87}{4},\ \tfrac{63}{2},\ \tfrac{195}{2},\ \tfrac{233}{4},\ \tfrac{179}{4}\Big].
\]

\subsection{Solution C (\texorpdfstring{$\lambda^\ast=\sqrt{22}/4$}{lambda*=sqrt(22)/4})}
\[
\begin{aligned}
\ket{0_L}^{(C)}=\frac{1}{4}\Big(
&\ket{0000000}-\ket{0000011}-\ket{0001101}+\ket{0001110}\\
&-i\ket{0110101}-i\ket{0110110}-i\ket{0111001}+i\ket{0111010}\\
&-i\ket{1010101}-i\ket{1010110}+i\ket{1011001}-i\ket{1011010}\\
&-i\ket{1100101}-i\ket{1100110}+i\ket{1101001}-i\ket{1101010}
\Big),
\qquad
\ket{1_L}^{(C)}=X^{\otimes 7}\ket{0_L}^{(C)}.
\end{aligned}
\]
The invariant is
\[
(\lambda^\ast)^2=\frac{11}{8},\qquad \lambda^\ast=\sqrt{\frac{11}{8}}=\frac{\sqrt{22}}{4}.
\]
Enumerator coefficient lists:
\[
A^{(C)}=\Big[1,\ 0,\ \tfrac{11}{8},\ \tfrac14,\ \tfrac{57}{4},\ \tfrac{21}{2},\ \tfrac{267}{8},\ \tfrac{13}{4}\Big],
\qquad
B^{(C)}=\Big[1,\ 0,\ \tfrac{11}{8},\ \tfrac{177}{8},\ \tfrac{125}{4},\ \tfrac{387}{4},\ \tfrac{467}{8},\ \tfrac{361}{8}\Big].
\]

\subsection{Minimal SymPy dictionary (drop-in)}
As a concrete example, Solution~B can be represented (for exact verification) by the coefficient dictionary
\begin{lstlisting}[language=Python]
import sympy as sp
from sympy import Rational, I

c0 = {
 "0000000":  Rational(1,4),
 "0000011":  Rational(1,4),
 "0001101":  Rational(1,4),
 "0001110":  Rational(1,4),
 "0110101":  I*Rational(1,4),
 "0110110":  I*Rational(1,4),
 "0111001": -I*Rational(1,4),
 "0111010": -I*Rational(1,4),
 "1010101": -I*Rational(1,4),
 "1010110": -I*Rational(1,4),
 "1011001": -I*Rational(1,4),
 "1011010": -I*Rational(1,4),
 "1100101": -I*Rational(1,4),
 "1100110":  I*Rational(1,4),
 "1101001": -I*Rational(1,4),
 "1101010":  I*Rational(1,4),
}
# c1 = {bitwise_complement(s): coef for s,coef in c0.items()}
# then run your wt<=2 KL check and compute A,B, lambda*^2=A1+A2.
\end{lstlisting}

\section{An explicit self-complementary BD16 code for \texorpdfstring{$\vec a=(1,1,1,2,2,3,5)$}{a=(1,1,1,2,2,3,5)} at \texorpdfstring{$m=8$}{m=8}}

We fix BD16 level \(m=8\) with residues \((b_0,b_1)=(0,7)\).
Consider the angle vector
\[
\vec a=(1,1,1,2,2,3,5)\in\mathbb{Z}_8^{\,7},
\qquad
\sum_{i=1}^7 a_i = 15 \equiv 7 \pmod 8,
\]
so the usual self-complementary convention \(\ket{1_L}=X^{\otimes 7}\ket{0_L}\) is compatible with the residue split
\(S_0(\vec a)\leftrightarrow S_7(\vec a)=\overline{S_0(\vec a)}\).
We take \(\ket{0_L}\) supported on the residue-\(0\) set
\[
S_0(\vec a)=\{\,s\in\{0,1\}^7:\ \vec a\cdot s \equiv 0 \pmod 8\,\}.
\]

\subsection{One closed-form \texorpdfstring{\(((7,2,3))\)}{((7,2,3))} solution}
An exact (fully Knill--Laflamme valid for all Pauli errors of weight \(\le 2\)) choice of codewords is
\begin{align*}
\ket{0_L}
=&\;\frac14\ket{0000000}+\frac14\ket{0000011}\\
&-\frac{i\sqrt3}{8}\Big(
\ket{0010101}+\ket{0011001}+\ket{0100101}+\ket{0101001}+\ket{1000101}+\ket{1001001}
\Big)\\
&-\frac{\sqrt6}{8}\Big(
\ket{0011110}+\ket{0101110}+\ket{1001110}
\Big)\\
&-\frac{\sqrt{10}}{8}\ket{1110001}
+\frac{i\sqrt5}{8}\Big(\ket{1110110}+\ket{1111010}\Big),
\end{align*}
and
\[
\ket{1_L}=X^{\otimes 7}\ket{0_L},
\]
i.e.\ explicitly (bitwise complements of the above basis states, with identical coefficients),
\begin{align*}
\ket{1_L}
=&\;\frac14\ket{1111111}+\frac14\ket{1111100}\\
&-\frac{i\sqrt3}{8}\Big(
\ket{1101010}+\ket{1100110}+\ket{1011010}+\ket{1010110}+\ket{0111010}+\ket{0110110}
\Big)\\
&-\frac{\sqrt6}{8}\Big(
\ket{1100001}+\ket{1010001}+\ket{0110001}
\Big)\\
&-\frac{\sqrt{10}}{8}\ket{0001110}
+\frac{i\sqrt5}{8}\Big(\ket{0001001}+\ket{0000101}\Big).
\end{align*}
This pair satisfies all Knill--Laflamme constraints for distance \(3\), and was verified exactly using symbolic arithmetic.

\subsection{Enumerator data}
For this solution,
\[
\lambda_*^2=A_1+A_2=\frac{369}{128},
\qquad
\lambda_*=\sqrt{\frac{369}{128}}=\frac{3\sqrt{82}}{16}.
\]
The Shor--Laflamme enumerator coefficient lists \([A_0,\dots,A_7]\) and \([B_0,\dots,B_7]\) are
\[
A=\Big[1,\ 0,\ \frac{369}{128},\ 0,\ \frac{975}{64},\ 0,\ \frac{5745}{128},\ 0\Big],
\qquad
B=\Big[1,\ 0,\ \frac{369}{128},\ \frac{3795}{128},\ \frac{975}{64},\ \frac{6957}{64},\ \frac{5745}{128},\ \frac{6867}{128}\Big].
\]
Equivalently,
\begin{align*}
A(z)&=1+\frac{369}{128}z^2+\frac{975}{64}z^4+\frac{5745}{128}z^6,\\
B(z)&=1+\frac{369}{128}z^2+\frac{3795}{128}z^3+\frac{975}{64}z^4+\frac{6957}{64}z^5+\frac{5745}{128}z^6+\frac{6867}{128}z^7.
\end{align*}

\section{A 128-element analytic BD16 family for \texorpdfstring{$\vec a=(1,1,1,1,3,4,4)$}{a=(1,1,1,1,3,4,4)} at \texorpdfstring{$m=8$}{m=8}}

We work in the BD16 setting with \(m=8\) and residues \((b_0,b_1)=(0,7)\), using the standard self-complement convention
\[
\ket{1_L}=X^{\otimes 7}\ket{0_L}.
\]
For the candidate
\[
\vec a=(1,1,1,1,3,4,4)\in\mathbb{Z}_8^{\,7},
\]
there is a fully analytic \(128\)-element exact \(((7,2,3))\) family with coefficients in \(\mathbb{Q}(\sqrt3,i)\),
with magnitudes in \(\{1/4,\sqrt3/4\}\) and phases among the \(4\)th roots of unity.

\subsection{Subset-sum residue class}
The residue constraint is
\[
s_1+s_2+s_3+s_4+3s_5+4s_6+4s_7\equiv 0\pmod 8.
\]
The residue-\(0\) class \(S_0(\vec a)\) has size \(12\) and is
\[
\begin{aligned}
S_0(\vec a)=\{&
0000000,\ 0000011,\ 0001101,\ 0001110,\ 0010101,\ 0010110,\\
&0100101,\ 0100110,\ 1000101,\ 1000110,\ 1111001,\ 1111010
\}.
\end{aligned}
\]
Moreover \(\sum_{j=1}^7 a_j=15\equiv 7\pmod8\), so bitwise complements of \(S_0(\vec a)\) lie in residue class \(7\).
Hence the self-complement map \(\ket{1_L}=X^{\otimes 7}\ket{0_L}\) matches the SS split \(S_0\leftrightarrow S_7\).

\subsection{A clean 128-element exact family}
Introduce the following paired basis strings in \(S_0(\vec a)\) (each pair swaps \((s_6,s_7)\)):
\[
\begin{aligned}
(P_1^{(7)},P_1^{(6)})&=(1000101,\ 1000110),&
(P_2^{(7)},P_2^{(6)})&=(0100101,\ 0100110),\\
(P_3^{(7)},P_3^{(6)})&=(0010101,\ 0010110),&
(P_4^{(7)},P_4^{(6)})&=(0001101,\ 0001110),\\
(D^{(7)},D^{(6)})&=(1111001,\ 1111010).
\end{aligned}
\]
Choose discrete parameters
\[
\alpha_B,\alpha_1,\alpha_2,\alpha_3,\alpha_4,\alpha_D\in\{\pm 1\},\qquad \tau\in\{0,1\}.
\]
Define
\[
\boxed{
\begin{aligned}
\ket{0_L(\alpha,\tau)} \;:=\;
&\ \frac14\ket{0000000}+\frac{\alpha_B}{4}\ket{0000011}\\[2pt]
&+\frac{i}{4}\sum_{j=1}^4 \alpha_j\Big(\ket{P_j^{(7)}} + (-1)^{\tau+1}\ket{P_j^{(6)}}\Big)\\[2pt]
&+\frac{i\sqrt3}{4}\,\alpha_D\Big(\ket{D^{(7)}} + (-1)^{\tau}\ket{D^{(6)}}\Big),\\[4pt]
\ket{1_L(\alpha,\tau)} \;:=\;&\ X^{\otimes 7}\ket{0_L(\alpha,\tau)}.
\end{aligned}
}
\]

\paragraph{Verified properties.}
For every choice of \((\alpha_B,\alpha_1,\ldots,\alpha_4,\alpha_D,\tau)\), the pair
\(\{\ket{0_L(\alpha,\tau)},\ket{1_L(\alpha,\tau)}\}\) satisfies:
\begin{itemize}
\item \textbf{SS support:} \(\mathrm{supp}(\ket{0_L})\subseteq S_0(\vec a)\) and \(\mathrm{supp}(\ket{1_L})\subseteq S_7(\vec a)\).
\item \textbf{Normalization:} there are \(10\) terms of magnitude \(1/4\) and \(2\) terms of magnitude \(\sqrt3/4\), hence
\[
10\cdot\frac{1}{16}+2\cdot\frac{3}{16}=1.
\]
\item \textbf{Full KL / distance \(3\):} for every Pauli error \(P\) of weight \(\le 2\),
\[
\langle j_L|P|k_L\rangle=\lambda_P\,\delta_{jk}.
\]
This was verified symbolically with exact arithmetic.
\end{itemize}
Thus the construction yields an exact \(128\)-element family in closed form.

\subsection{A paste-ready representative}
Setting all signs to \(+1\) and \(\tau=0\) gives the explicit member
\[
\boxed{
\begin{aligned}
\ket{0_L}=
&\ \frac{1}{4}\ket{0000000}
+\frac{1}{4}\ket{0000011}
+\frac{i}{4}\ket{0001101}
-\frac{i}{4}\ket{0001110}
+\frac{i}{4}\ket{0010101}
-\frac{i}{4}\ket{0010110}\\
&+\frac{i}{4}\ket{0100101}
-\frac{i}{4}\ket{0100110}
+\frac{i}{4}\ket{1000101}
-\frac{i}{4}\ket{1000110}
+\frac{i\sqrt3}{4}\ket{1111001}
+\frac{i\sqrt3}{4}\ket{1111010},\\[3pt]
\ket{1_L}=&\ X^{\otimes 7}\ket{0_L}.
\end{aligned}
}
\]
Here all phases lie in \(\{\pm 1,\pm i\}\), and all magnitudes lie in \(\{1/4,\sqrt3/4\}\).

\subsection{Family invariants}
Using the convention \((\lambda^\ast)^2=A_1+A_2\), the value of \(\lambda^\ast\) is constant across the entire family:
\[
(\lambda^\ast)^2=\frac{31}{8}=\frac{62}{16},
\qquad
\lambda^\ast=\sqrt{\frac{31}{8}}=\frac{\sqrt{62}}{4}.
\]
Moreover, the Shor--Laflamme enumerator coefficients are family-invariant:
\[
A=\Big[1,\ 0,\ \frac{31}{8},\ 0,\ \frac{53}{4},\ 0,\ \frac{367}{8},\ 0\Big],
\qquad
B=\Big[1,\ 0,\ \frac{31}{8},\ \frac{261}{8},\ \frac{53}{4},\ \frac{411}{4},\ \frac{367}{8},\ \frac{453}{8}\Big].
\]
Equivalently,
\begin{align*}
A(z)&=1+\frac{31}{8}z^2+\frac{53}{4}z^4+\frac{367}{8}z^6,\\
B(z)&=1+\frac{31}{8}z^2+\frac{261}{8}z^3+\frac{53}{4}z^4+\frac{411}{4}z^5+\frac{367}{8}z^6+\frac{453}{8}z^7.
\end{align*}

\subsection{Remarks on searches beyond the family}
Within the restricted ``paper-style'' ansatz
\[
c_s=\omega_s\frac{\sqrt{w_s}}{4},
\qquad
\omega_s\in\{\pm1,\pm i\},
\qquad
w_s\in\{0,1,2,3,4\},
\]
together with the LP balance constraints
\[
\sum_{s\in S_0} w_s=16,\qquad \sum_{s\in S_0} w_s\, s_j = 8\quad (j=1,\dots,7),
\]
an exhaustive enumeration over all integer-feasible weight patterns and \(\{\pm1,\pm i\}\) phases produced no additional full-KL solutions beyond the above family. This does not preclude further solutions outside this phase regime (e.g.\ allowing \(8\)th-root phases), but it provides a strong classification within the \(\{\pm1,\pm i\}\) setting.

\section{A continuous analytic BD16 family for \texorpdfstring{$\vec a=(1,1,1,1,3,3,5)$}{a=(1,1,1,1,3,3,5)} at \texorpdfstring{$m=8$}{m=8}}

We work in the BD16 setting with \(m=8\) and residues \((b_0,b_1)=(0,7)\), using the standard self-complement convention
\[
\ket{1_L}=X^{\otimes 7}\ket{0_L}.
\]
Consider the angle vector
\[
\vec a=(1,1,1,1,3,3,5)\in\mathbb{Z}_8^{\,7}.
\]

\subsection{Residue class \(S_0(\vec a)\)}
With \(b_0=0\), the subset-sum residue class is
\[
S_0(\vec a)=\{\,s\in\{0,1\}^7:\ \vec a\cdot s \equiv 0 \pmod 8\,\},
\]
and explicitly
\[
\begin{aligned}
S_0(\vec a)=\{&
0000000,\ 0000011,\ 0000101,\ 0011110,\ 0101110,\ 0110110,\ 0111001,\\
&1001110,\ 1010110,\ 1011001,\ 1100110,\ 1101001,\ 1110001
\}.
\end{aligned}
\]
We take \(\ket{0_L}\) supported on \(S_0(\vec a)\), and \(\ket{1_L}\) supported on the complement class \(S_7(\vec a)=\overline{S_0(\vec a)}\).

\subsection{A clean closed-form representative}
For any real \(\theta\), define
\begin{align*}
\ket{0_L(\theta)}
=&\;\frac14\Bigl(\ket{0000000}-\ket{0000011}-\ket{0000101}\Bigr)\\[2mm]
&+e^{i\theta}\Biggl[
\frac14\Bigl(\ket{0011110}+\ket{1100110}\Bigr)
+\frac18\Bigl(\ket{0101110}+\ket{1010110}\Bigr)\\
&
+\frac38\Bigl(\ket{0110110}+\ket{1001110}\Bigr)
+\frac{\sqrt6}{8}\Bigl(\ket{1011001}+\ket{1101001}-\ket{0111001}-\ket{1110001}\Bigr)
\Biggr].
\end{align*}
Then set
\[
\ket{1_L(\theta)}:=X^{\otimes 7}\ket{0_L(\theta)}.
\]
The pair \(\{\ket{0_L(\theta)},\ket{1_L(\theta)}\}\) satisfies the Knill--Laflamme conditions for all Pauli errors of weight \(\le 2\),
hence it is an exact \(((7,2,3))\) code in this BD16 single-syndrome residue-class setting.  Taking \(\theta=0\) yields an
all-real representative.

\subsection{The full continuous family}
More generally, let \(x>0\) and \(y>0\) satisfy
\[
x^2+y^2+xy=7,
\qquad\text{and define}\qquad
z:=x+y.
\]
(Equivalently \(x^2+y^2+z^2=14\) with \(z=x+y\).)
Then for any real \(\theta\), an exact solution is given by
\begin{align*}
\ket{0_L(\theta;x,y)}
=&\;\frac14\Bigl(\ket{0000000}-\ket{0000011}-\ket{0000101}\Bigr)\\[2mm]
&+e^{i\theta}\Biggl[
\frac{y}{8}\Bigl(\ket{0011110}+\ket{1100110}\Bigr)
+\frac{x}{8}\Bigl(\ket{0101110}+\ket{1010110}\Bigr)\\
&
+\frac{z}{8}\Bigl(\ket{0110110}+\ket{1001110}\Bigr)
+\frac{\sqrt6}{8}\Bigl(\ket{1011001}+\ket{1101001}-\ket{0111001}-\ket{1110001}\Bigr)
\Biggr],
\end{align*}
with
\[
\ket{1_L(\theta;x,y)}:=X^{\otimes 7}\ket{0_L(\theta;x,y)}.
\]
This yields a continuous family of exact \(((7,2,3))\) codes parameterized by \((\theta;x,y)\) subject to \(x^2+y^2+xy=7\).

A convenient one-parameter parametrization is obtained by setting \(y=r x\) with \(r>0\), giving
\[
x=\sqrt{\frac{7}{1+r+r^2}},\qquad y=r x,\qquad z=(1+r)x,
\]
which can be substituted into the above formula. The ``clean'' representative above corresponds to \(r=2\), i.e.\
\((x,y,z)=(1,2,3)\).

\paragraph{Weight enumerators and $\lambda^\ast$ for the $(x,y,z,\theta)$ family.}
Impose $z=x+y$ and $x^2+y^2+xy=7$.  Let $c:=\cos(2\theta)$.  Then the Shor--Laflamme
(Rains-type) enumerators are
\[
A=\Big[
1,\ 0,\ \frac{177}{64},\ 0,\ \frac{3}{32}\bigl(139+26c\bigr),\ \frac{117}{16}\bigl(1-c\bigr),\ 
\frac{3}{64}\bigl(799+156c\bigr),\ \frac{39}{16}\bigl(1-c\bigr)
\Big],
\]
\[
B=\Big[
1,\ 0,\ \frac{177}{64},\ \frac{3}{64}\bigl(573+52c\bigr),\ \frac{3}{32}\bigl(269-104c\bigr),\ 
\frac{9}{32}\bigl(337+52c\bigr),\ \frac{3}{64}\bigl(1163-208c\bigr),\ \frac{3}{64}\bigl(1085+52c\bigr)
\Big].
\]
In particular $A_1=0$ and $A_2=177/64$, so
\[
(\lambda^\ast)^2=A_1+A_2=\frac{177}{64},
\qquad
\lambda^\ast=\frac{\sqrt{177}}{8},
\]
independent of $(x,y,z,\theta)$ within the family.



\section{Checklist summary}

This note focuses on the BD16 (transversal \(T\)) specialization at \(m=8\), with residues \((b_0,b_1)=(0,7)\), and the SS--LP candidate angle vectors listed in Sec.~3 (BD16 candidate list).

\subsection{Candidate vectors treated in this note (Full-KL verified)}

Among the SS--LP-passing BD16 candidates, the following \(\vec a\) have been taken through the Full-KL stage in this note (either via an explicit construction from the cited paper, or via new exact solutions verified symbolically):

\begin{itemize}
\item \(\vec a=(1,2,2,2,2,3,3)\): the explicit two-parameter family from Sec.~6.3.1 of the cited paper (free phases \(\theta_1,\theta_2\)), with \(\lambda^\ast=\sqrt{21/8}\) and weight enumerators depending on \(a=\cos(2\theta_1)+\cos(2\theta_2)\), \(b=\cos(2\theta_1-2\theta_2)\).
\item \(\vec a=(1,1,2,2,2,2,5)\): the detailed exact solution and the 192-element discrete family;\\ all have coefficients in \(\mathbb{Q}(\sqrt2,\sqrt3,i)\) and share \(\lambda^\ast=\sqrt{33}/4\).
\item \(\vec a=(0,1,2,2,3,3,4)\): three closed-form exact solutions (Solutions A/B/C) \\
with \(\lambda^\ast\in\{\sqrt{51}/4,\ \sqrt{23}/4,\ \sqrt{30}/4\}\).
\item \(\vec a=(1,1,2,2,3,3,3)\): an exact one-parameter family \(\{\ket{0_L(\theta)},\ket{1_L(\theta)}\}\) with \(\lambda^\ast=\frac{9\sqrt2}{8}\) and weight enumerators depending on \(c=\cos(2\theta)\).
\item \(\vec a=(1,1,2,2,2,3,4)\): three uniform-magnitude (phase-only) exact solutions with \((\lambda^\ast)^2\in\{7/16,11/16,19/16\}\), together with additional non-uniform \(\sqrt2\)-reweighting families whose \((\lambda^\ast)^2\) values range over
\[
\left\{
\frac{7}{16},\ \frac{11}{16},\ \frac{17}{16},\ \frac{19}{16},\ \frac{21}{16},\ \frac{25}{16},\ \frac{17}{8},\ \frac{7}{4},\ 2,\ \frac{35}{16},\ \frac{9}{4}
\right\}.
\]
\item \(\vec a=(1,1,1,2,3,3,4)\): an explicit self-complementary \(((7,2,3))\) BD16 code at level \(m=8\) with residues \((b_0,b_1)=(0,7)\) (i.e.\ \(\ket{1_L}=X^{\otimes 7}\ket{0_L}\)), admitting a closed-form solution supported on \(8\) basis states; it has \(\lambda^\ast=\frac{5\sqrt6}{8}\) (so \((\lambda^\ast)^2=\frac{75}{32}\)) and weight enumerators
\(A(z)=1+\frac{75}{32}z^2+\frac{261}{16}z^4+\frac{1419}{32}z^6\),
\(B(z)=1+\frac{75}{32}z^2+\frac{897}{32}z^3+\frac{261}{16}z^4+\frac{1791}{16}z^5+\frac{1419}{32}z^6+\frac{1665}{32}z^7\);
moreover, fixing the same support and magnitudes yields \(64\) distinct phase variants, with further solutions generated by permuting qubits \((1,2,3)\) and \((5,6)\).
\item \(\vec a=(1,1,1,2,2,4,4)\): three uniform-magnitude (phase-only) exact \(((7,2,3))\) solutions at BD16 level \(m=8\) (with \((b_0,b_1)=(0,7)\) and \(\ket{1_L}=X^{\otimes 7}\ket{0_L}\)), with
\((\lambda^\ast)^2\in\left\{1,\ \tfrac{5}{4},\ \tfrac{11}{8}\right\}\)
(equivalently \(\lambda^\ast\in\left\{1,\ \tfrac{\sqrt5}{2},\ \tfrac{\sqrt{22}}{4}\right\}\)),
and Shor--Laflamme enumerator coefficient lists
\[
A^{(A)}=\Big[1,0,1,\tfrac14,15,\tfrac{21}{2},33,\tfrac{13}{4}\Big],\quad
B^{(A)}=\Big[1,0,1,21,32,99,58,44\Big],
\]
\[
A^{(B)}=\Big[1,0,\tfrac54,\tfrac14,\tfrac{29}{2},\tfrac{21}{2},\tfrac{133}{4},\tfrac{13}{4}\Big],\quad
B^{(B)}=\Big[1,0,\tfrac54,\tfrac{87}{4},\tfrac{63}{2},\tfrac{195}{2},\tfrac{233}{4},\tfrac{179}{4}\Big],
\]
\[
A^{(C)}=\Big[1,0,\tfrac{11}{8},\tfrac14,\tfrac{57}{4},\tfrac{21}{2},\tfrac{267}{8},\tfrac{13}{4}\Big],\quad
B^{(C)}=\Big[1,0,\tfrac{11}{8},\tfrac{177}{8},\tfrac{125}{4},\tfrac{387}{4},\tfrac{467}{8},\tfrac{361}{8}\Big].
\]
\item \(\vec a=(1,1,1,2,2,3,5)\): an explicit closed-form self-complementary \(((7,2,3))\) BD16 solution at level \(m=8\) (with \((b_0,b_1)=(0,7)\) and \(\ket{1_L}=X^{\otimes 7}\ket{0_L}\)), with
\[
(\lambda^\ast)^2=\frac{369}{128}\qquad\text{(so }\lambda^\ast=\frac{3\sqrt{82}}{16}\text{)}
\]
and weight enumerators
\[
A(z)=1+\frac{369}{128}z^2+\frac{975}{64}z^4+\frac{5745}{128}z^6,\qquad
B(z)=1+\frac{369}{128}z^2+\frac{3795}{128}z^3+\frac{975}{64}z^4+\frac{6957}{64}z^5+\frac{5745}{128}z^6+\frac{6867}{128}z^7.
\]
\item \(\vec a=(1,1,1,1,3,4,4)\): a fully analytic \(128\)-element exact \(((7,2,3))\) BD16 family at level \(m=8\) with \((b_0,b_1)=(0,7)\) and \(\ket{1_L}=X^{\otimes 7}\ket{0_L}\), with coefficients in \(\mathbb{Q}(\sqrt3,i)\) (phases in \(\{\pm1,\pm i\}\) and magnitudes in \(\{1/4,\sqrt3/4\}\)); the family has constant invariant
\[
(\lambda^\ast)^2=\frac{31}{8}\qquad\text{(so }\lambda^\ast=\frac{\sqrt{62}}{4}\text{)},
\]
and constant Shor--Laflamme enumerators
\[
A=\Big[1,\ 0,\ \frac{31}{8},\ 0,\ \frac{53}{4},\ 0,\ \frac{367}{8},\ 0\Big],\qquad
B=\Big[1,\ 0,\ \frac{31}{8},\ \frac{261}{8},\ \frac{53}{4},\ \frac{411}{4},\ \frac{367}{8},\ \frac{453}{8}\Big].
\]
\item \(\vec a=(1,1,1,1,3,3,5)\): a continuous family with $(\lambda^\ast)^2=\frac{177}{64}$.
\[
A=\Big[
1,\ 0,\ \frac{177}{64},\ 0,\ \frac{3}{32}\bigl(139+26c\bigr),\ \frac{117}{16}\bigl(1-c\bigr),\ 
\frac{3}{64}\bigl(799+156c\bigr),\ \frac{39}{16}\bigl(1-c\bigr)
\Big],
\]
\[
B=\Big[
1,\ 0,\ \frac{177}{64},\ \frac{3}{64}\bigl(573+52c\bigr),\ \frac{3}{32}\bigl(269-104c\bigr),\ 
\frac{9}{32}\bigl(337+52c\bigr),\ \frac{3}{64}\bigl(1163-208c\bigr),\ \frac{3}{64}\bigl(1085+52c\bigr)
\Big].
\]
\end{itemize}

\subsection{SS--LP candidates not yet analyzed here}

The remaining SS--LP-passing BD16 candidate vectors from Sec.~3 have not been taken through the Full-KL stage in this note (i.e., we have not produced or certified distance-3 codewords for them here):
\[
\begin{aligned}
&(1,1,1,2,2,2,6),\ (0,1,1,2,3,3,5).
\end{aligned}
\]

\subsection{What has been verified for the treated vectors}

For each treated vector listed above, we check:
\begin{itemize}
\item \textbf{SS support:} basis supports satisfy the residue-class condition for the stated \(\vec a\) (with \(m=8\), \(b_0=0\), \(b_1=7\)).
\item \textbf{Full KL:} the Knill--Laflamme equalities hold for all Pauli errors of weight \(\le 2\) (hence distance \(3\)).
\item \textbf{Invariants:} Shor--Laflamme weight enumerators \(A,B\) and the signature norm \(\lambda^\ast\) are recorded.
\end{itemize}

\section{\texorpdfstring{$\vec a=(1,1,1,2,2,2,6)$}{a=(1,1,1,2,2,2,6)}}

Throughout we work at BD16 level \(m=8\) with residues \((b_0,b_1)=(0,7)\), and we restrict to the
\emph{single-syndrome residue-class} structure
\[
\ket{0_L}\subset S_0(\vec a),\qquad \ket{1_L}\subset S_7(\vec a),
\qquad S_7(\vec a)=\overline{S_0(\vec a)}.
\]
When additionally imposing self-complementarity, we take the usual convention
\[
\ket{1_L}=X^{\otimes 7}\ket{0_L}.
\]

\subsection{Case \texorpdfstring{$\vec a=(1,1,1,2,2,2,6)$}{a=(1,1,1,2,2,2,6)}: residue support and LP constraints}

With \(m=8\) and \(b_0=0\), the residue-\(0\) support is
\[
S_0(\vec a)=\Big\{
\underbrace{0000000}_{s_0},
\underbrace{0000011}_{s_1},
\underbrace{0000101}_{s_2},
\underbrace{0001001}_{s_3},
\underbrace{0110001}_{s_4},
\underbrace{0111110}_{s_5},
\underbrace{1010001}_{s_6},
\underbrace{1011110}_{s_7},
\underbrace{1100001}_{s_8},
\underbrace{1101110}_{s_9}
\Big\},
\]
\[
|S_0|=10
\]
and \(S_7(\vec a)\) is the bitwise-complement set.

Assume the self-complement form \(\ket{1_L}=X^{\otimes 7}\ket{0_L}\) and write
\[
\ket{0_L}=\sum_{k=0}^{9} c_k\ket{s_k},\qquad p_k:=|c_k|^2.
\]
Then the \(Z\)-type Knill--Laflamme constraints force the squared magnitudes into a small two-parameter family:
\[
p_0=p_1=p_2=p_3=\frac{1}{16},
\]
and the remaining six are determined by \((u,v)\) via
\[
\begin{aligned}
p_4&=u+v-\frac{3}{16},&
p_5&=\frac{7}{16}-u-v,\\
p_6&=\frac{1}{4}-u,&
p_7&=u,\\
p_8&=\frac{1}{4}-v,&
p_9&=v,
\end{aligned}
\]
with feasibility region
\[
0\le u\le \frac14,\qquad 0\le v\le \frac14,\qquad \frac{3}{16}\le u+v\le \frac{7}{16}.
\]

\subsection{Numerical and discrete evidence for an obstruction}

I pushed the search substantially harder than before in two complementary regimes.

\paragraph{Continuous phases (no restriction).}
I minimized the full distance-3 KL residual over \((u,v)\) in the feasible region and over all relative phases of the
\(10\) amplitudes (fixing a global phase). The best value reached reproducibly was on the order of
\[
4.82\times 10^{-1},
\]
i.e.\ not close to \(0\).

\paragraph{Exhaustive discrete ``analytic-style'' scan.}
I exhaustively checked the restricted ansatz
\[
u,v\in\left\{0,\frac1{16},\frac2{16},\frac3{16},\frac4{16}\right\}
\quad\text{(subject to feasibility, \(18\) pairs total),}
\qquad
\text{phases in }\{\pm 1,\pm i\},
\]
and used the \(XX\)-implied phase restrictions to reduce the naive \(4^9\) phase search per \((u,v)\) to at most \(4096\) cases.
This yields an exhaustive scan of \(73{,}728\) candidates within this sub-ansatz. None satisfied the full KL constraints;
the best residual observed was about \(0.5\).

\paragraph{Beyond self-complementarity.}
I also re-confirmed (in smaller runs) the more general case where \(\ket{0_L}\) and \(\ket{1_L}\) have independent amplitudes on
\(S_0(\vec a)\) and \(S_7(\vec a)\) (dropping the \(X^{\otimes 7}\) relation). The best residual remained around
\[
3.5\times 10^{-1}.
\]

Taken together, the evidence strongly suggests that \(\vec a=(1,1,1,2,2,2,6)\) is genuinely obstructed
within the BD16 single-syndrome residue-class setting; in particular, the discrete
\(\{\pm 1,\pm i\}\) phase-only ansatz with \((u,v)\in\frac1{16}\mathbb{Z}\) is now exhaustively ruled out.

\subsection{No full-KL solution for $\vec a=(1,1,1,2,2,2,6)$ under the complementary ansatz}

Let $m=8$, $(b_0,b_1)=(0,7)$ and $\vec a=(1,1,1,2,2,2,6)$.
Assume the complementary convention $\ket{1_L}=X^{\otimes 7}\ket{0_L}$ and $\text{supp}(\ket{0_L})\subseteq S_0(\vec a)$.

A direct subset-sum check gives
\[
S_0=\{0000000,0000011,0000101,0001001,0110001,0111110,1010001,1011110,1100001,1101110\}.
\]
Write $\ket{0_L}=\sum_{s\in S_0} c_s\ket{s}$ and $p_s:=|c_s|^2$.

\paragraph{Z-type KL.}
For each $j$, KL for $Z_j$ implies $\langle 0_L|Z_j|0_L\rangle=0$, hence a linear system in $(p_s)_{s\in S_0}$.
Solving it yields
\[
p_{0000000}=p_{0000011}=p_{0000101}=p_{0001001}=\frac{1}{16},
\]
and with $u:=p_{1011110}$, $v:=p_{1101110}$,
\[
p_{1010001}=\frac14-u,\quad p_{1100001}=\frac14-v,\quad
p_{0110001}=u+v-\frac{3}{16},\quad p_{0111110}=\frac{7}{16}-u-v.
\]

\paragraph{Off-diagonal KL on $X_1,X_2,X_3$.}
Let
\[
A=1010001,\ B=1101110,\ C=1011110,\ D=1100001,\ E=0110001,\ F=0111110,
\]
and abbreviate $(a,b,c,d,e,f):=(c_A,c_B,c_C,c_D,c_E,c_F)$.
From the SS structure one checks the induced pairings:
$A\leftrightarrow B$, $C\leftrightarrow D$ under $X_1$;
$E\leftrightarrow B$, $F\leftrightarrow D$ under $X_2$;
$E\leftrightarrow C$, $F\leftrightarrow A$ under $X_3$.
Using the KL constraints for $X_i$ and for suitable $Z_kX_i$ (weight $2$) one deduces that
$b^*a,d^*c,b^*e,d^*f,c^*e,a^*f\in\mathbb R$ and
\[
b^*a=-d^*c,\qquad b^*e=-d^*f,\qquad c^*e=-a^*f.
\]
Since all these products are real and connect the six amplitudes, after a global phase we may take
$a,b,c,d,e,f\in\mathbb R$, so the relations become
\[
ba=-dc,\qquad be=-df,\qquad ce=-af.
\]
Multiplying the first by $f$ and the second by $c$ gives $baf=bce$, hence (as $b\neq0$) $af=ce$,
but the third gives $ce=-af$, so $af=0$.

\paragraph{Nonzero check.}
From the Z-type parametrization and the magnitude consequences of $ba=-dc$ and $be=-df$ one checks
that $a,b,f\neq 0$ (otherwise one forces negative probabilities such as $p_F<0$ or $p_E<0$).
Thus $af=0$ is impossible, contradicting full KL.

Therefore no full-KL \(((7,2,3))\) code exists for $\vec a=(1,1,1,2,2,2,6)$ within the complementary SS ansatz.

\subsection{Lean-friendly no--go for full Knill--Laflamme under the complementary ansatz:
$\vec a=(1,1,1,2,2,2,6)$ (mod $8$)}
\label{subsec:lean-nogo-1112226}

\paragraph{Goal (pure finite-dimensional linear algebra).}
Let $\Bit7:=\{0,1\}^7$ and write a bitstring as $s=s_1\cdots s_7\in\Bit7$.
Define bitwise XOR by $(s\oplus t)_i := s_i+t_i \pmod 2$ and bitwise complement $\bar s:=s\oplus 1^7$.

Fix the following \emph{explicit} finite set (a list of $10$ bitstrings):
\[
S_0:=\{s_0,\dots,s_9\}\subset\Bit7,
\]
where
\[
\begin{aligned}
&s_0=0000000,\ s_1=0000011,\ s_2=0000101,\ s_3=0001001,\ s_4=0110001,\\
&s_5=0111110,\ s_6=1010001,\ s_7=1011110,\ s_8=1100001,\ s_9=1101110.
\end{aligned}
\]
Let $S_7:=\overline{S_0}:=\{\bar s:\ s\in S_0\}$.

Let $(c_0,\dots,c_9)\in\mathbb C^{10}$ be unknown amplitudes and define two vectors in
$\mathbb C^{\Bit7}$ by their coordinates in the standard basis $\{\ket{s}:s\in\Bit7\}$:
\begin{equation}
\ket{0_L}:=\sum_{k=0}^9 c_k\,\ket{s_k},
\qquad
\ket{1_L}:=\sum_{k=0}^9 c_k\,\ket{\overline{s_k}}.
\label{eq:lean-def-01}
\end{equation}
Equivalently, $\ket{1_L}=X^{\otimes 7}\ket{0_L}$, i.e.\ $\ket{1_L}$ is obtained by complementing every bit.
Assume normalization $\braket{0_L}{0_L}=1$, i.e.
\begin{equation}
\sum_{k=0}^9 |c_k|^2=1.
\label{eq:lean-normalization}
\end{equation}

\paragraph{Pauli operators as explicit signed permutations.}
For $x,z\in\Bit7$, define a linear operator $P(x,z):\mathbb C^{\Bit7}\to\mathbb C^{\Bit7}$ by its action on basis vectors
\begin{equation}
P(x,z)\ket{s} := i^{x\cdot z}\,(-1)^{z\cdot s}\,\ket{s\oplus x},
\label{eq:lean-pauli-action}
\end{equation}
where $x\cdot z:=\sum_i x_i z_i\in\{0,\dots,7\}$ (in $\mathbb Z$) and $z\cdot s:=\sum_i z_i s_i\pmod 2$.
Define the Pauli weight
\[
\wt(x,z):=\bigl|\supp(x)\cup\supp(z)\bigr|
\quad\text{where}\quad
\supp(x):=\{i:\ x_i=1\}.
\]

\paragraph{Full distance-$3$ Knill--Laflamme system.}
The \emph{full} KL system we want to solve is:
\begin{equation}
\boxed{
\begin{aligned}
&\forall\, x,z\in\{0,1\}^7,\ \mathrm{wt}(x,z)\le 2:\quad
\bra{0_L}P(x,z)\ket{1_L}=0,\\
&\hspace{4.1cm}
\bra{0_L}P(x,z)\ket{0_L}=\bra{1_L}P(x,z)\ket{1_L}.
\end{aligned}}
\label{eq:lean-full-KL}
\end{equation}
Everything here is standard finite-dimensional linear algebra over $\mathbb C$ with an explicit basis.

\paragraph{What we prove.}
We prove that the finite system \eqref{eq:lean-full-KL} has \emph{no} solution $(c_0,\dots,c_9)\in\mathbb C^{10}$.
Since \eqref{eq:lean-full-KL} implies every \emph{subsystem}, it suffices to derive a contradiction from a small subset of
the equations in \eqref{eq:lean-full-KL}.

%------------------------------------------------------------
\subsubsection*{Step 1: the $Z_j$ diagonal equations force a 2-parameter family of probabilities}

Define $p_k:=|c_k|^2\in\mathbb R_{\ge 0}$.

\begin{lemma}[Diagonal $Z$-constraints]
\label{lem:lean-Z-system}
Assume \eqref{eq:lean-full-KL}. Then for each $j\in\{1,\dots,7\}$ one has
\begin{equation}
\bra{0_L}Z_j\ket{0_L}=0
\qquad\Longleftrightarrow\qquad
\sum_{\substack{0\le k\le 9\\ (s_k)_j=1}} p_k=\frac12,
\label{eq:lean-Zj}
\end{equation}
together with normalization $\sum_{k=0}^9 p_k=1$.
Solving these eight linear equations yields
\begin{equation}
p_0=p_1=p_2=p_3=\frac{1}{16},
\label{eq:lean-four-fixed}
\end{equation}
and, writing
\[
u:=p_7=p_{1011110},\qquad v:=p_9=p_{1101110},
\]
the remaining six probabilities are
\begin{equation}
p_6=\frac14-u,\quad p_8=\frac14-v,\quad
p_4=u+v-\frac{3}{16},\quad p_5=\frac{7}{16}-u-v.
\label{eq:lean-Z-solution}
\end{equation}
In particular $p_k\ge 0$ forces
\[
0\le u\le \frac14,\qquad 0\le v\le \frac14,\qquad \frac{3}{16}\le u+v\le \frac{7}{16}.
\]
\end{lemma}

\begin{proof}
For $Z_j=P(0,e_j)$, one has $Z_j\ket{s}=(-1)^{s_j}\ket{s}$, so
$\bra{0_L}Z_j\ket{0_L}=\sum_{k=0}^9 (-1)^{(s_k)_j}p_k$ and $\bra{1_L}Z_j\ket{1_L}=-\bra{0_L}Z_j\ket{0_L}$ because complementation flips the parity in the exponent.
Thus $\bra{0_L}Z_j\ket{0_L}=\bra{1_L}Z_j\ket{1_L}$ forces $\bra{0_L}Z_j\ket{0_L}=0$, giving \eqref{eq:lean-Zj}.
The linear solve \eqref{eq:lean-four-fixed}--\eqref{eq:lean-Z-solution} is a routine elimination (Lean target: \texttt{linarith} after expanding \eqref{eq:lean-Zj} for $j=1,\dots,7$ and \eqref{eq:lean-normalization}). 
\end{proof}

%------------------------------------------------------------
\subsubsection*{Step 2: nine explicit off-diagonal equations force three ``real product'' relations}

We now name the six amplitudes that appear in the chosen off-diagonal equations:
\[
\begin{aligned}
&a:=c_{1010001}=c_6,\qquad b:=c_{1101110}=c_9,\qquad c:=c_{1011110}=c_7,\\
&d:=c_{1100001}=c_8,\qquad e:=c_{0110001}=c_4,\qquad f:=c_{0111110}=c_5.
\end{aligned}
\]

\begin{lemma}[The $x=e_1$ subsystem]
\label{lem:lean-x1}
Assume \eqref{eq:lean-full-KL}. Consider the three weight$\le 2$ operators
\[
X_1,\qquad Z_2X_1,\qquad Z_4X_1.
\]
Then their off-diagonal KL constraints $\bra{0_L}(\cdot)\ket{1_L}=0$ expand to the three equations
\begin{align}
&\overline b\,a+\overline c\,d+\overline a\,b+\overline d\,c=0,
\label{eq:lean-x1-0}\\
&-\overline b\,a+\overline c\,d+\overline a\,b-\overline d\,c=0,
\label{eq:lean-x1-2}\\
&-\overline b\,a-\overline c\,d+\overline a\,b+\overline d\,c=0.
\label{eq:lean-x1-4}
\end{align}
Consequently, defining $u_1:=\overline b\,a$ and $v_1:=\overline d\,c$, one deduces
\begin{equation}
u_1\in\mathbb R,\qquad v_1\in\mathbb R,\qquad v_1=-u_1,
\quad\text{i.e.}\quad
\overline b\,a=-\overline d\,c\in\mathbb R.
\label{eq:lean-prod1}
\end{equation}
\end{lemma}

\begin{proof}
The expansion is a direct finite sum computation using \eqref{eq:lean-pauli-action} and \eqref{eq:lean-def-01}.
(Lean implementation: expand $\bra{0_L}P(x,z)\ket{1_L}$ using the orthonormal basis; only terms with matching basis vectors survive.)

To solve the three equations, note that $\overline a\,b=\overline{(\overline b\,a)}=\overline{u_1}$
and $\overline c\,d=\overline{(\overline d\,c)}=\overline{v_1}$. Then \eqref{eq:lean-x1-0} becomes
$u_1+\overline{u_1}+v_1+\overline{v_1}=0$, i.e.\ $\Re(u_1+v_1)=0$.
Equation \eqref{eq:lean-x1-2} is $-(u_1+v_1)+\overline{(u_1+v_1)}=0$, i.e.\ $u_1+v_1\in\mathbb R$.
Hence $u_1+v_1=0$ so $v_1=-u_1$.
Finally \eqref{eq:lean-x1-4} becomes $(\overline{u_1}-u_1)+ (v_1-\overline{v_1})=0$,
so $\Im(u_1)=\Im(v_1)$. With $v_1=-u_1$ this forces $\Im(u_1)=0$, hence $u_1\in\mathbb R$ and $v_1=-u_1\in\mathbb R$.
\end{proof}

\begin{lemma}[The $x=e_2$ subsystem]
\label{lem:lean-x2}
Assume \eqref{eq:lean-full-KL}. Consider
\[
X_2,\qquad Z_1X_2,\qquad Z_4X_2.
\]
Their off-diagonal KL constraints expand to
\begin{align}
&\overline f\,d+\overline d\,f+\overline b\,e+\overline e\,b=0,
\label{eq:lean-x2-0}\\
&\overline f\,d-\overline d\,f-\overline b\,e+\overline e\,b=0,
\label{eq:lean-x2-1}\\
&-\overline f\,d+\overline d\,f-\overline b\,e+\overline e\,b=0.
\label{eq:lean-x2-4}
\end{align}
Consequently,
\begin{equation}
\overline b\,e=-\overline d\,f\in\mathbb R.
\label{eq:lean-prod2}
\end{equation}
\end{lemma}

\begin{proof}
Same proof pattern as Lemma~\ref{lem:lean-x1}, with $u_2:=\overline b\,e$ and $v_2:=\overline d\,f$.
\end{proof}

\begin{lemma}[The $x=e_3$ subsystem]
\label{lem:lean-x3}
Assume \eqref{eq:lean-full-KL}. Consider
\[
X_3,\qquad Z_1X_3,\qquad Z_4X_3.
\]
Their off-diagonal KL constraints expand to
\begin{align}
&\overline f\,a+\overline a\,f+\overline c\,e+\overline e\,c=0,
\label{eq:lean-x3-0}\\
&\overline f\,a-\overline a\,f-\overline c\,e+\overline e\,c=0,
\label{eq:lean-x3-1}\\
&-\overline f\,a+\overline a\,f-\overline c\,e+\overline e\,c=0.
\label{eq:lean-x3-4}
\end{align}
Consequently,
\begin{equation}
\overline c\,e=-\overline a\,f\in\mathbb R.
\label{eq:lean-prod3}
\end{equation}
\end{lemma}

\begin{proof}
Same pattern, with $u_3:=\overline c\,e$ and $v_3:=\overline a\,f$.
\end{proof}

%------------------------------------------------------------
\subsubsection*{Step 3: the product relations force a phase alignment and then an algebraic contradiction}

\begin{lemma}[Nonzero amplitudes forced by \eqref{eq:lean-Z-solution} and \eqref{eq:lean-prod1}--\eqref{eq:lean-prod3}]
\label{lem:lean-nonzero}
Assume \eqref{eq:lean-full-KL}. Then
\[
a\neq 0,\qquad b\neq 0,\qquad c\neq 0,\qquad d\neq 0,\qquad e\neq 0,\qquad f\neq 0.
\]
\end{lemma}

\begin{proof}
We use Lemma~\ref{lem:lean-Z-system} and Lemmas~\ref{lem:lean-x1}--\ref{lem:lean-x3}.

\emph{(1) $b\neq 0$.}
If $b=0$ then $p_9=|b|^2=v=0$, so by \eqref{eq:lean-Z-solution} we get $p_8=\frac14-v=\frac14$, hence $d\neq 0$.
But \eqref{eq:lean-prod1} gives $0=\overline b\,a=-\overline d\,c$, so $c=0$ and thus $u=p_7=|c|^2=0$.
Then \eqref{eq:lean-Z-solution} gives $p_4=u+v-\frac{3}{16}=-\frac{3}{16}<0$, contradicting $p_4=|e|^2\ge 0$.

\emph{(2) $e\neq 0$.}
If $e=0$ then $p_4=0$, so \eqref{eq:lean-Z-solution} gives $u+v=\frac{3}{16}$ and hence
$p_5=\frac{7}{16}-u-v=\frac14$, so $f\neq 0$.
Now \eqref{eq:lean-prod2} gives $0=\overline b\,e=-\overline d\,f$, so $d=0$, i.e.\ $p_8=0$ and thus $v=\frac14$.
But then $u=\frac{3}{16}-v=-\frac{1}{16}<0$, impossible since $u=p_7=|c|^2\ge 0$.

\emph{(3) $c\neq 0$.}
If $c=0$ then $u=p_7=0$. By \eqref{eq:lean-prod1}, $\overline b\,a=0$, and since $b\neq 0$ we get $a=0$.
But $a=0$ means $p_6=|a|^2=0$, so \eqref{eq:lean-Z-solution} gives $\frac14-u=0$, i.e.\ $u=\frac14$, contradicting $u=0$.

\emph{(4) $d\neq 0$.}
If $d=0$ then $v=\frac14$ and \eqref{eq:lean-prod2} gives $\overline b\,e=0$, hence $e=0$ since $b\neq 0$,
contradicting (2).

\emph{(5) $a\neq 0$.}
If $a=0$ then $p_6=0$, so \eqref{eq:lean-Z-solution} gives $u=\frac14$ and thus $c\neq 0$.
But \eqref{eq:lean-prod3} gives $0=\overline a\,f=-\overline c\,e$, hence $e=0$, contradicting (2).

\emph{(6) $f\neq 0$.}
If $f=0$ then $p_5=0$ so \eqref{eq:lean-Z-solution} gives $u+v=\frac{7}{16}$ and hence $p_4=\frac14$, so $e\neq 0$.
But \eqref{eq:lean-prod2} gives $0=-\overline d\,f=\overline b\,e$, so $e=0$ since $b\neq 0$, contradiction.
\end{proof}

\begin{lemma}[Phase-fixing and reduction to real variables]
\label{lem:lean-phase-fix}
Assume \eqref{eq:lean-full-KL}. Then, after multiplying all $c_k$ by a single complex scalar $\omega$ with $|\omega|=1$
(global phase), we may assume
\[
a,b,c,d,e,f\in\mathbb R
\]
while preserving all equations in Lemmas~\ref{lem:lean-Z-system}--\ref{lem:lean-x3}.
\end{lemma}

\begin{proof}
By Lemma~\ref{lem:lean-nonzero}, $b\neq 0$. Let $\omega:=\overline b/|b|$, so $|\omega|=1$ and $\omega b=|b|\in\mathbb R_{>0}$.
Replacing each $c_k$ by $\omega c_k$ preserves:
(i) each $p_k=|c_k|^2$ (hence Lemma~\ref{lem:lean-Z-system}),
and (ii) each bilinear form $\overline{c_\alpha}c_\beta$ (hence Lemmas~\ref{lem:lean-x1}--\ref{lem:lean-x3}),
because $\overline{\omega c_\alpha}\,(\omega c_\beta)=\overline\omega\,\omega\,\overline{c_\alpha}c_\beta=\overline{c_\alpha}c_\beta$.

Now assume $b\in\mathbb R_{>0}$. From \eqref{eq:lean-prod1}, $\overline b\,a=ba\in\mathbb R$ and $b\neq 0$, so $a\in\mathbb R$.
From \eqref{eq:lean-prod2}, $\overline b\,e=be\in\mathbb R$ and $b\neq 0$, so $e\in\mathbb R$.
From \eqref{eq:lean-prod3}, $\overline c\,e\in\mathbb R$ and $e\neq 0$ (Lemma~\ref{lem:lean-nonzero}), so $\overline c\in\mathbb R$ hence $c\in\mathbb R$.
From \eqref{eq:lean-prod1}, $\overline d\,c\in\mathbb R$ and $c\neq 0$, so $d\in\mathbb R$.
From \eqref{eq:lean-prod2}, $\overline d\,f\in\mathbb R$ and $d\neq 0$, so $f\in\mathbb R$.
\end{proof}

\begin{lemma}[Final algebraic contradiction]
\label{lem:lean-final-contradiction}
Assume \eqref{eq:lean-full-KL}. Then $af=0$, contradicting Lemma~\ref{lem:lean-nonzero}.
\end{lemma}

\begin{proof}
Apply Lemma~\ref{lem:lean-phase-fix} and work in the phase where $a,b,c,d,e,f\in\mathbb R$.
Then \eqref{eq:lean-prod1}--\eqref{eq:lean-prod3} become
\begin{equation}
ba=-dc,\qquad be=-df,\qquad ce=-af.
\label{eq:lean-real-relations}
\end{equation}
Multiply $ba=-dc$ by $f$ and $be=-df$ by $c$:
\[
baf=-dcf,\qquad bce=-dfc=-dcf.
\]
Hence $baf=bce$, and since $b\neq 0$ we obtain $af=ce$.
But $ce=-af$ from \eqref{eq:lean-real-relations}, so $af=-af$ and therefore $af=0$.
This contradicts $a\neq 0$ and $f\neq 0$ from Lemma~\ref{lem:lean-nonzero}.
\end{proof}

\begin{theorem}[No solution to the full KL system for $S_0$ above]
\label{thm:lean-nogo-1112226}
There is no choice of amplitudes $(c_0,\dots,c_9)\in\mathbb C^{10}$ such that
$\ket{0_L},\ket{1_L}$ defined by \eqref{eq:lean-def-01} satisfy the full system \eqref{eq:lean-full-KL}.
\end{theorem}

\begin{proof}
Assume a solution exists. Then Lemma~\ref{lem:lean-Z-system} holds, as do Lemmas~\ref{lem:lean-x1}--\ref{lem:lean-x3}.
Lemma~\ref{lem:lean-nonzero} shows $a,f\neq 0$.
Lemma~\ref{lem:lean-final-contradiction} then yields $af=0$, a contradiction.
\end{proof}

\begin{remark}[Minimal KL subset used]
The no-go uses only:
(i) the seven diagonal constraints for $Z_1,\dots,Z_7$ plus normalization,
and (ii) nine off-diagonal constraints for
$X_1,Z_2X_1,Z_4X_1$, $X_2,Z_1X_2,Z_4X_2$, $X_3,Z_1X_3,Z_4X_3$.
All are weight $\le 2$ Paulis, hence are part of \eqref{eq:lean-full-KL}.
\end{remark}

%============================================================
\subsection{A purely algebraic no--go for the complementary residue support
($S_0$ of size $10$ for $\vec a=(1,1,1,2,2,2,6)$ mod $8$)}
\label{subsec:pure-math-nogo-1112226}

\paragraph{Goal.}
This subsection is written in \emph{pure finite-dimensional linear algebra}:
we define an explicit finite-dimensional complex inner product space, an explicit finite list of
basis vectors, an explicit family of \emph{signed permutation} operators, and an explicit finite system
of bilinear equations (a subsystem of the full distance-$3$ Knill--Laflamme constraints).
We prove that this subsystem has no solution, hence the full distance-$3$ system has no solution either.

%------------------------------------------------------------
\subsubsection*{Data, variables, and vectors}

\paragraph{The ambient space.}
Let $B:=\{0,1\}^7$ (equivalently $\mathbb F_2^7$).
Let $V:=\mathbb C^{B}$, i.e.\ the vector space of functions $B\to\mathbb C$.
Write the standard orthonormal basis as $\{e_s:s\in B\}$ where $e_s(t)=\mathbf 1_{t=s}$, and
use the standard inner product
\[
\langle u,v\rangle:=\sum_{s\in B}\overline{u(s)}\,v(s).
\]
Define bitwise XOR $s\oplus t\in B$ by $(s\oplus t)_i=s_i+t_i\pmod 2$ and complement $\bar s:=s\oplus 1^7$.

\paragraph{The fixed support set.}
Fix the explicit subset $S_0=\{s_0,\dots,s_9\}\subset B$:
\[
\begin{aligned}
&s_0=0000000,\quad s_1=0000011,\quad s_2=0000101,\quad s_3=0001001,\quad s_4=0110001,\\
&s_5=0111110,\quad s_6=1010001,\quad s_7=1011110,\quad s_8=1100001,\quad s_9=1101110.
\end{aligned}
\]
Let $S_7:=\overline{S_0}:=\{\bar s:\ s\in S_0\}$.

\paragraph{Unknown coefficients and the two vectors.}
Let $c_0,\dots,c_9\in\mathbb C$ be unknown scalars and define
\begin{equation}
\psi_0 := \sum_{k=0}^9 c_k\, e_{s_k}\in V,
\qquad
\psi_1 := \sum_{k=0}^9 c_k\, e_{\overline{s_k}}\in V.
\label{eq:pure-def-psi01}
\end{equation}
Define the \emph{squared magnitudes} (nonnegative reals)
\[
p_k:=|c_k|^2\in\mathbb R_{\ge 0}\qquad (k=0,\dots,9).
\]

%------------------------------------------------------------
\subsubsection*{The signed-permutation operators and the equation system}

\paragraph{Two families of basic operators.}
For each coordinate index $i\in\{1,\dots,7\}$ define:
\begin{itemize}
\item The \emph{sign operator} $D_i:V\to V$ by
\[
D_i(e_s):=(-1)^{s_i}\,e_s.
\]
This is diagonal in the basis $\{e_s\}$ with entries $\pm 1$.
\item The \emph{flip operator} $F_i:V\to V$ by
\[
F_i(e_s):=e_{s\oplus e_i},
\]
where $e_i\in B$ is the standard basis vector with a single $1$ in position $i$.
This is a permutation matrix in the basis $\{e_s\}$.
\end{itemize}
We also consider their compositions $D_j\circ F_i$, which are again signed permutation operators.

\paragraph{The equation system (a subsystem implied by full distance-$3$ constraints).}
We consider the following finite system of equations in the unknowns $c_0,\dots,c_9$:
\begin{equation}
\label{eq:pure-eq-system}
\boxed{
\begin{aligned}
&\text{(Norm)}\quad &&\langle \psi_0,\psi_0\rangle=1.\\[1mm]
&\text{(Seven diagonal equations)}\quad &&\forall i\in\{1,\dots,7\}:\ \ \langle \psi_0,D_i\psi_0\rangle=\langle \psi_1,D_i\psi_1\rangle.\\[1mm]
&\text{(Nine off-diagonal equations)}\quad
&&\langle \psi_0, F_1\psi_1\rangle=\langle \psi_0, (D_2\!\circ F_1)\psi_1\rangle=\langle \psi_0,(D_4\!\circ F_1)\psi_1\rangle=0,\\
&&&\langle \psi_0, F_2\psi_1\rangle=\langle \psi_0, (D_1\!\circ F_2)\psi_1\rangle=\langle \psi_0,(D_4\!\circ F_2)\psi_1\rangle=0,\\
&&&\langle \psi_0, F_3\psi_1\rangle=\langle \psi_0, (D_1\!\circ F_3)\psi_1\rangle=\langle \psi_0,(D_4\!\circ F_3)\psi_1\rangle=0.
\end{aligned}}
\end{equation}
Every operator appearing in \eqref{eq:pure-eq-system} is a signed permutation matrix and each equation is a polynomial
constraint of degree $2$ in the real and imaginary parts of $(c_0,\dots,c_9)$.

\paragraph{What we prove.}
\begin{theorem}[Unsatisfiability of the subsystem]
\label{thm:pure-nogo-subsystem}
There do not exist complex numbers $c_0,\dots,c_9\in\mathbb C$ such that the vectors
$\psi_0,\psi_1$ from \eqref{eq:pure-def-psi01} satisfy all equations in \eqref{eq:pure-eq-system}.
\end{theorem}

Since the \emph{full} distance-$3$ condition implies \eqref{eq:pure-eq-system}, this already rules out
any full distance-$3$ solution in the complementary support model.

%------------------------------------------------------------
\subsubsection*{Proof of Theorem~\ref{thm:pure-nogo-subsystem}}

\begin{proof}
Assume for contradiction that $c_0,\dots,c_9$ satisfy \eqref{eq:pure-eq-system}.

%------------------------------------------------------------
\paragraph{Step 1: the seven diagonal equations reduce to a linear system in $p_k=|c_k|^2$.}

\begin{lemma}[Diagonal sign equations become linear constraints on $p_k$]
\label{lem:pure-diag-linear}
For each $i\in\{1,\dots,7\}$ one has
\[
\langle \psi_0,D_i\psi_0\rangle=\sum_{k=0}^9 (-1)^{(s_k)_i}\,p_k,
\qquad
\langle \psi_1,D_i\psi_1\rangle=-\sum_{k=0}^9 (-1)^{(s_k)_i}\,p_k,
\]
hence the diagonal equation $\langle \psi_0,D_i\psi_0\rangle=\langle \psi_1,D_i\psi_1\rangle$ is equivalent to
\begin{equation}
\sum_{k=0}^9 (-1)^{(s_k)_i}\,p_k = 0.
\label{eq:pure-diag-eq}
\end{equation}
Together with $\sum_{k=0}^9 p_k=1$ (normalization), this is an $8\times 10$ linear system in the $p_k$.
\end{lemma}

\begin{proof}
The first identity is immediate from $D_i(e_{s_k})=(-1)^{(s_k)_i}e_{s_k}$ and orthonormality of the basis.
For the second identity, note that $(\overline{s_k})_i = 1-(s_k)_i$, hence $(-1)^{(\overline{s_k})_i}=-(-1)^{(s_k)_i}$,
so the diagonal expectation changes sign when we pass from $S_0$ to its complement $S_7$.
\end{proof}

\begin{lemma}[Explicit solution of the diagonal linear system]
\label{lem:pure-diag-solve}
Assuming the $7$ diagonal equations \eqref{eq:pure-diag-eq} for $i=1,\dots,7$ and $\sum_{k=0}^9 p_k=1$,
one obtains
\begin{equation}
p_0=p_1=p_2=p_3=\frac{1}{16}.
\label{eq:pure-four-fixed}
\end{equation}
Moreover, if we set
\[
U:=p_7\quad\text{and}\quad V:=p_9,
\]
then the remaining $p_k$ are determined by
\begin{equation}
p_6=\frac14-U,\quad p_8=\frac14-V,\quad
p_4=U+V-\frac{3}{16},\quad p_5=\frac{7}{16}-U-V,
\label{eq:pure-diag-param}
\end{equation}
together with $p_7=U$ and $p_9=V$.
In particular, $p_k\ge 0$ forces
\[
0\le U\le\frac14,\qquad 0\le V\le\frac14,\qquad \frac{3}{16}\le U+V\le\frac{7}{16}.
\]
\end{lemma}

\begin{proof}
This is routine linear elimination on the explicit $8$ equations.
(It is a good target for computer-checked linear arithmetic: after expanding the $7$ equations
$\sum_k (-1)^{(s_k)_i}p_k=0$ using the explicit list of $s_k$, one can eliminate variables and obtain
\eqref{eq:pure-four-fixed}--\eqref{eq:pure-diag-param} by $\texttt{linarith}$.)
\end{proof}

%------------------------------------------------------------
\paragraph{Step 2: the nine off-diagonal equations yield three ``real product'' relations among six coefficients.}

\paragraph{Naming the six coefficients that appear.}
Let
\[
a:=c_6\ (s_6=1010001),\quad
c:=c_7\ (s_7=1011110),\quad
d:=c_8\ (s_8=1100001),\quad
b:=c_9\ (s_9=1101110),\quad
e:=c_4\ (s_4=0110001),\quad
f:=c_5\ (s_5=0111110).
\]
(So $p_6=|a|^2$, $p_7=|c|^2$, $p_8=|d|^2$, $p_9=|b|^2$, $p_4=|e|^2$, $p_5=|f|^2$.)

\begin{lemma}[The three $F_1$-equations imply $\overline{b}a=-\overline{d}c\in\mathbb R$]
\label{lem:pure-prod1}
Assume the three off-diagonal equations
\[
\langle \psi_0,F_1\psi_1\rangle=\langle \psi_0,(D_2\!\circ F_1)\psi_1\rangle=\langle \psi_0,(D_4\!\circ F_1)\psi_1\rangle=0.
\]
Then
\begin{equation}
\overline{b}\,a=-\overline{d}\,c\in\mathbb R.
\label{eq:pure-prod1}
\end{equation}
\end{lemma}

\begin{proof}
A direct basis expansion shows that only the pairings
\[
s_6\leftrightarrow s_9
\quad\text{and}\quad
s_7\leftrightarrow s_8
\]
contribute to $\langle \psi_0,F_1\psi_1\rangle$ (this is a finite check using the explicit $S_0$).
Concretely one finds
\[
\langle \psi_0,F_1\psi_1\rangle
=
\overline{b}a+\overline{c}d+\overline{a}b+\overline{d}c.
\]
Similarly, the sign operator $D_2$ flips the sign of exactly those terms whose \emph{row index} in $S_0$ has second bit $1$,
and $D_4$ flips the sign of those whose row index has fourth bit $1$; with the explicit bit patterns of $s_6,s_7,s_8,s_9$ this yields:
\begin{align*}
0&=\langle \psi_0,F_1\psi_1\rangle
=\overline{b}a+\overline{c}d+\overline{a}b+\overline{d}c,\\
0&=\langle \psi_0,(D_2\!\circ F_1)\psi_1\rangle
=-\overline{b}a+\overline{c}d+\overline{a}b-\overline{d}c,\\
0&=\langle \psi_0,(D_4\!\circ F_1)\psi_1\rangle
=-\overline{b}a-\overline{c}d+\overline{a}b+\overline{d}c.
\end{align*}
Now set $u_1:=\overline{b}a$ and $v_1:=\overline{d}c$. Note that $\overline{a}b=\overline{u_1}$ and
$\overline{c}d=\overline{v_1}$. The first equation becomes
$u_1+\overline{u_1}+v_1+\overline{v_1}=0$, i.e.\ $\Re(u_1+v_1)=0$.
The second equation becomes $-(u_1+v_1)+\overline{(u_1+v_1)}=0$, i.e.\ $u_1+v_1\in\mathbb R$.
Hence $u_1+v_1=0$ and $v_1=-u_1$.
Substitute into the third equation to obtain $\Im(u_1)=0$, hence $u_1\in\mathbb R$ and $v_1=-u_1\in\mathbb R$.
\end{proof}

\begin{lemma}[The three $F_2$-equations imply $\overline{b}e=-\overline{d}f\in\mathbb R$]
\label{lem:pure-prod2}
Assume the three off-diagonal equations
\[
\langle \psi_0,F_2\psi_1\rangle=\langle \psi_0,(D_1\!\circ F_2)\psi_1\rangle=\langle \psi_0,(D_4\!\circ F_2)\psi_1\rangle=0.
\]
Then
\begin{equation}
\overline{b}\,e=-\overline{d}\,f\in\mathbb R.
\label{eq:pure-prod2}
\end{equation}
\end{lemma}

\begin{proof}
Exactly the same algebra as Lemma~\ref{lem:pure-prod1}, after checking (by finite enumeration using $S_0$)
that the only contributing pairings are $s_4\leftrightarrow s_9$ and $s_5\leftrightarrow s_8$.
\end{proof}

\begin{lemma}[The three $F_3$-equations imply $\overline{c}e=-\overline{a}f\in\mathbb R$]
\label{lem:pure-prod3}
Assume the three off-diagonal equations
\[
\langle \psi_0,F_3\psi_1\rangle=\langle \psi_0,(D_1\!\circ F_3)\psi_1\rangle=\langle \psi_0,(D_4\!\circ F_3)\psi_1\rangle=0.
\]
Then
\begin{equation}
\overline{c}\,e=-\overline{a}\,f\in\mathbb R.
\label{eq:pure-prod3}
\end{equation}
\end{lemma}

\begin{proof}
Same proof pattern, with the (finite) pairing check giving the contributing pairs $s_4\leftrightarrow s_7$ and $s_5\leftrightarrow s_6$.
\end{proof}

%------------------------------------------------------------
\paragraph{Step 3: show $a,b,c,d,e,f$ are all nonzero.}

\begin{lemma}[Forced nonvanishing]
\label{lem:pure-nonzero}
Under the diagonal parametrization of Lemma~\ref{lem:pure-diag-solve} and the product relations
\eqref{eq:pure-prod1}--\eqref{eq:pure-prod3}, one has
\[
a\neq 0,\quad b\neq 0,\quad c\neq 0,\quad d\neq 0,\quad e\neq 0,\quad f\neq 0.
\]
\end{lemma}

\begin{proof}
We use only: (i) $p_k=|c_k|^2\ge 0$, (ii) Lemma~\ref{lem:pure-diag-solve}, and (iii) Lemmas~\ref{lem:pure-prod1}--\ref{lem:pure-prod3}.

\emph{(1) $b\neq 0$.}
If $b=0$ then $p_9=0$, so $V=0$.
Then \eqref{eq:pure-diag-param} gives $p_8=\frac14$, so $d\neq 0$.
From \eqref{eq:pure-prod1}, $\overline{d}c=-\overline{b}a=0$, so $c=0$ hence $U=p_7=0$.
Then \eqref{eq:pure-diag-param} gives $p_4=U+V-\frac{3}{16}=-\frac{3}{16}<0$, impossible.

\emph{(2) $e\neq 0$.}
If $e=0$ then $p_4=0$, hence $U+V=\frac{3}{16}$ and $p_5=\frac{7}{16}-U-V=\frac14$, so $f\neq 0$.
From \eqref{eq:pure-prod2}, $\overline{d}f=-\overline{b}e=0$, so $d=0$ hence $p_8=0$.
But $p_8=\frac14-V$ gives $V=\frac14$, then $U=\frac{3}{16}-V=-\frac{1}{16}<0$, impossible.

\emph{(3) $c\neq 0$.}
If $c=0$ then $U=0$. From \eqref{eq:pure-prod1}, $\overline{b}a=0$ and since $b\neq 0$ by (1), we get $a=0$.
Then $p_6=|a|^2=0$, but $p_6=\frac14-U=\frac14$, contradiction.

\emph{(4) $d\neq 0$.}
If $d=0$ then $p_8=0$ so $V=\frac14$.
Then \eqref{eq:pure-prod2} gives $\overline{b}e=0$, hence $e=0$ since $b\neq 0$, contradicting (2).

\emph{(5) $a\neq 0$.}
If $a=0$ then $p_6=0$ so $U=\frac14$ and thus $c\neq 0$ (since $p_7=U$).
From \eqref{eq:pure-prod3}, $\overline{c}e=-\overline{a}f=0$, hence $e=0$, contradicting (2).

\emph{(6) $f\neq 0$.}
If $f=0$ then $p_5=0$, so $U+V=\frac{7}{16}$ and $p_4=\frac14$, hence $e\neq 0$.
From \eqref{eq:pure-prod2}, $\overline{b}e=-\overline{d}f=0$, forcing $e=0$ (since $b\neq 0$), contradiction.
\end{proof}

%------------------------------------------------------------
\paragraph{Step 4: a global phase normalization makes $a,b,c,d,e,f$ real.}

\begin{lemma}[Global phase reduction]
\label{lem:pure-phase}
Multiplying all coefficients $(c_0,\dots,c_9)$ by a fixed scalar $\omega\in\mathbb C$ with $|\omega|=1$
preserves \eqref{eq:pure-eq-system}. Hence, we may assume (without loss of generality) that
\[
b\in\mathbb R_{>0}.
\]
Under this assumption, the product relations \eqref{eq:pure-prod1}--\eqref{eq:pure-prod3} and Lemma~\ref{lem:pure-nonzero}
imply
\[
a,b,c,d,e,f\in\mathbb R.
\]
\end{lemma}

\begin{proof}
If $c_k'=\omega c_k$ with $|\omega|=1$, then $|c_k'|^2=|c_k|^2$, so the normalization and diagonal equations are preserved.
Also $\overline{c_i'}c_j'=\overline{\omega c_i}\,(\omega c_j)=\overline\omega\,\omega\,\overline{c_i}c_j=\overline{c_i}c_j$,
so every off-diagonal equation of the form $\langle \psi_0,(\text{signed permutation})\psi_1\rangle=0$ is preserved.
Thus we may choose $\omega:=\overline{b}/|b|$ to make $b$ real and positive.

Assume $b\in\mathbb R_{>0}$. From \eqref{eq:pure-prod1}, $\overline{b}a=ba\in\mathbb R$ and $b\neq 0$, hence $a\in\mathbb R$.
From \eqref{eq:pure-prod2}, $\overline{b}e=be\in\mathbb R$ and $b\neq 0$, hence $e\in\mathbb R$.
From \eqref{eq:pure-prod3}, $\overline{c}e\in\mathbb R$ and $e\neq 0$ (Lemma~\ref{lem:pure-nonzero}), hence $\overline{c}\in\mathbb R$ so $c\in\mathbb R$.
Then \eqref{eq:pure-prod1} gives $\overline{d}c\in\mathbb R$ and $c\neq 0$, hence $d\in\mathbb R$.
Finally \eqref{eq:pure-prod2} gives $\overline{d}f\in\mathbb R$ and $d\neq 0$, hence $f\in\mathbb R$.
\end{proof}

%------------------------------------------------------------
\paragraph{Step 5: the real relations force a contradiction.}

\begin{lemma}[Final contradiction]
\label{lem:pure-final}
Under the conclusions of Lemma~\ref{lem:pure-phase}, the relations \eqref{eq:pure-prod1}--\eqref{eq:pure-prod3} imply
\[
af=0.
\]
This contradicts Lemma~\ref{lem:pure-nonzero}.
\end{lemma}

\begin{proof}
In the phase where $a,b,c,d,e,f\in\mathbb R$, the relations
$\overline{b}a=-\overline{d}c$, $\overline{b}e=-\overline{d}f$, $\overline{c}e=-\overline{a}f$
become
\begin{equation}
ba=-dc,\qquad be=-df,\qquad ce=-af.
\label{eq:pure-real-rel}
\end{equation}
Multiply $ba=-dc$ by $f$ and $be=-df$ by $c$:
\[
baf=-dcf,\qquad bce=-dfc=-dcf.
\]
Thus $baf=bce$. Since $b\neq 0$ (Lemma~\ref{lem:pure-nonzero}), we get $af=ce$.
But the third equation in \eqref{eq:pure-real-rel} says $ce=-af$. Hence $af=-af$, so $af=0$.
\end{proof}

\paragraph{Conclusion.}
By Lemma~\ref{lem:pure-final} we obtain a contradiction, so the assumed solution to \eqref{eq:pure-eq-system} cannot exist.
This proves Theorem~\ref{thm:pure-nogo-subsystem}.
\end{proof}

\begin{remark}[What is (and is not) used]
The proof is purely algebraic and uses only:
\begin{itemize}
\item the explicit finite set $S_0\subset\{0,1\}^7$,
\item the definitions \eqref{eq:pure-def-psi01} of $\psi_0,\psi_1$,
\item the signed-permutation operators $D_i$, $F_i$, and $D_j\circ F_i$,
\item the $17$ bilinear equations in \eqref{eq:pure-eq-system}.
\end{itemize}
No other structure (subset-sum, codes, or ``quantum'' interpretation) is required for the contradiction.
\end{remark}

%============================================================
\subsection{A completely real-coordinate no--go certificate for the
\texorpdfstring{$(1+7+9)$}{(1+7+9)}-equation subsystem
(norm + 7 diagonal + 9 off-diagonal)}
\label{subsec:real-coordinate-nogo-1112226-subsystem}

This subsection rewrites the subsystem \emph{purely over $\mathbb R$} by expanding the
$10$ complex amplitudes into $20$ real variables and then showing that the resulting
finite system of \emph{real quadratic} equations is unsatisfiable.

%------------------------------------------------------------
\paragraph{Fixed data (support set).}
Let $S_0=\{s_0,\dots,s_9\}\subset\{0,1\}^7$ be the explicit list
\[
\begin{aligned}
&s_0=0000000,\ s_1=0000011,\ s_2=0000101,\ s_3=0001001,\ s_4=0110001,\\
&s_5=0111110,\ s_6=1010001,\ s_7=1011110,\ s_8=1100001,\ s_9=1101110,
\end{aligned}
\]
and let $S_7:=\overline{S_0}=\{\bar s:\ s\in S_0\}$ where $\bar s$ is bitwise complement.

%------------------------------------------------------------
\paragraph{Real unknowns and their complex packaging.}
Introduce $20$ real variables
\[
(v_1,\dots,v_{20})\in\mathbb R^{20}.
\]
Define $10$ complex amplitudes $c_0,\dots,c_9\in\mathbb C$ by
\begin{align}
c_0&:=v_1+i v_2, &
c_1&:=v_3+i v_4, &
c_2&:=v_5+i v_6, &
c_3&:=v_7+i v_8, \label{eq:real-coords-ck-1}\\
c_4&:=v_9+i v_{10}, &
c_5&:=v_{11}+i v_{12}, &
c_6&:=v_{13}+i v_{14}, &
c_7&:=v_{15}+i v_{16}, \label{eq:real-coords-ck-2}\\
c_8&:=v_{17}+i v_{18}, &
c_9&:=v_{19}+i v_{20}. &&
\label{eq:real-coords-ck-3}
\end{align}
Define the squared magnitudes (nonnegative reals)
\begin{equation}
p_k:=|c_k|^2=v_{2k+1}^2+v_{2k+2}^2\qquad (k=0,\dots,9).
\label{eq:real-coords-pk}
\end{equation}
Equivalently,
\[
\begin{aligned}
&p_0=v_1^2+v_2^2,\quad p_1=v_3^2+v_4^2,\quad p_2=v_5^2+v_6^2,\quad p_3=v_7^2+v_8^2,\\
&p_4=v_9^2+v_{10}^2,\quad p_5=v_{11}^2+v_{12}^2,\quad p_6=v_{13}^2+v_{14}^2,\quad p_7=v_{15}^2+v_{16}^2,\\
&p_8=v_{17}^2+v_{18}^2,\quad p_9=v_{19}^2+v_{20}^2.
\end{aligned}
\]

It will be convenient to name the six amplitudes that appear in the off-diagonal subsystem:
\begin{equation}
a:=c_6=v_{13}+i v_{14},\quad
b:=c_9=v_{19}+i v_{20},\quad
c:=c_7=v_{15}+i v_{16},\quad
d:=c_8=v_{17}+i v_{18},\quad
e:=c_4=v_{9}+i v_{10},\quad
f:=c_5=v_{11}+i v_{12}.
\label{eq:real-coords-abcdef}
\end{equation}

%------------------------------------------------------------
\paragraph{Counting variables and equations.}
The unknowns are the $20$ real variables $v_1,\dots,v_{20}$.
The subsystem consists of $1$ normalization equation, $7$ diagonal equations, and $9$ off-diagonal equations,
for a total of
\[
\boxed{\text{$20$ real variables and $17$ real quadratic equations.}}
\]

%------------------------------------------------------------
\paragraph{The explicit real quadratic system (17 equations).}
We now list the full system as explicit real polynomial equations in the $v_i$.

\smallskip
\noindent\textbf{(N) Normalization (1 equation).}
\begin{equation}
\boxed{
(v_1^2+v_2^2)+(v_3^2+v_4^2)+\cdots+(v_{19}^2+v_{20}^2)=1.
}
\tag{N}\label{eq:real:N}
\end{equation}

\smallskip
\noindent\textbf{(D1)--(D7) Seven diagonal equations (7 equations).}
Using the explicit $S_0$ above, the diagonal constraints are exactly:
\begin{align}
&(v_{13}^2+v_{14}^2)+(v_{15}^2+v_{16}^2)+(v_{17}^2+v_{18}^2)+(v_{19}^2+v_{20}^2)=\tfrac12,
\tag{D1}\label{eq:real:D1}\\
&(v_{9}^2+v_{10}^2)+(v_{11}^2+v_{12}^2)+(v_{17}^2+v_{18}^2)+(v_{19}^2+v_{20}^2)=\tfrac12,
\tag{D2}\label{eq:real:D2}\\
&(v_{9}^2+v_{10}^2)+(v_{11}^2+v_{12}^2)+(v_{13}^2+v_{14}^2)+(v_{15}^2+v_{16}^2)=\tfrac12,
\tag{D3}\label{eq:real:D3}\\
&(v_{7}^2+v_{8}^2)+(v_{11}^2+v_{12}^2)+(v_{15}^2+v_{16}^2)+(v_{19}^2+v_{20}^2)=\tfrac12,
\tag{D4}\label{eq:real:D4}\\
&(v_{5}^2+v_{6}^2)+(v_{11}^2+v_{12}^2)+(v_{15}^2+v_{16}^2)+(v_{19}^2+v_{20}^2)=\tfrac12,
\tag{D5}\label{eq:real:D5}\\
&(v_{3}^2+v_{4}^2)+(v_{11}^2+v_{12}^2)+(v_{15}^2+v_{16}^2)+(v_{19}^2+v_{20}^2)=\tfrac12,
\tag{D6}\label{eq:real:D6}\\
&(v_{3}^2+v_{4}^2)+(v_{5}^2+v_{6}^2)+(v_{7}^2+v_{8}^2)
+(v_{9}^2+v_{10}^2)+(v_{13}^2+v_{14}^2)+(v_{17}^2+v_{18}^2)=\tfrac12.
\tag{D7}\label{eq:real:D7}
\end{align}

\smallskip
\noindent\textbf{(X1--X3) Nine off-diagonal equations (9 equations).}
These are the real-coordinate expansions of the chosen off-diagonal constraints:
\begin{align}
&v_{13}v_{19}+v_{14}v_{20}+v_{15}v_{17}+v_{16}v_{18}=0,
\tag{X1-0}\label{eq:real:X1-0}\\
&v_{13}v_{20}-v_{14}v_{19}+v_{15}v_{18}-v_{16}v_{17}=0,
\tag{X1-2}\label{eq:real:X1-2}\\
&v_{13}v_{20}-v_{14}v_{19}-v_{15}v_{18}+v_{16}v_{17}=0,
\tag{X1-4}\label{eq:real:X1-4}\\[1mm]
&v_{11}v_{17}+v_{12}v_{18}+v_{19}v_{9}+v_{20}v_{10}=0,
\tag{X2-0}\label{eq:real:X2-0}\\
&v_{11}v_{18}-v_{12}v_{17}-v_{19}v_{10}+v_{20}v_{9}=0,
\tag{X2-1}\label{eq:real:X2-1}\\
&-v_{11}v_{18}+v_{12}v_{17}-v_{19}v_{10}+v_{20}v_{9}=0,
\tag{X2-4}\label{eq:real:X2-4}\\[1mm]
&v_{11}v_{13}+v_{12}v_{14}+v_{15}v_{9}+v_{16}v_{10}=0,
\tag{X3-0}\label{eq:real:X3-0}\\
&v_{11}v_{14}-v_{12}v_{13}-v_{15}v_{10}+v_{16}v_{9}=0,
\tag{X3-1}\label{eq:real:X3-1}\\
&-v_{11}v_{14}+v_{12}v_{13}-v_{15}v_{10}+v_{16}v_{9}=0.
\tag{X3-4}\label{eq:real:X3-4}
\end{align}

%------------------------------------------------------------
\paragraph{No-go statement.}
\begin{proposition}[Unsatisfiability over $\mathbb R$]
\label{prop:real-unsat-subsystem}
There is no real vector $(v_1,\dots,v_{20})\in\mathbb R^{20}$ satisfying simultaneously
\eqref{eq:real:N}, \eqref{eq:real:D1}--\eqref{eq:real:D7}, and \eqref{eq:real:X1-0}--\eqref{eq:real:X3-4}.
\end{proposition}

\begin{proof}
Assume for contradiction that $(v_1,\dots,v_{20})\in\mathbb R^{20}$ satisfies all $17$ equations.
Define $p_k$ by \eqref{eq:real-coords-pk}, so $p_k\ge 0$ for all $k$.

%------------------------------------------------------------
\paragraph{Step 1: Solve the diagonal subsystem as a linear system in the $p_k$.}
We rewrite \eqref{eq:real:N} and \eqref{eq:real:D1}--\eqref{eq:real:D7} using the shorthand $p_k=|c_k|^2$.

\smallskip
\noindent\emph{(i) First deduce $p_1=p_2=p_3$.}
Subtract \eqref{eq:real:D5} from \eqref{eq:real:D4} to get $p_3-p_2=0$, so $p_3=p_2$.
Subtract \eqref{eq:real:D6} from \eqref{eq:real:D5} to get $p_2-p_1=0$, so $p_2=p_1$.
Hence
\begin{equation}
p_1=p_2=p_3=:t.
\label{eq:real-step1-t}
\end{equation}

\smallskip
\noindent\emph{(ii) Deduce $p_6+p_7=p_8+p_9=\tfrac14$.}
Subtract \eqref{eq:real:D3} from \eqref{eq:real:D2} to obtain
\[
(p_4+p_5+p_8+p_9)-(p_4+p_5+p_6+p_7)=0
\quad\Longrightarrow\quad
p_8+p_9=p_6+p_7.
\]
Together with \eqref{eq:real:D1}, which says $p_6+p_7+p_8+p_9=\tfrac12$, we conclude
\begin{equation}
p_6+p_7=\tfrac14,
\qquad
p_8+p_9=\tfrac14.
\label{eq:real-step1-halves}
\end{equation}

\smallskip
\noindent\emph{(iii) Deduce $p_4+p_5=\tfrac14$.}
Plug \eqref{eq:real-step1-halves} into \eqref{eq:real:D2}:
\[
p_4+p_5+(p_8+p_9)=\tfrac12
\quad\Longrightarrow\quad
p_4+p_5+\tfrac14=\tfrac12,
\]
so
\begin{equation}
p_4+p_5=\tfrac14.
\label{eq:real-step1-p4p5}
\end{equation}

\smallskip
\noindent\emph{(iv) Use normalization to deduce $p_0+3t=\tfrac14$.}
From \eqref{eq:real:N} and \eqref{eq:real-step1-t}--\eqref{eq:real-step1-p4p5} we have
\[
1
=p_0+(p_1+p_2+p_3)+(p_4+p_5)+(p_6+p_7)+(p_8+p_9)
=p_0+3t+\tfrac14+\tfrac14+\tfrac14,
\]
hence
\begin{equation}
p_0+3t=\tfrac14.
\label{eq:real-step1-p0t}
\end{equation}

\smallskip
\noindent\emph{(v) Determine $t$ (and hence $p_0$).}
Equation \eqref{eq:real:D6} is $t+p_5+p_7+p_9=\tfrac12$.
Using \eqref{eq:real-step1-halves} and \eqref{eq:real-step1-p4p5}, write
\[
p_5=\tfrac14-p_4,\qquad p_7=\tfrac14-p_6,\qquad p_9=\tfrac14-p_8.
\]
Then \eqref{eq:real:D6} becomes
\[
t+\Bigl(\tfrac14-p_4\Bigr)+\Bigl(\tfrac14-p_6\Bigr)+\Bigl(\tfrac14-p_8\Bigr)=\tfrac12
\quad\Longrightarrow\quad
t=(p_4+p_6+p_8)-\tfrac14.
\]
Meanwhile \eqref{eq:real:D7} reads
\[
p_1+p_2+p_3+p_4+p_6+p_8=\tfrac12
\quad\Longleftrightarrow\quad
3t+(p_4+p_6+p_8)=\tfrac12.
\]
Substitute $t=(p_4+p_6+p_8)-\tfrac14$ into the latter:
\[
3\Bigl((p_4+p_6+p_8)-\tfrac14\Bigr)+(p_4+p_6+p_8)=\tfrac12
\quad\Longrightarrow\quad
4(p_4+p_6+p_8)=\tfrac54.
\]
So $p_4+p_6+p_8=\tfrac{5}{16}$ and therefore $t=\tfrac{5}{16}-\tfrac14=\tfrac{1}{16}$.
Thus
\begin{equation}
p_1=p_2=p_3=\tfrac{1}{16}.
\label{eq:real-step1-p123}
\end{equation}
Finally, \eqref{eq:real-step1-p0t} gives $p_0=\tfrac14-3\cdot \tfrac{1}{16}=\tfrac{1}{16}$:
\begin{equation}
p_0=p_1=p_2=p_3=\tfrac{1}{16}.
\label{eq:real-step1-fourfixed}
\end{equation}

\smallskip
\noindent\emph{(vi) Parametrize the remaining $p_k$.}
Set
\begin{equation}
U:=p_7=v_{15}^2+v_{16}^2,
\qquad
V:=p_9=v_{19}^2+v_{20}^2.
\label{eq:real-step1-UV}
\end{equation}
Then \eqref{eq:real-step1-halves} implies
\[
p_6=\tfrac14-U,\qquad p_8=\tfrac14-V.
\]
Using \eqref{eq:real:D6} together with \eqref{eq:real-step1-p123}, we get
\[
\tfrac{1}{16}+p_5+U+V=\tfrac12
\quad\Longrightarrow\quad
p_5=\tfrac{7}{16}-U-V.
\]
Finally \eqref{eq:real-step1-p4p5} yields
\[
p_4=\tfrac14-p_5=U+V-\tfrac{3}{16}.
\]
Summarizing,
\begin{equation}
\boxed{
\begin{aligned}
&p_0=p_1=p_2=p_3=\tfrac{1}{16},\qquad p_7=U,\qquad p_9=V,\\
&p_6=\tfrac14-U,\qquad p_8=\tfrac14-V,\qquad
p_4=U+V-\tfrac{3}{16},\qquad p_5=\tfrac{7}{16}-U-V.
\end{aligned}}
\label{eq:real-step1-param}
\end{equation}
Since each $p_k$ is a sum of squares, $p_k\ge 0$ forces
\begin{equation}
0\le U\le \tfrac14,\qquad 0\le V\le \tfrac14,\qquad \tfrac{3}{16}\le U+V\le \tfrac{7}{16}.
\label{eq:real-step1-constraints}
\end{equation}

%------------------------------------------------------------
\paragraph{Step 2: The nine off-diagonal equations force six ``dot/cross'' relations.}
From the paired equations in each $(X1)$, $(X2)$, $(X3)$ block, we extract the following consequences.

\smallskip
\noindent\emph{(a) From \eqref{eq:real:X1-2} and \eqref{eq:real:X1-4}.}
Adding and subtracting these two equations gives:
\begin{align}
v_{13}v_{20}-v_{14}v_{19}&=0,
\tag{X1-cross-$ba$}\label{eq:real:X1-cross-ba}\\
v_{15}v_{18}-v_{16}v_{17}&=0.
\tag{X1-cross-$cd$}\label{eq:real:X1-cross-cd}
\end{align}
We keep \eqref{eq:real:X1-0} as the corresponding ``dot'' equation.

\smallskip
\noindent\emph{(b) From \eqref{eq:real:X2-1} and \eqref{eq:real:X2-4}.}
Adding and subtracting gives:
\begin{align}
v_{20}v_{9}-v_{19}v_{10}&=0,
\tag{X2-cross-$be$}\label{eq:real:X2-cross-be}\\
v_{11}v_{18}-v_{12}v_{17}&=0.
\tag{X2-cross-$df$}\label{eq:real:X2-cross-df}
\end{align}
We keep \eqref{eq:real:X2-0} as the corresponding ``dot'' equation.

\smallskip
\noindent\emph{(c) From \eqref{eq:real:X3-1} and \eqref{eq:real:X3-4}.}
Adding and subtracting gives:
\begin{align}
v_{16}v_{9}-v_{15}v_{10}&=0,
\tag{X3-cross-$ce$}\label{eq:real:X3-cross-ce}\\
v_{11}v_{14}-v_{12}v_{13}&=0.
\tag{X3-cross-$af$}\label{eq:real:X3-cross-af}
\end{align}
We keep \eqref{eq:real:X3-0} as the corresponding ``dot'' equation.

%------------------------------------------------------------
\paragraph{Step 3: The diagonal parametrization plus the off-diagonal relations force $a,b,c,d,e,f\neq 0$.}
Recall $a,b,c,d,e,f$ from \eqref{eq:real-coords-abcdef}.  We show that none of the corresponding pairs
\[
(v_{13},v_{14}),\ (v_{19},v_{20}),\ (v_{15},v_{16}),\ (v_{17},v_{18}),\ (v_9,v_{10}),\ (v_{11},v_{12})
\]
can be $(0,0)$.

\smallskip
\noindent\emph{(1) $b\neq 0$.}
Assume $b=0$, i.e.\ $v_{19}=v_{20}=0$. Then $V=p_9=0$.
By \eqref{eq:real-step1-param}, $p_8=\tfrac14-V=\tfrac14$, so $(v_{17},v_{18})\neq (0,0)$.
With $v_{19}=v_{20}=0$, equation \eqref{eq:real:X1-0} becomes
\[
v_{15}v_{17}+v_{16}v_{18}=0
\]
and \eqref{eq:real:X1-cross-cd} becomes
\[
v_{15}v_{18}-v_{16}v_{17}=0.
\]
Thus the vector $(v_{15},v_{16})$ is simultaneously orthogonal and parallel to $(v_{17},v_{18})\neq(0,0)$, forcing
$(v_{15},v_{16})=(0,0)$, i.e.\ $c=0$ and $U=p_7=0$.
Then \eqref{eq:real-step1-param} gives
\[
p_4=U+V-\tfrac{3}{16}=-\tfrac{3}{16}<0,
\]
contradicting $p_4=v_9^2+v_{10}^2\ge 0$. Hence $b\neq 0$.

\smallskip
\noindent\emph{(2) $e\neq 0$.}
Assume $e=0$, i.e.\ $v_9=v_{10}=0$ so $p_4=0$.
Then \eqref{eq:real-step1-param} gives $U+V=\tfrac{3}{16}$ and hence
\[
p_5=\tfrac{7}{16}-U-V=\tfrac14,
\]
so $(v_{11},v_{12})\neq (0,0)$ (i.e.\ $f\neq 0$).
With $v_9=v_{10}=0$, equation \eqref{eq:real:X2-0} becomes
\[
v_{11}v_{17}+v_{12}v_{18}=0
\]
and \eqref{eq:real:X2-cross-df} becomes
\[
v_{11}v_{18}-v_{12}v_{17}=0.
\]
Thus $(v_{17},v_{18})$ is simultaneously orthogonal and parallel to the nonzero vector $(v_{11},v_{12})$, forcing
$(v_{17},v_{18})=(0,0)$, i.e.\ $d=0$ and $p_8=0$.
But $p_8=\tfrac14-V$ then gives $V=\tfrac14$, so $U=\tfrac{3}{16}-\tfrac14=-\tfrac{1}{16}<0$,
contradicting $U=p_7\ge 0$. Hence $e\neq 0$.

\smallskip
\noindent\emph{(3) $c\neq 0$.}
Assume $c=0$, i.e.\ $v_{15}=v_{16}=0$ so $U=p_7=0$.
Then \eqref{eq:real-step1-param} gives $p_6=\tfrac14-U=\tfrac14$, so $(v_{13},v_{14})\neq (0,0)$ (i.e.\ $a\neq 0$).
But with $v_{15}=v_{16}=0$, equations \eqref{eq:real:X1-0} and \eqref{eq:real:X1-cross-ba} become
\[
v_{13}v_{19}+v_{14}v_{20}=0,\qquad
v_{13}v_{20}-v_{14}v_{19}=0,
\]
so $(v_{13},v_{14})$ is orthogonal and parallel to $(v_{19},v_{20})$.
Since $b\neq 0$ from (1), $(v_{19},v_{20})\neq (0,0)$ and hence $(v_{13},v_{14})=(0,0)$,
contradicting $p_6=\tfrac14$. Hence $c\neq 0$.

\smallskip
\noindent\emph{(4) $d\neq 0$.}
Assume $d=0$, i.e.\ $v_{17}=v_{18}=0$ so $p_8=0$ and $V=\tfrac14$ by \eqref{eq:real-step1-param}.
With $v_{17}=v_{18}=0$, equations \eqref{eq:real:X2-0} and \eqref{eq:real:X2-cross-be} become
\[
v_{19}v_{9}+v_{20}v_{10}=0,\qquad
v_{20}v_{9}-v_{19}v_{10}=0,
\]
so $(v_9,v_{10})$ is orthogonal and parallel to $(v_{19},v_{20})$.
Since $b\neq 0$, this forces $(v_9,v_{10})=(0,0)$, i.e.\ $e=0$, contradicting (2).
Hence $d\neq 0$.

\smallskip
\noindent\emph{(5) $a\neq 0$.}
Assume $a=0$, i.e.\ $v_{13}=v_{14}=0$ so $p_6=0$.
Then \eqref{eq:real-step1-param} gives $U=\tfrac14$, so $c\neq 0$ (since $p_7=U=\tfrac14$).
With $v_{13}=v_{14}=0$, equations \eqref{eq:real:X3-0} and \eqref{eq:real:X3-cross-ce} become
\[
v_{15}v_{9}+v_{16}v_{10}=0,\qquad
v_{16}v_{9}-v_{15}v_{10}=0,
\]
so $(v_9,v_{10})$ is orthogonal and parallel to $(v_{15},v_{16})$.
Since $c\neq 0$, this forces $(v_9,v_{10})=(0,0)$, i.e.\ $e=0$, contradicting (2).
Hence $a\neq 0$.

\smallskip
\noindent\emph{(6) $f\neq 0$.}
Assume $f=0$, i.e.\ $v_{11}=v_{12}=0$ so $p_5=0$.
Then \eqref{eq:real-step1-param} gives $U+V=\tfrac{7}{16}$ and hence $p_4=\tfrac14$,
so $e\neq 0$ (i.e.\ $(v_9,v_{10})\neq(0,0)$).
With $v_{11}=v_{12}=0$, equations \eqref{eq:real:X2-0} and \eqref{eq:real:X2-cross-be} again become
\[
v_{19}v_{9}+v_{20}v_{10}=0,\qquad
v_{20}v_{9}-v_{19}v_{10}=0,
\]
forcing $(v_9,v_{10})=(0,0)$ since $b\neq 0$, contradiction.
Hence $f\neq 0$.

Therefore,
\begin{equation}
a\neq 0,\quad b\neq 0,\quad c\neq 0,\quad d\neq 0,\quad e\neq 0,\quad f\neq 0.
\label{eq:real-step3-nonzero}
\end{equation}

%------------------------------------------------------------
\paragraph{Step 4: Global phase symmetry lets us impose $b\in\mathbb R_{>0}$ and forces $a,c,d,e,f\in\mathbb R$.}
The system \eqref{eq:real:N}--\eqref{eq:real:X3-4} is invariant under a \emph{common rotation} of each pair
$(v_{2k+1},v_{2k+2})$ by the same angle $\theta$:
\begin{equation}
(v_{2k+1},v_{2k+2})\longmapsto
(v_{2k+1}',v_{2k+2}')
:=
\bigl(v_{2k+1}\cos\theta - v_{2k+2}\sin\theta,\ \ v_{2k+1}\sin\theta + v_{2k+2}\cos\theta\bigr).
\label{eq:real-global-rotation}
\end{equation}
Indeed, \eqref{eq:real-global-rotation} is exactly the change induced by multiplying every complex $c_k$
by the same unit scalar $e^{i\theta}$.
It preserves each $p_k=v_{2k+1}^2+v_{2k+2}^2$ and preserves all bilinear combinations
appearing in \eqref{eq:real:X1-0}--\eqref{eq:real:X3-4}.
Since $b\neq 0$ by \eqref{eq:real-step3-nonzero}, choose $\theta$ so that
\begin{equation}
b=v_{19}+i v_{20}\ \ \text{is real and positive, i.e.}\ \ v_{20}=0,\ v_{19}>0.
\label{eq:real-b-real-positive}
\end{equation}

Now use the ``cross'' equations from Step 2:
\begin{itemize}
\item From \eqref{eq:real:X1-cross-ba} and \eqref{eq:real-b-real-positive}:
$0=v_{13}v_{20}-v_{14}v_{19}=-v_{14}v_{19}$, hence $v_{14}=0$ and $a=v_{13}\in\mathbb R$.
\item From \eqref{eq:real:X2-cross-be} and \eqref{eq:real-b-real-positive}:
$0=v_{20}v_{9}-v_{19}v_{10}=-v_{19}v_{10}$, hence $v_{10}=0$ and $e=v_{9}\in\mathbb R$.
Moreover $e\neq 0$ implies $v_9\neq 0$.
\item From \eqref{eq:real:X3-cross-ce} with $v_{10}=0$ and $v_9\neq 0$:
$0=v_{16}v_{9}-v_{15}v_{10}=v_{16}v_{9}$, hence $v_{16}=0$ and $c=v_{15}\in\mathbb R$.
Moreover $c\neq 0$ implies $v_{15}\neq 0$.
\item From \eqref{eq:real:X1-cross-cd} with $v_{16}=0$ and $v_{15}\neq 0$:
$0=v_{15}v_{18}-v_{16}v_{17}=v_{15}v_{18}$, hence $v_{18}=0$ and $d=v_{17}\in\mathbb R$.
Moreover $d\neq 0$ implies $v_{17}\neq 0$.
\item From \eqref{eq:real:X2-cross-df} with $v_{18}=0$ and $v_{17}\neq 0$:
$0=v_{11}v_{18}-v_{12}v_{17}=-v_{12}v_{17}$, hence $v_{12}=0$ and $f=v_{11}\in\mathbb R$.
\end{itemize}
Therefore, after the global rotation we may assume
\begin{equation}
a=v_{13},\quad b=v_{19},\quad c=v_{15},\quad d=v_{17},\quad e=v_9,\quad f=v_{11}
\quad\in\mathbb R\setminus\{0\}.
\label{eq:real-step4-all-real}
\end{equation}

%------------------------------------------------------------
\paragraph{Step 5: The remaining dot equations force $af=0$, contradiction.}
Under \eqref{eq:real-step4-all-real}, the dot equations simplify:

\smallskip
\noindent From \eqref{eq:real:X1-0} with $v_{14}=v_{20}=v_{16}=v_{18}=0$ we get
\[
v_{13}v_{19}+v_{15}v_{17}=0
\quad\Longleftrightarrow\quad
ba=-dc.
\]
From \eqref{eq:real:X2-0} with $v_{12}=v_{18}=v_{20}=v_{10}=0$ we get
\[
v_{11}v_{17}+v_{19}v_{9}=0
\quad\Longleftrightarrow\quad
be=-df.
\]
From \eqref{eq:real:X3-0} with $v_{12}=v_{14}=v_{16}=v_{10}=0$ we get
\[
v_{11}v_{13}+v_{15}v_{9}=0
\quad\Longleftrightarrow\quad
ce=-af.
\]
Thus
\begin{equation}
ba=-dc,\qquad be=-df,\qquad ce=-af,
\qquad\text{with }a,b,c,d,e,f\in\mathbb R\setminus\{0\}.
\label{eq:real-step5-relations}
\end{equation}

Multiply $ba=-dc$ by $f$ and $be=-df$ by $c$:
\[
baf=-dcf,\qquad bce=-dfc=-dcf,
\]
hence $baf=bce$. Since $b\neq 0$, cancel $b$ to obtain $af=ce$.
But the third relation in \eqref{eq:real-step5-relations} says $ce=-af$.
Therefore $af=-af$, so
\[
af=0,
\]
contradicting $a\neq 0$ and $f\neq 0$ from \eqref{eq:real-step4-all-real}.

This contradiction shows the assumed real solution $(v_1,\dots,v_{20})$ cannot exist.
\end{proof}


\section{Negative evidence for \texorpdfstring{$\vec a=(0,1,1,2,3,3,5)$}{a=(0,1,1,2,3,3,5)} in the BD16 single-syndrome setting}

We consider the BD16 setting with \(m=8\) and angle vector
\[
\vec a=(0,1,1,2,3,3,5)\in\mathbb{Z}_8^{\,7}.
\]
As in the rest of these notes, we interpret the task as searching for a distance-\(3\) \(((7,2,3))\) code satisfying:
(i) \(\ket{0_L}\) supported only on the residue class
\[
S_0(\vec a)=\{\,s\in\{0,1\}^7:\ \vec a\cdot s \equiv 0 \pmod 8\,\},
\]
(ii) \(\ket{1_L}\) supported only on the complementary class \(S_7(\vec a)\) (typically via self-complementarity
\(\ket{1_L}=X^{\otimes 7}\ket{0_L}\), since \(\sum_j a_j\equiv 7\pmod 8\)), and
(iii) full Knill--Laflamme conditions for all Pauli errors of weight \(\le 2\).

\subsection{Subset-sum support sets}
With dot product
\(
\vec a\cdot s = 0\cdot s_1 + 1\cdot s_2 + 1\cdot s_3 + 2\cdot s_4 + 3\cdot s_5 + 3\cdot s_6 + 5\cdot s_7
\pmod 8,
\)
the residue-\(0\) class is
\[
\begin{aligned}
S_0(\vec a)=\{&
0000000,\ 0000011,\ 0000101,\ 0001110,\ 0011001,\ 0101001,\ 0110110,\\
&1000000,\ 1000011,\ 1000101,\ 1001110,\ 1011001,\ 1101001,\ 1110110
\}.
\end{aligned}
\]
Since \(\sum_j a_j=15\equiv 7\pmod 8\), the bitwise-complement map sends \(S_0(\vec a)\) to \(S_7(\vec a)\), namely
\[
\begin{aligned}
S_7(\vec a)=\{\overline{s}:s\in S_0(\vec a)\}=\{&
1111111,\ 1111100,\ 1111010,\ 1110001,\ 1100110,\ 1010110,\ 1001001,\\
&0111111,\ 0111100,\ 0111010,\ 0110001,\ 0100110,\ 0010110,\ 0001001
\}.
\end{aligned}
\]
A notable structural feature is \(a_1=0\), hence \(S_0(\vec a)\) comes in \(7\) pairs \((0u,1u)\) sharing the same last
six bits \(u\).

\subsection{LP-stage feasibility (Z-balance) but failure at full KL}
Imposing the natural pair symmetry suggested by \(a_1=0\),
\[
|c_{0u}|^2=|c_{1u}|^2\qquad\text{for each paired support }(0u,1u),
\]
the \(Z\)-balance LP constraints (in particular \(\langle Z_i\rangle_{0_L}=0\) for all \(i\)) force a \emph{unique} distribution
over the \(7\) pairs:
\[
p_u := \Pr(\text{pair }u)=\left(\frac1{16},\frac1{16},\frac1{16},\frac2{16},\frac3{16},\frac3{16},\frac5{16}\right)
\]
for
\[
u\in\{000000,\ 000011,\ 000101,\ 001110,\ 011001,\ 101001,\ 110110\}.
\]
Equivalently, each individual basis state \(0u\) or \(1u\) would have probability \(p_u/2\), yielding magnitudes in the set
\[
|c|\in\left\{\frac{\sqrt2}{8},\ \frac14,\ \frac{\sqrt6}{8},\ \frac{\sqrt{10}}{8}\right\}.
\]
This provides a very ``analytic-looking'' LP-stage solution, but extensive searches did not find any choice of phases (including
broad continuous searches) that upgrades this LP solution to a full Knill--Laflamme \(((7,2,3))\) code.

\subsection{Search regimes attempted (and what failed)}
\paragraph{Discrete ``paper-style'' ansatz (square-root rationals \(\times\) 4th roots).}
I searched over many LP-feasible integer-weight patterns of the form
\[
|c_s|^2=\frac{w_s}{16},\qquad w_s\in\{0,1,2,3,4\},
\]
including \(10\) distinct sparse LP-feasible patterns with support size \(8\), and then exhaustively brute-forced phases
\[
c_s=\omega_s\frac{\sqrt{w_s}}{4},\qquad \omega_s\in\{\pm 1,\pm i\}.
\]
No full-KL solutions were found for any such pattern (even after filtering on all weight-\(1\) Pauli constraints first, then
checking full weight-\(\le 2\) KL).

\paragraph{Continuous numerical solving with free complex amplitudes.}
To test whether failure was an artifact of restricting phases to \(\{\pm1,\pm i\}\), I ran Levenberg--Marquardt least-squares
searches allowing arbitrary complex coefficients (with normalization enforced internally) for:
(i) the self-complement case \(\ket{1_L}=X^{\otimes 7}\ket{0_L}\), and
(ii) the more general case where \(\ket{0_L}\) on \(S_0(\vec a)\) and \(\ket{1_L}\) on \(S_7(\vec a)\) are independent.
In both setups, repeated random restarts converged to the same nonzero minimum rather than to a near-zero residual.

A simple diagnostic at the best point found indicated that the largest off-diagonal violation occurred for a weight-\(1\)
\(X\) error (e.g.\ \(X\) on qubit \(3\)), with magnitude \(\sim 0.24\), and a large diagonal mismatch also appeared for a
weight-\(1\) \(Y\) error (e.g.\ \(Y\) on qubit \(1\)). These violations are not consistent with being ``near'' a true solution.

\subsection{Bottom line}
Under the standard BD16 interpretation (support restricted to \(S_0(\vec a)\) and \(S_7(\vec a)\), and full KL for all Pauli
errors of weight \(\le 2\)), no analytic \(((7,2,3))\) solution was found for \(\vec a=(0,1,1,2,3,3,5)\).
The combination of exhaustive discrete searches in the usual exact-coefficient ansatz and repeated continuous searches with
fully free complex amplitudes provides strong evidence that this angle vector is an SS/LP false positive: it passes the SS and
\(Z\)-LP filters but does not admit a distance-\(3\) \(((7,2,3))\) code satisfying full Knill--Laflamme constraints within these
support restrictions.

\subsection{A numerical no-solution experiment for \texorpdfstring{$\vec a=(0,1,1,2,3,3,5)$}{a=(0,1,1,2,3,3,5)} in the BD16 SS-support model}

We consider the BD16 setting with \(m=8\), residues \((b_0,b_1)=(0,7)\), and angle vector
\[
\vec a=(0,1,1,2,3,3,5)\in\mathbb{Z}_8^{\,7}.
\]
The goal is to test (numerically, with strong robustness checks) whether there exists a distance-\(3\) \(((7,2,3))\) code
satisfying the \emph{SS support constraint}
\[
\ket{0_L}\subset S_0(\vec a),\qquad \ket{1_L}\subset S_7(\vec a),
\]
and (optionally) the \emph{complementary/self-complementary} convention \(\ket{1_L}=X^{\otimes 7}\ket{0_L}\).
The experiment minimizes an explicit sum-of-squares Knill--Laflamme loss \(\mathcal{L}\); a strictly positive minimum provides
strong evidence that the full KL system is inconsistent within the searched parameter space.

\subsubsection{SS support sets}
Define the residue map
\[
\mathrm{res}(s):=\sum_{j=1}^7 a_j s_j \pmod 8,\qquad s\in\{0,1\}^7.
\]
For \(\vec a=(0,1,1,2,3,3,5)\), the residue-\(0\) class has size \(14\):
\[
\begin{aligned}
S_0(\vec a)=\{&
0000000,\ 0000011,\ 0000101,\ 0001110,\ 0011001,\ 0101001,\ 0110110,\\
&1000000,\ 1000011,\ 1000101,\ 1001110,\ 1011001,\ 1101001,\ 1110110
\}.
\end{aligned}
\]
Since \(\sum_j a_j=15\equiv 7\pmod 8\), the bitwise complement map sends \(S_0(\vec a)\) to \(S_7(\vec a)\), i.e.
\[
\overline{s}\in S_7(\vec a)\quad\text{for all } s\in S_0(\vec a),
\qquad
S_7(\vec a)=\{\overline{s}:s\in S_0(\vec a)\}.
\]
A structural feature is \(a_1=0\), hence \(S_0(\vec a)\) decomposes into \(7\) pairs \((0u,1u)\) sharing the same last six bits \(u\).

\subsubsection{Parameterizations of the codewords}
We tested two increasingly general parameter spaces.

\paragraph{Case A (complementary).}
\[
\ket{0_L}=\sum_{s\in S_0} c_s \ket{s},\qquad
\ket{1_L}=X^{\otimes 7}\ket{0_L}=\sum_{s\in S_0} c_s \ket{\overline{s}},
\qquad
\sum_{s\in S_0}|c_s|^2=1.
\]
Orthogonality \(\braket{0_L|1_L}=0\) holds automatically since \(S_0\cap S_7=\varnothing\).
This is the full continuous complementary SS-support family: \(14\) complex amplitudes modulo normalization and global phase.

\paragraph{Case B (non-complementary, still SS-support).}
\[
\ket{0_L}=\sum_{s\in S_0} c_s \ket{s},\qquad
\ket{1_L}=\sum_{s\in S_0} d_s \ket{\overline{s}},
\qquad
\sum|c_s|^2=\sum|d_s|^2=1.
\]
This strictly enlarges the search space; failure here is stronger evidence than failure under Case A.

\subsubsection{Full KL constraints and the optimized KL loss}
For distance \(3\), it suffices to enforce the Knill--Laflamme equalities for all Pauli strings
\(P\in\{I,X,Y,Z\}^{\otimes 7}\) of weight \(\mathrm{wt}(P)\le 2\):
\[
\bra{0_L}P\ket{1_L}=0,
\qquad
\bra{0_L}P\ket{0_L}=\bra{1_L}P\ket{1_L}.
\]
We minimize the nonnegative sum-of-squares loss
\[
\mathcal{L}(\ket{0_L},\ket{1_L})
:=\sum_{\substack{P:\ \mathrm{wt}(P)\le 2\\P\neq I^{\otimes 7}}}
\Big(
|\bra{0_L}P\ket{1_L}|^2
+
|\bra{0_L}P\ket{0_L}-\bra{1_L}P\ket{1_L}|^2
\Big).
\]
By construction,
\[
\mathcal{L}\ge 0
\quad\text{and}\quad
\mathcal{L}=0\ \Longleftrightarrow\ \text{all full KL equalities hold for every }\mathrm{wt}(P)\le 2.
\]
Thus, finding (robustly) that the global minimum \(\min\mathcal{L}\) is bounded away from \(0\) is strong numerical evidence
that no full-KL \(((7,2,3))\) code exists within the searched SS-support model.

\subsubsection{Observed numerical minima}
We optimized \(\mathcal{L}\) using a combination of (i) differential evolution (global search, multiple seeds, with polishing) and
(ii) multi-start local optimization (Powell) followed by BFGS refinement.

\paragraph{Complementary search (Case A).}
Across multiple random seeds and many independent restarts, the optimizer consistently converged to the same floor:
\[
\boxed{\mathcal{L}_{\min}\approx 0.273738237810.}
\]
Empirically, differential evolution with seeds \(0,1,2,3\) returned \(\mathcal{L}\approx 0.273738238\) agreeing to \(\sim 10^{-9}\),
and independent Powell restarts converged to the same value; BFGS polishing moved only at the \(10^{-10}\) level.

\paragraph{Non-complementary search (Case B).}
Repeating the same optimization with independent \(\ket{1_L}\) coefficients on \(S_7\) again converged to the same floor:
\[
\boxed{\mathcal{L}_{\min}\approx 0.273738237810.}
\]
Numerically, the obstruction therefore does not appear to be an artifact of forcing \(\ket{1_L}=X^{\otimes 7}\ket{0_L}\);
rather, it appears already at the SS-support level.

\subsubsection{Diagnostics at the best point}
At the best complementary minimizer, the largest remaining violations among all \(\mathrm{wt}(P)\le 2\) Paulis were
\[
\max_{P}|\bra{0_L}P\ket{1_L}|\approx 0.16737,
\qquad
\max_{P}\bigl|\bra{0_L}P\ket{0_L}-\bra{1_L}P\ket{1_L}\bigr|\approx 0.28027.
\]
The dominant offenders were weight-\(1\) Paulis (notably off-diagonal failures for \(X_2,X_3\) and similarly \(Y_2,Y_3\), and
diagonal-equality mismatches including \(Y_1\) and some \(Z_j\)).

A striking numerical signature of incompatibility is that there exist points which make the weight-\(1\) constraints essentially exact
(loss \(\sim 10^{-9}\)) while causing large violations at weight \(2\), and conversely points which make the weight-\(2\) constraints
essentially exact while leaving weight-\(1\) constraints nonzero. Minimizing the full system yields a compromise floor
\(\approx 0.273738\), consistent with two constraint families that are individually satisfiable but not jointly satisfiable under the SS supports.

\subsubsection{Reproducible SciPy script}
The following compact script constructs \(S_0(\vec a)\), parameterizes complementary SS-support codes, defines \(\mathcal{L}\),
runs differential evolution, and reports the best loss along with the largest constraint violations:
\begin{lstlisting}[language=Python]
import numpy as np
from itertools import product
from scipy.optimize import differential_evolution

# ----------------------------
# Problem data: BD16, m=8
# ----------------------------
n = 7
m = 8
a = [0,1,1,2,3,3,5]

def residue(bitstr):
    return sum(ai*int(b) for ai,b in zip(a, bitstr)) % m

def complement(s):
    return "".join("1" if b=="0" else "0" for b in s)

all_bits = ["".join(b) for b in product("01", repeat=n)]
S0 = [s for s in all_bits if residue(s) == 0]
idx0 = np.array([int(s,2) for s in S0], dtype=np.int16)
idx1 = np.array([int(complement(s),2) for s in S0], dtype=np.int16)

# ----------------------------
# Pauli action precompute
# ----------------------------
paulis = ['I','X','Y','Z']
def wt(P): return sum(1 for ch in P if ch!='I')

P_list = [''.join(P) for P in product(paulis, repeat=n) if wt(''.join(P))<=2]
idx_noI = [i for i,P in enumerate(P_list) if P!='IIIIIII']

def apply_single(p, bit):
    if p=='I': return 1+0j, bit
    if p=='X': return 1+0j, 1-bit
    if p=='Z': return (1+0j if bit==0 else -1+0j), bit
    if p=='Y': return (1j if bit==0 else -1j), 1-bit
    raise ValueError(p)

def apply_pauli_to_index(P, idx):
    bits=[(idx>>(n-1-i))&1 for i in range(n)]
    phase=1+0j
    out_bits=[]
    for i,p in enumerate(P):
        ph, nb = apply_single(p, bits[i])
        phase *= ph
        out_bits.append(nb)
    out=0
    for b in out_bits:
        out = (out<<1)|b
    return phase, out

targets = np.empty((len(P_list), 2**n), dtype=np.int16)
phases  = np.empty((len(P_list), 2**n), dtype=np.complex128)
for i,P in enumerate(P_list):
    for idx in range(2**n):
        ph,t = apply_pauli_to_index(P, idx)
        targets[i,idx] = t
        phases[i,idx]  = ph

targets_noI = targets[idx_noI]
phases_noI  = phases[idx_noI]

def build_vectors_from_c(c):
    psi0 = np.zeros(128, dtype=np.complex128)
    psi1 = np.zeros(128, dtype=np.complex128)
    psi0[idx0] = c
    psi1[idx1] = c   # complementary: same coeff on complements
    return psi0, psi1

def all_mat_elems(psiA, psiB):
    A_targets = psiA[targets_noI]  # (nP,128)
    return np.sum(np.conj(A_targets) * (phases_noI * psiB[None,:]), axis=1)

def normalize_fix_phase(c):
    c = np.asarray(c, dtype=np.complex128)
    norm = np.linalg.norm(c)
    if norm == 0: return None
    c = c / norm
    k = np.argmax(np.abs(c))
    if abs(c[k]) > 1e-12:
        c = c * np.exp(-1j*np.angle(c[k]))
    return c

# x in R^{28} encodes 14 complex coeffs
def KL_loss(x):
    x = np.asarray(x, dtype=float)
    c = x[:14] + 1j*x[14:]
    c = normalize_fix_phase(c)
    if c is None: return 1e6
    psi0, psi1 = build_vectors_from_c(c)
    m01 = all_mat_elems(psi0, psi1)
    m00 = all_mat_elems(psi0, psi0)
    m11 = all_mat_elems(psi1, psi1)
    diff = m00 - m11
    return float(np.vdot(m01,m01).real + np.vdot(diff,diff).real)

bounds = [(-1,1)]*28
de = differential_evolution(KL_loss, bounds, maxiter=30, popsize=8, seed=0, polish=True, disp=False)
print("S0 size =", len(S0))
print("best KL_loss =", de.fun)
\end{lstlisting}

\subsubsection{Bottom line}
In this SS-support model for \(\vec a=(0,1,1,2,3,3,5)\), the KL loss \(\mathcal{L}\) repeatedly converges to a stable nonzero floor
\(\approx 0.273738\) under both complementary and non-complementary parameterizations, and does not approach numerical zero
(\(\sim 10^{-10}\)) as would be expected if a true full-KL \(((7,2,3))\) solution existed in the searched space.
This provides strong numerical evidence that \((0,1,1,2,3,3,5)\) is an SS/LP false positive: it passes SS and \(Z\)-LP filters but
does not survive the full KL stage under the SS support constraint.

\subsection{A Lean-friendly analytic no--go (in $28$ real variables) for full KL under the complementary ansatz:
$\vec a=(0,1,1,2,3,3,5)$ (mod $8$)}
\label{subsec:lean-nogo-19eq}

\paragraph{Support sets and complementary code states.}
Fix $m=8$ and
\[
\vec a=(0,1,1,2,3,3,5),\qquad b_0=0,\qquad b_1=\sum_{i=1}^7 a_i\equiv 7\pmod 8.
\]
For $s=s_1\cdots s_7\in\{0,1\}^7$ define $\vec a\cdot s:=\sum_i a_i s_i$ (over $\mathbb Z$, reduced mod $8$),
and let $\bar s:=s\oplus 1^7$ be the bitwise complement.
Set
\[
S_0:=\{s:\ \vec a\cdot s\equiv 0\pmod 8\},\qquad
S_1:=\{s:\ \vec a\cdot s\equiv 7\pmod 8\}.
\]
A direct enumeration gives
\[
\begin{aligned}
S_0=\{&
0000000,\ 0000011,\ 0000101,\ 0001110,\ 0011001,\ 0101001,\ 0110110,\\
&1000000,\ 1000011,\ 1000101,\ 1001110,\ 1011001,\ 1101001,\ 1110110\},
\end{aligned}
\]
and $S_1=\overline{S_0}$.
We work with the \emph{complementary ansatz}
\[
\ket{0_L}=\sum_{s\in S_0} c_s\ket{s},\qquad
\ket{1_L}=X^{\otimes 7}\ket{0_L}=\sum_{s\in S_0} c_s\ket{\bar s}.
\]

\paragraph{Paulis, weight, and the KL conditions.}
For $x,z\in\{0,1\}^7$, write
\[
P(x,z):=i^{x\cdot z}X^x Z^z,\qquad
\wt(x,z):=\bigl|\supp(x)\cup\supp(z)\bigr|.
\]
Distance $3$ requires the Knill--Laflamme (KL) conditions for all Paulis with $\wt(x,z)\le 2$:
\[
\bra{i_L}P(x,z)\ket{j_L}=\lambda_{x,z}\,\delta_{ij}\qquad(i,j\in\{0,1\}).
\]
In particular:
\begin{itemize}
\item for each weight-$\le2$ Pauli, the \emph{off-diagonal} condition $\bra{0_L}P(x,z)\ket{1_L}=0$ must hold;
\item for each \emph{diagonal} Pauli, $\bra{0_L}P(x,z)\ket{0_L}=\bra{1_L}P(x,z)\ket{1_L}$ must hold.
\end{itemize}
Under the complementary ansatz, diagonal equality forces $\bra{0_L}Z_i\ket{0_L}=0$ for each $Z_i$
(since $Z_i$ has odd $Z$-weight), and similarly forces $\bra{0_L}Y_1\ket{0_L}=0$.

\paragraph{A small (but sufficient) KL subsystem.}
A full-KL solution must satisfy \emph{all} weight-$\le2$ KL equations, hence in particular it must satisfy
the KL equations for the following explicit weight-$\le2$ Paulis:
\[
Z_i\ (i=1,\dots,7),\quad Y_1,\quad
X_2,\ X_2Z_1,\ X_2Z_3,\quad
X_3,\ X_3Z_1,\ X_3Z_2,\quad
X_1X_2,\ X_1X_2Z_1,\quad
X_1X_3,\ X_1X_3Z_1.
\]
We now rewrite \emph{exactly} these KL constraints as $19$ explicit quadratic equations
in $28$ real variables $v_1,\dots,v_{28}$, and prove that system has no real solution.
This is the formulation used directly in Lean.

%============================================================
\subsubsection*{The $28$ real variables and the $19$ explicit quadratic equations}

\paragraph{Real variables and amplitudes.}
Introduce $28$ real variables $v_1,\dots,v_{28}\in\mathbb R$, and define the $14$ amplitudes
(corresponding to the $14$ strings in $S_0$ in the order listed above) by
\begin{align*}
c_{0000000}&:=v_1+i v_2,&
c_{0000011}&:=v_3+i v_4,&
c_{0000101}&:=v_5+i v_6,&
c_{0001110}&:=v_7+i v_8,\\
c_{0011001}&:=v_9+i v_{10},&
c_{0101001}&:=v_{11}+i v_{12},&
c_{0110110}&:=v_{13}+i v_{14},&
c_{1000000}&:=v_{15}+i v_{16},\\
c_{1000011}&:=v_{17}+i v_{18},&
c_{1000101}&:=v_{19}+i v_{20},&
c_{1001110}&:=v_{21}+i v_{22},&
c_{1011001}&:=v_{23}+i v_{24},\\
c_{1101001}&:=v_{25}+i v_{26},&
c_{1110110}&:=v_{27}+i v_{28}.&&
\end{align*}
Define the nonnegative quadratic forms (``probabilities'')
\[
p_s:=|c_s|^2=(\Re c_s)^2+(\Im c_s)^2\in\mathbb R_{\ge 0}.
\]
Equivalently:
\[
\begin{aligned}
&p_{0000000}=v_1^2+v_2^2,\quad p_{0000011}=v_3^2+v_4^2,\quad p_{0000101}=v_5^2+v_6^2,\quad p_{0001110}=v_7^2+v_8^2,\\
&p_{0011001}=v_9^2+v_{10}^2,\quad p_{0101001}=v_{11}^2+v_{12}^2,\quad p_{0110110}=v_{13}^2+v_{14}^2,\quad p_{1000000}=v_{15}^2+v_{16}^2,\\
&p_{1000011}=v_{17}^2+v_{18}^2,\quad p_{1000101}=v_{19}^2+v_{20}^2,\quad p_{1001110}=v_{21}^2+v_{22}^2,\quad p_{1011001}=v_{23}^2+v_{24}^2,\\
&p_{1101001}=v_{25}^2+v_{26}^2,\quad p_{1110110}=v_{27}^2+v_{28}^2.
\end{aligned}
\]

\paragraph{(E1) Normalization.}
KL requires $\langle 0_L\mid 0_L\rangle=1$, i.e.
\begin{equation}
\label{eq:poly-norm}
v_1^2+v_2^2+\cdots+v_{27}^2+v_{28}^2 \;=\; 1.
\tag{E1}
\end{equation}

\paragraph{(E2)--(E8) The seven $Z_i$ constraints.}
For each $i$, diagonal KL forces $\bra{0_L}Z_i\ket{0_L}=0$.
Since $Z_i\ket{s}=(-1)^{s_i}\ket{s}$, this yields a linear equation in the $p_s$.
Substituting $p_s=v_{\Re}^2+v_{\Im}^2$ and multiplying by $2$ (to clear $\tfrac12$) gives:
\begin{align}
\label{eq:poly-Z1}
&2\Big((v_{15}^2+v_{16}^2)+(v_{17}^2+v_{18}^2)+(v_{19}^2+v_{20}^2)+(v_{21}^2+v_{22}^2)+(v_{23}^2+v_{24}^2)+(v_{25}^2+v_{26}^2)+(v_{27}^2+v_{28}^2)\Big)=1,\tag{E2}\\
\label{eq:poly-Z2}
&2\Big((v_{11}^2+v_{12}^2)+(v_{13}^2+v_{14}^2)+(v_{25}^2+v_{26}^2)+(v_{27}^2+v_{28}^2)\Big)=1,\tag{E3}\\
\label{eq:poly-Z3}
&2\Big((v_{9}^2+v_{10}^2)+(v_{13}^2+v_{14}^2)+(v_{23}^2+v_{24}^2)+(v_{27}^2+v_{28}^2)\Big)=1,\tag{E4}\\
\label{eq:poly-Z4}
&2\Big((v_{7}^2+v_{8}^2)+(v_{9}^2+v_{10}^2)+(v_{11}^2+v_{12}^2)+(v_{21}^2+v_{22}^2)+(v_{23}^2+v_{24}^2)+(v_{25}^2+v_{26}^2)\Big)=1,\tag{E5}\\
\label{eq:poly-Z5}
&2\Big((v_{5}^2+v_{6}^2)+(v_{7}^2+v_{8}^2)+(v_{13}^2+v_{14}^2)+(v_{19}^2+v_{20}^2)+(v_{21}^2+v_{22}^2)+(v_{27}^2+v_{28}^2)\Big)=1,\tag{E6}\\
\label{eq:poly-Z6}
&2\Big((v_{3}^2+v_{4}^2)+(v_{7}^2+v_{8}^2)+(v_{13}^2+v_{14}^2)+(v_{17}^2+v_{18}^2)+(v_{21}^2+v_{22}^2)+(v_{27}^2+v_{28}^2)\Big)=1,\tag{E7}\\
\label{eq:poly-Z7}
&2\Big((v_{3}^2+v_{4}^2)+(v_{5}^2+v_{6}^2)+(v_{9}^2+v_{10}^2)+(v_{11}^2+v_{12}^2)+(v_{17}^2+v_{18}^2)+(v_{19}^2+v_{20}^2)+(v_{23}^2+v_{24}^2)+(v_{25}^2+v_{26}^2)\Big)=1.\tag{E8}
\end{align}

\paragraph{(E9) The $Y_1$ diagonal constraint.}
Diagonal KL forces $\bra{0_L}Y_1\ket{0_L}=0$.
Pairing $(0u,1u)$ strings in $S_0$ yields
\[
0=\bra{0_L}Y_1\ket{0_L}=2\sum_u \Im\!\bigl(\overline{c_{0u}}c_{1u}\bigr),
\]
and rewriting each $\Im(\overline{(a+ib)}(c+id))$ as $ad-bc$ gives:
\begin{equation}
\label{eq:poly-Y1}
(v_1v_{16}-v_2v_{15})+(v_3v_{18}-v_4v_{17})+(v_5v_{20}-v_6v_{19})+(v_7v_{22}-v_8v_{21})
+(v_9v_{24}-v_{10}v_{23})+(v_{11}v_{26}-v_{12}v_{25})+(v_{13}v_{28}-v_{14}v_{27})=0.
\tag{E9}
\end{equation}

\paragraph{(E10)--(E12) Off-diagonal KL from $X_2$, $X_2Z_1$, $X_2Z_3$.}
For each listed weight-$\le2$ Pauli, KL requires $\bra{0_L}P\ket{1_L}=0$.
For $P\in\{X_2,\ X_2Z_1,\ X_2Z_3\}$ (i.e.\ $x=e_2$ and $z\in\{0,e_1,e_3\}$),
the four-term intersection sum collapses to bilinear constraints on
$\overline{c_{0101001}}c_{1110110}$ and $\overline{c_{0110110}}c_{1101001}$.
Splitting into real/imaginary parts and writing in the $v_i$ gives:
\begin{align}
\label{eq:poly-e2-z0}
&(v_{11}v_{27}+v_{12}v_{28})+(v_{13}v_{25}+v_{14}v_{26})=0,\tag{E10}\\
\label{eq:poly-e2-ze1}
&(v_{11}v_{28}-v_{12}v_{27})+(v_{13}v_{26}-v_{14}v_{25})=0,\tag{E11}\\
\label{eq:poly-e2-ze3}
&(v_{11}v_{28}-v_{12}v_{27})-(v_{13}v_{26}-v_{14}v_{25})=0.\tag{E12}
\end{align}

\paragraph{(E13)--(E15) Off-diagonal KL from $X_3$, $X_3Z_1$, $X_3Z_2$.}
Likewise, for $P\in\{X_3,\ X_3Z_1,\ X_3Z_2\}$ (i.e.\ $x=e_3$ and $z\in\{0,e_1,e_2\}$),
the off-diagonal sums collapse to bilinear constraints on
$\overline{c_{0011001}}c_{1110110}$ and $\overline{c_{0110110}}c_{1011001}$, giving:
\begin{align}
\label{eq:poly-e3-z0}
&(v_{9}v_{27}+v_{10}v_{28})+(v_{13}v_{23}+v_{14}v_{24})=0,\tag{E13}\\
\label{eq:poly-e3-ze1}
&(v_{9}v_{28}-v_{10}v_{27})+(v_{13}v_{24}-v_{14}v_{23})=0,\tag{E14}\\
\label{eq:poly-e3-ze2}
&(v_{9}v_{28}-v_{10}v_{27})-(v_{13}v_{24}-v_{14}v_{23})=0.\tag{E15}
\end{align}

\paragraph{(E16)--(E17) Off-diagonal KL from $X_1X_2$, $X_1X_2Z_1$.}
For $P\in\{X_1X_2,\ X_1X_2Z_1\}$ (i.e.\ $x=e_1+e_2$ and $z\in\{0,e_1\}$),
the off-diagonal sums reduce to real-part constraints on
$\overline{c_{1101001}}c_{1110110}$ and $\overline{c_{0101001}}c_{0110110}$, giving:
\begin{align}
\label{eq:poly-e1e2-z0}
&(v_{25}v_{27}+v_{26}v_{28})+(v_{11}v_{13}+v_{12}v_{14})=0,\tag{E16}\\
\label{eq:poly-e1e2-ze1}
&(v_{25}v_{27}+v_{26}v_{28})-(v_{11}v_{13}+v_{12}v_{14})=0.\tag{E17}
\end{align}

\paragraph{(E18)--(E19) Off-diagonal KL from $X_1X_3$, $X_1X_3Z_1$.}
Similarly, for $P\in\{X_1X_3,\ X_1X_3Z_1\}$ (i.e.\ $x=e_1+e_3$ and $z\in\{0,e_1\}$),
we obtain:
\begin{align}
\label{eq:poly-e1e3-z0}
&(v_{23}v_{27}+v_{24}v_{28})+(v_{9}v_{13}+v_{10}v_{14})=0,\tag{E18}\\
\label{eq:poly-e1e3-ze1}
&(v_{23}v_{27}+v_{24}v_{28})-(v_{9}v_{13}+v_{10}v_{14})=0.\tag{E19}
\end{align}

\paragraph{Summary.}
Equations \eqref{eq:poly-norm}--\eqref{eq:poly-e1e3-ze1} are $19$ explicit real polynomial equations,
each of total degree $2$, in the $28$ real unknowns $(v_1,\dots,v_{28})$.
Every full-KL complementary solution must satisfy them.

%============================================================
\subsubsection*{No-solution proof (purely in the $v_i$ equations)}

\begin{theorem}[The $19$ quadratic equations have no real solution]
\label{thm:no-solution-19eq}
There do not exist real numbers $v_1,\dots,v_{28}$ satisfying
\eqref{eq:poly-norm}--\eqref{eq:poly-e1e3-ze1}.
\end{theorem}

\begin{proof}
Assume for contradiction that $(v_1,\dots,v_{28})$ satisfies all $19$ equations.
Let the quadratic forms $p_s$ be as defined above, so $p_s\ge 0$.

\paragraph{Step 1: Linear elimination of the $Z$-system (pair sums).}
Equations \eqref{eq:poly-norm} and \eqref{eq:poly-Z1}--\eqref{eq:poly-Z7} are linear in the $p_s$.
A routine linear elimination (Lean: \texttt{linarith}) yields the \emph{pair-sum} identities
\begin{equation}
\label{eq:pairsums-v}
\begin{aligned}
&p_{0000000}+p_{1000000}=\frac{1}{16},\quad
p_{0000011}+p_{1000011}=\frac{1}{16},\quad
p_{0000101}+p_{1000101}=\frac{1}{16},\quad
p_{0001110}+p_{1001110}=\frac{1}{8},\\
&p_{0011001}+p_{1011001}=\frac{3}{16},\quad
p_{0101001}+p_{1101001}=\frac{3}{16},\quad
p_{0110110}+p_{1110110}=\frac{5}{16}.
\end{aligned}
\end{equation}
Equivalently, with the shorthand
\[
p_0:=p_{0000000},\ p_1:=p_{0000011},\dots,\ p_{13}:=p_{1110110},
\]
we have
\begin{equation}
\label{eq:pairsums-short}
p_0+p_7=\frac{1}{16},\quad p_1+p_8=\frac{1}{16},\quad p_2+p_9=\frac{1}{16},\quad p_3+p_{10}=\frac{1}{8},\quad
p_4+p_{11}=\frac{3}{16},\quad p_5+p_{12}=\frac{3}{16},\quad p_6+p_{13}=\frac{5}{16}.
\end{equation}

\paragraph{Step 2: The $x=e_2$ block forces a multiplicative identity.}
From \eqref{eq:poly-e2-ze1} and \eqref{eq:poly-e2-ze3}, adding and subtracting gives
\[
(v_{11}v_{28}-v_{12}v_{27})=0,\qquad (v_{13}v_{26}-v_{14}v_{25})=0.
\]
From \eqref{eq:poly-e2-z0} we also have
\[
(v_{11}v_{27}+v_{12}v_{28})=-(v_{13}v_{25}+v_{14}v_{26}).
\]
Now use the real identity
\[
(ac+bd)^2+(ad-bc)^2=(a^2+b^2)(c^2+d^2)
\]
(with $a=v_{11},b=v_{12},c=v_{27},d=v_{28}$ and similarly for $(v_{13},v_{14},v_{25},v_{26}$)).
Squaring and adding the ``real-part'' and ``imag-part'' equations yields
\[
(v_{11}^2+v_{12}^2)(v_{27}^2+v_{28}^2)=(v_{13}^2+v_{14}^2)(v_{25}^2+v_{26}^2),
\]
i.e.
\begin{equation}
\label{eq:mult-e2-v}
p_5\,p_{13}=p_6\,p_{12}.
\end{equation}

\paragraph{Step 3: The $x=e_3$ block forces a second multiplicative identity.}
Exactly the same argument applied to \eqref{eq:poly-e3-z0}--\eqref{eq:poly-e3-ze2} gives
\begin{equation}
\label{eq:mult-e3-v}
p_4\,p_{13}=p_6\,p_{11}.
\end{equation}

\paragraph{Step 4: Linearize \eqref{eq:mult-e2-v}--\eqref{eq:mult-e3-v} using the pair sums.}
From \eqref{eq:pairsums-short}, $p_5=\frac{3}{16}-p_{12}$ and $p_6=\frac{5}{16}-p_{13}$.
Substitute into \eqref{eq:mult-e2-v}:
\[
\Bigl(\frac{3}{16}-p_{12}\Bigr)p_{13}=\Bigl(\frac{5}{16}-p_{13}\Bigr)p_{12}
\ \Longrightarrow\
\frac{3}{16}p_{13}=\frac{5}{16}p_{12}
\ \Longrightarrow\
p_{12}=\frac35\,p_{13}.
\]
Similarly, using $p_4=\frac{3}{16}-p_{11}$ in \eqref{eq:mult-e3-v} gives $p_{11}=\frac35\,p_{13}$.
Hence
\begin{equation}
\label{eq:linearized-v}
p_{11}=p_{12}=\frac35\,p_{13},
\qquad
p_4=p_5=\frac35\,p_6.
\end{equation}

\paragraph{Step 5: Uniform lower bounds on $p_{13}$ and $p_6$.}
Let $t:=p_{13}\ge 0$.
Then $p_{11}+p_{12}+p_{13}=\frac{11}{5}t$ by \eqref{eq:linearized-v}.
Using $p_7\ge 0$ and $p_7=\frac12-(p_8+p_9+p_{10}+p_{11}+p_{12}+p_{13})$ (from the linear system),
we get
\[
\frac{11}{5}t\le \frac12\quad\Longrightarrow\quad t\le \frac{5}{22},
\]
and hence from $p_6+p_{13}=\frac{5}{16}$ in \eqref{eq:pairsums-short},
\[
p_6=\frac{5}{16}-t\ge \frac{5}{16}-\frac{5}{22}=\frac{15}{176}.
\]
For the lower bound on $t$, use $p_0\ge 0$ and
$p_0=(p_8+p_9+p_{10}+p_{11}+p_{12}+p_{13})-\frac{7}{16}$:
since $p_8\le\frac{1}{16}$, $p_9\le\frac{1}{16}$, $p_{10}\le\frac18$ by \eqref{eq:pairsums-short} and nonnegativity,
we have $p_8+p_9+p_{10}\le\frac14$, so
\[
0\le \frac14+\frac{11}{5}t-\frac{7}{16}
\quad\Longrightarrow\quad
t\ge \frac{15}{176}.
\]
Therefore
\begin{equation}
\label{eq:bounds-v}
p_{13}\ge \frac{15}{176}>0,
\qquad
p_6\ge \frac{15}{176}>0.
\end{equation}

\paragraph{Step 6: Real-part vanishings from \eqref{eq:poly-e1e2-z0}--\eqref{eq:poly-e1e3-ze1}.}
Adding and subtracting \eqref{eq:poly-e1e2-z0} and \eqref{eq:poly-e1e2-ze1} gives
\begin{equation}
\label{eq:re-vanish-1}
(v_{25}v_{27}+v_{26}v_{28})=0,\qquad (v_{11}v_{13}+v_{12}v_{14})=0.
\end{equation}
Similarly, adding/subtracting \eqref{eq:poly-e1e3-z0} and \eqref{eq:poly-e1e3-ze1} gives
\begin{equation}
\label{eq:re-vanish-2}
(v_{23}v_{27}+v_{24}v_{28})=0,\qquad (v_{9}v_{13}+v_{10}v_{14})=0.
\end{equation}

\paragraph{Step 7: A global phase rotation (a simultaneous planar rotation) fixes $v_{28}=0$.}
Since $p_{13}=v_{27}^2+v_{28}^2>0$ by \eqref{eq:bounds-v}, we have $(v_{27},v_{28})\neq (0,0)$.
All $19$ equations are invariant under multiplying every amplitude $c_s$ by a common phase $e^{i\theta}$,
which in real coordinates is the simultaneous rotation
\[
(v_{2k+1},v_{2k+2})\mapsto (v_{2k+1}\cos\theta-v_{2k+2}\sin\theta,\ v_{2k+1}\sin\theta+v_{2k+2}\cos\theta)
\quad(k=0,\dots,13).
\]
Choosing $\theta$ so that $(v_{27},v_{28})$ rotates to $(\sqrt{p_{13}},0)$, we may assume \emph{without loss of generality} that
\begin{equation}
\label{eq:gauge-v}
v_{28}=0,\qquad v_{27}>0.
\end{equation}

Under the gauge \eqref{eq:gauge-v}, equations \eqref{eq:poly-e2-ze1}, \eqref{eq:poly-e3-ze1},
\eqref{eq:re-vanish-1}, and \eqref{eq:re-vanish-2} immediately force
\[
v_{12}=0,\qquad v_{10}=0,\qquad v_{25}=0,\qquad v_{23}=0,\qquad v_{13}=0.
\]
Also $p_6=v_{13}^2+v_{14}^2=v_{14}^2>0$ by \eqref{eq:bounds-v}, hence
\begin{equation}
\label{eq:v14-nonzero}
v_{14}\neq 0.
\end{equation}
Finally, since $p_4=\frac35 p_6>0$ and $p_5=\frac35 p_6>0$ by \eqref{eq:linearized-v},
we have $v_9\neq 0$ and $v_{11}\neq 0$ (because $p_4=v_9^2+v_{10}^2=v_9^2$ and $p_5=v_{11}^2+v_{12}^2=v_{11}^2$ under the gauge).

\paragraph{Step 8: Split the $Y_1$ equation into a small part $U$ and a large part $T$.}
Write the $Y_1$ equation \eqref{eq:poly-Y1} as
\[
U+T=0
\]
where
\[
\begin{aligned}
U&:=(v_1v_{16}-v_2v_{15})+(v_3v_{18}-v_4v_{17})+(v_5v_{20}-v_6v_{19})+(v_7v_{22}-v_8v_{21}),\\
T&:=(v_9v_{24}-v_{10}v_{23})+(v_{11}v_{26}-v_{12}v_{25})+(v_{13}v_{28}-v_{14}v_{27}).
\end{aligned}
\]
Under the gauge consequences above ($v_{10}=v_{12}=v_{13}=v_{23}=v_{25}=v_{28}=0$), this simplifies to
\[
T=v_9v_{24}+v_{11}v_{26}-v_{14}v_{27}.
\]
Now use \eqref{eq:poly-e2-z0} and \eqref{eq:poly-e3-z0} in the gauge:
\[
\eqref{eq:poly-e2-z0}:\quad v_{11}v_{27}+v_{14}v_{26}=0
\ \Longrightarrow\ 
v_{26}=-\frac{v_{11}v_{27}}{v_{14}},
\qquad
\eqref{eq:poly-e3-z0}:\quad v_{9}v_{27}+v_{14}v_{24}=0
\ \Longrightarrow\ 
v_{24}=-\frac{v_{9}v_{27}}{v_{14}},
\]
where division by $v_{14}$ is valid by \eqref{eq:v14-nonzero}.
Substituting into $T$ gives the \emph{exact} identity
\[
T=-\frac{v_{27}}{v_{14}}\,(v_9^2+v_{11}^2+v_{14}^2),
\]
hence
\begin{equation}
\label{eq:T-exact-v}
|T|=\frac{|v_{27}|}{|v_{14}|}\,(v_9^2+v_{11}^2+v_{14}^2).
\end{equation}
Rewrite in terms of $p$'s under the gauge:
\[
|v_{27}|=\sqrt{p_{13}},\quad |v_{14}|=\sqrt{p_6},\quad v_9^2=p_4,\quad v_{11}^2=p_5,\quad v_{14}^2=p_6,
\]
so \eqref{eq:T-exact-v} becomes
\[
|T|=\frac{\sqrt{p_{13}}}{\sqrt{p_6}}\,(p_4+p_5+p_6).
\]
Using \eqref{eq:linearized-v}, $p_4=p_5=\frac35 p_6$, so $p_4+p_5+p_6=\frac{11}{5}p_6$, and therefore
\begin{equation}
\label{eq:T-sqrt-v}
|T|=\frac{11}{5}\sqrt{p_6p_{13}}.
\end{equation}
With \eqref{eq:bounds-v}, $\sqrt{p_6p_{13}}\ge \frac{15}{176}$, hence
\begin{equation}
\label{eq:T-lower-v}
|T|\ge \frac{11}{5}\cdot \frac{15}{176}=\frac{3}{16}.
\end{equation}

\paragraph{Step 9: A uniform upper bound for $|U|$.}
Each summand of $U$ has the form $ad-bc$ and satisfies
\[
|ad-bc|\le \sqrt{(a^2+b^2)(c^2+d^2)}\le \frac{(a^2+b^2)+(c^2+d^2)}{2}.
\]
Apply this to the four pairs $(v_1,v_2)$ with $(v_{15},v_{16})$, $(v_3,v_4)$ with $(v_{17},v_{18})$, etc.
Using the pair sums \eqref{eq:pairsums-short}:
\[
p_0+p_7=\frac{1}{16},\quad p_1+p_8=\frac{1}{16},\quad p_2+p_9=\frac{1}{16},\quad p_3+p_{10}=\frac18,
\]
we obtain
\[
|v_1v_{16}-v_2v_{15}|\le \frac{p_0+p_7}{2}=\frac{1}{32},\quad
|v_3v_{18}-v_4v_{17}|\le \frac{1}{32},\quad
|v_5v_{20}-v_6v_{19}|\le \frac{1}{32},\quad
|v_7v_{22}-v_8v_{21}|\le \frac{1}{16}.
\]
By the triangle inequality,
\begin{equation}
\label{eq:U-upper-v}
|U|\le \frac{1}{32}+\frac{1}{32}+\frac{1}{32}+\frac{1}{16}=\frac{5}{32}.
\end{equation}

\paragraph{Step 10: Contradiction.}
From $U+T=0$ we have $|T|=|U|$.
But \eqref{eq:T-lower-v} gives $|T|\ge \frac{3}{16}=\frac{6}{32}$ while \eqref{eq:U-upper-v} gives $|U|\le \frac{5}{32}$.
This is impossible. Hence no real solution to \eqref{eq:poly-norm}--\eqref{eq:poly-e1e3-ze1} exists.
\end{proof}

%============================================================
\subsubsection*{Conclusion: no complementary full-KL solution}

\begin{corollary}[No complementary full--KL solution for $\vec a=(0,1,1,2,3,3,5)$]
\label{cor:nogo-fullKL-a0112335}
There is no choice of amplitudes $(c_s)_{s\in S_0}$ such that
\[
\ket{0_L}=\sum_{s\in S_0} c_s\ket{s},\qquad
\ket{1_L}=X^{\otimes 7}\ket{0_L},
\]
satisfies the full Knill--Laflamme conditions for all Pauli operators $P(x,z)$ of weight $\le 2$.
\end{corollary}

\begin{proof}
Any full-KL complementary solution must satisfy the KL constraints for each of the specific weight-$\le2$ Paulis
used to derive \eqref{eq:poly-norm}--\eqref{eq:poly-e1e3-ze1}; hence it would give a real solution
to the $19$ quadratic equations. This contradicts Theorem~\ref{thm:no-solution-19eq}.
\end{proof}


\end{document}
